\documentclass[11pt,a4paper]{article}
\usepackage[T1]{fontenc}
\usepackage[utf8]{inputenc}
\usepackage{lmodern,amsmath,amssymb,amsthm,mathtools,microtype,mathrsfs}
\usepackage[margin=25mm,top=23mm,bottom=25mm]{geometry}
\usepackage{booktabs,array,longtable,enumitem,xcolor,fancyhdr}
\usepackage[colorlinks=true,linkcolor=blue!45!black,urlcolor=blue!45!black,citecolor=blue!45!black]{hyperref}
\hypersetup{pdftitle={Reading the Alpoge--Claude RH and Dirichlet-GRH proof: working version 0.57},pdfauthor={ChatGPT},pdfsubject={Partial human-readable reconstruction with explicit arithmetic inputs},pdfkeywords={Riemann hypothesis, Lean, theta lifts, prime sums, Arakelov slopes}}
\setlength{\parindent}{0pt}
\setlength{\parskip}{5pt plus 1pt}
\setlist{topsep=3pt,itemsep=3pt,leftmargin=*}
\setlength{\headheight}{14pt}
\pagestyle{fancy}\fancyhf{}
\fancyhead[L]{\small RH and Dirichlet GRH: a reconstruction}
\fancyhead[R]{\small Working version 0.57}
\fancyfoot[C]{\thepage}
\renewcommand{\headrulewidth}{0.3pt}
\numberwithin{equation}{section}
\newtheorem{theorem}{Theorem}[section]
\newtheorem{proposition}[theorem]{Proposition}
\newtheorem{lemma}[theorem]{Lemma}
\newtheorem{corollary}[theorem]{Corollary}
\theoremstyle{remark}\newtheorem{remark}[theorem]{Remark}
\newcommand{\Rea}{\operatorname{Re}}
\newcommand{\HS}{\mathrm{HS}}
\newcommand{\tr}{\operatorname{tr}}
\newcommand{\codim}{\operatorname{codim}}
\newcommand{\covol}{\operatorname{covol}}
\newcommand{\OK}{\mathcal O_K}
\newcommand{\src}[2]{\href{https://github.com/CocoSgt/riemann-hypothesis-lean/blob/7cbf7a0c32fca1badf05931879a37c3f8a2b03b4/src/#1}{\texttt{#2}}}
\newcommand{\status}[1]{\textbf{Status:} #1}
\begin{document}
\begin{center}
{\LARGE\bfseries RH and Dirichlet GRH: new human proofs}\\[7pt]
{\large Version 0.57 --- 24 November 2026}\\[5pt]
{\large Local functors, arithmetic components, and exact geometric periods}
\end{center}
\textbf{Scope.} This edition contains the current audit and all new chapters.
It is a partial human reconstruction of the accepted formal proof.
Earlier named claims are in the cumulative manuscript and preserved components.
Each result below states its original inputs and exact remaining producers.
\setcounter{tocdepth}{1}
\tableofcontents
\clearpage
\begingroup\emergencystretch=3em
\section{Current proof boundary after the component, period and local-functor joins}
\label{sec:dependency-audit-v57}

\paragraph{The purpose of this audit.}
We continue to accept the pinned formal proof as correct.  The question
here is which portions now have a readable mathematical proof, and which
original constructions still need one.  This edition adds 101 claims
to the 3645 preserved claims, for 3746 in all.  Every new claim has
a separate review bound to its exact proof text and source statements.
The records distinguish a complete original implication from a construction
of all of that implication's inputs.

\paragraph{What has moved.}
The arithmetic chapter now reaches the full original rank-two component
theorem.  Its local chains, dimension inequalities, characteristic-prime
curves, homogeneous incidence and seed localization are proved on the
actual rings.  Its original finite-weight, automorphic propagation and
seed assumptions are retained.  The rank-three parity and exceptional
component arguments remain distinct.

The period chapter proves the original universal Gram expansion at its
precise analytic inputs.  It establishes the scalar absolute majorant
before changing sums and integrals, and retains the original rational
returns, measures, finite phase, column convention and normalizing scalar.
It does not infer a uniform spectral estimate from an operator identity.

The geometric chapters prove the full determinant-jet window and carry it
to actual integral row germs, with an explicit unit determinant factor.
The original Cartier normal form remains the input to that last
application.  They also remove the full original section-growth input:
one very ample power has exactly polynomial section dimension at every
positive multiple.  The classical theorems used for this replacement and
the dictionary identifying their objects are printed in the proof.

The local period-ring chapter constructs a genuinely fixed integral basis,
the nonlinear Frobenius conjugacy, and the original etale functor with
every lattice morphism.  The earlier linear solution-descent theorem
alone did not supply this conjugacy.  The new construction supplies the
actual etale input to the previously proved conditional Wach theorem.
Herr class maps, quantitative comparison and the common global family
remain to be reconstructed at their full scope.

\paragraph{Whole obligations and shared objects.}
The selected sufficient graph retains its 121 nodes and 161 edges.
There remain 24 whole selected obligations: 17 class C and 7
class B.  Separately, 17 are shared by RH and GRH and 7
belong to the extra unrestricted Dirichlet-GRH route.  There is no new
whole selected closure.  The earlier full
\texttt{GLOBAL\_\allowbreak LOCAL\_\allowbreak GAIN} closure is preserved.

The unchanged whole count does not undo the completed subtheorems.
These large obligations combine several constructions and their
compatibility conditions.  A completed local ingredient becomes useful
only through its original maps on the common objects; its existence on
an independently chosen replacement object would not supply that join.
The register below records both the supplied ingredient and the exact
remaining construction, so subsequent work can begin at the first
unproved input rather than repeating earlier results.

\paragraph{Reading and review conventions.}
Read ACC for arithmetic components, PGR for periods, DJW for determinant
jets and integral row tests, SGV for the section-growth replacement,
and FET for faithful invariants and the local functor.  Published
foundations are explicit inputs, not unlabelled proof gaps.  The proof
and source hashes ensure that a review refers to the distributed text;
they are not a substitute for mathematical review or formal verification.
Draft TeX checks syntax and references.  No PDF rendering or new Lean
build is asserted.

\subsection{Current register of the complete remaining contracts}

\paragraph{1. GLOBAL\_\allowbreak DENOMINATOR — C — Actual LinkN local table and binding producer}

Closed(Des $\mathcal C$ 3) LinkN: the actual etaSum is bounded with exact dZ0 and same-host vvTerm, subject only to the literal ShaFinOdd and summability guards.

\emph{Shared objects and choice order.} one FamP; one class $\mathcal A$ and its target-compatible subclass; one coefficient field per family datum; one x and i at each application; original selected rows, words and levels; thresholds chosen before the common family; actual D.nd; actual etaJ and Aj; sha\_\allowbreak eq

\emph{Completed inputs (version 0.57).} FET1–FET26 full faithful invariants, fixed basis, nonlinear conjugacy, EtBasis and WEt; actual input to conditional WachEx, WPR1–WPR30 (full prescribed prism and original BS plus LocTied implies WachEx; no period-ring supplier left), IFD1–IFD26 (full original inhomogeneous invariant matrix descent; exact one-period-supplier Wach return), WIS1–WIS26 (full four-clause actual inclusion and fixed/near-parameter completeness; exact two-supplier Wach return), WLD1-WLD26 (full generic Wach lattice constructors; exact conditional three-supplier all-layer return), SSC1-SSC17 (full original Selmer SignatureChoice and family setup retaining only coherent Wach comparison), DFD1-DFD15, DNE1-DNE14 (full original ConclF-to-LinkN consumer)

\emph{Remaining producer scope.} FET now proves the full original FaithH, actual cyclotomic-kernel fixed basis, nonlinear EtalePhiMatrixDecompletionToWach conjugacy, EtBasis and WEt on the original period rings for every original prime and lattice. Perfectoid cyclotomic fixed-field recovery is proved through sharp coordinates, with the finite-extension theorem explicitly retained as a classical foundation. Finite trace-one averaging and adic lifting construct an integral fixed basis without division by the finite quotient order. The nonlinear conjugacy is distinct from the earlier linear solution descent IFD25: a parameter contraction first gives a basis modulo p, and IFD24 supplies each subsequent p-adic matrix correction. FET descends the actual Galois action and every lattice morphism, preserving all WEt fields and scalar towers. WPR30 therefore gives WachEx for this actual functor on each supplied good layer under its original BS and LocTied/IsStd premises. The actual ExistsLocalWachFunctor assembly still needs the group, series and inclusion records from NormCoeAlgebraMap.lean, OneAddXMulDerivativeBinomial.lean and TeichPow.lean; the downstream WIS/WPR/FET statements retain these records as supplied data. Reuse their completed generic maps while assembling these original records. Then prove the Herr class maps in ExistsUniqueFits.lean and the comparison producer in ChiGamPowNeOne.lean. Preserve torsion, full inverse-image finite local conditions, coefficient-normalized lengths and both orders of the signed clusters. Full HW/Fon/Cmp is not yet claimed. Every original local row/table estimate must then hold in one common ConclF family, retaining ShaFinOdd, exceptional sets, vvTerm summability and uniform thresholds. There is no additional outer-layer transition-system hypothesis. SSC, DFD12 and DNE remain available, with DNE K=2000(|Ca|+1)\textasciicircum{}3 fixed before the family. Whole GLOBAL\_\allowbreak DENOMINATOR remains open.

\emph{Source.} \href{https://github.com/CocoSgt/riemann-hypothesis-lean/blob/7cbf7a0c32fca1badf05931879a37c3f8a2b03b4/src/Formalism/SharpRestriction/GlobalTheoremInputs.lean#L45-L48}{LinkN}.

\emph{Source.} \href{https://github.com/CocoSgt/riemann-hypothesis-lean/blob/7cbf7a0c32fca1badf05931879a37c3f8a2b03b4/src/Formalism/SharpRestriction/GlobalTheoremInputs.lean#L40-L44}{LinkN\_\allowbreak at}.

\emph{Source.} \href{https://github.com/CocoSgt/riemann-hypothesis-lean/blob/7cbf7a0c32fca1badf05931879a37c3f8a2b03b4/src/Development/AutomorphicArithmetic/CharacterBoxes/DivinNOfHypV.lean#L117-L134}{divinN\_\allowbreak of\_\allowbreak hypV}.



\paragraph{2. GLOBAL\_\allowbreak RETAINED\_\allowbreak DEFECT — C — Original retained-row defect bound LinkIRR}

Rows.DelpNonneg and delp$\le$exp(-x\textasciicircum{}(ez/7)) uniformly for every late actual row datum. CURRENT VERSION0.47 QUALIFICATION: Finish the same-D/common-family adelic setting after ASG: normalized nonzero Bessel W and decay, measured-field block, original Siegel-Weil rescaling and measure/character ties, Bessel uniqueness/local ties, then the original AutTie/RowDictTie/LabTie. RowsFinDim is now proved universally by RFD. Supply the actual CountIdempotentsClass theorem or its three sufficient arithmetic cylinder counts, with one simultaneous family threshold; reuse the completed ARO structure. Completed new original producers: ARO1-ARO30, RFD1-RFD9, ASG1-ASG18. Any earlier narrative naming these completed producers as outstanding is historical; this qualification is the current scope authority. CURRENT VERSION0.48 QUALIFICATION: Finish the same-D/common-family adelic setting after ASG and AMD: normalized nonzero Bessel W and decay, original orbital/fiber and Siegel-Weil measure/character ties, Bessel uniqueness/local ties, and the full original transport theorem, then AutTie/RowDictTie/LabTie. The complete MExists measured-data block and unipotent covolume normalization are now supplied, and universal RowsFinDim remains available. Supply the actual CountIdempotentsClass theorem or its three sufficient arithmetic cylinder counts with one simultaneous family threshold. Completed current producers or conditional consumers: ARO1-ARO30, RFD1-RFD9, ASG1-ASG18, AMD1-AMD17. Earlier frontier descriptions are historical; this qualification is the current scope authority. Original whole-contract premises and quantifiers are retained. CURRENT VERSION0.49 QUALIFICATION: Prove the exact original BesselNVGap nonzero integral on its supplied discrete-series, compact-representation and character data. CoefWeightedLoc, WlocExists with its original nonvanishing antecedent, WFromLoc, the normalized global product and all-later-measured-block decay are complete. On the inherited Forms0 scope this leaves BesselNVGap for ConcreteSettingExists. Continue original orbital/fiber and Siegel-Weil measure/character ties, the required Bessel uniqueness/local ties and the full original transport theorem, then AutTie/RowDictTie/LabTie. Supply the actual CountIdempotentsClass theorem or its three sufficient arithmetic cylinder counts with one simultaneous family threshold. Completed current producers or conditional consumers: ARO1-ARO30, RFD1-RFD9, ASG1-ASG18, AMD1-AMD17, BWP1-BWP24 at their exact BesselNVGap/Forms0 scope. Earlier frontier descriptions are historical; this qualification is the current scope authority. Original whole-contract premises and quantifiers are retained.

\emph{Shared objects and choice order.} one FamP; one class $\mathcal A$ and its target-compatible subclass; one coefficient field per family datum; one x and i at each application; original selected rows, words and levels; thresholds chosen before the common family; AllRows A; LabelCountGoodInjective; AllJoinA; CountAllRows

\emph{Completed inputs (version 0.57).} SSC full Selmer-signature comparison, DFD common-family subresults, DNE full separate denominator consumer, WRS1-WRS10 full selected-class local gain, Inherited first DivDef3 and row-pair integrability

\emph{Remaining producer scope.} Complete the original F3 local coefficient bridges and F4 common-overfield readings, with valuation/norm/GoodCoef/iq and the entire coefficient field, and actual complex Hecke standard/no-Jordan plus arithmetic character data on the same form datum. Construct the actual base rows, host/frame lengths, nonzero chains and group ties, MultEq support, eigensystem/carrier/parts, all-embedding logLam, HD/d3b and NEb. Keep the stronger second DivRows3 scope distinct from the completed first DivDef3 and row-pair integrability. The full arithmetic entry and conditional common-frame construction do not by themselves install HOST\_\allowbreak NUMBER\_\allowbreak S, HOST\_\allowbreak FRAME\_\allowbreak LENGTH, TASK\_\allowbreak HOST\_\allowbreak VAN\_\allowbreak E\_\allowbreak INSTALL, TASK\_\allowbreak HOST\_\allowbreak XPP\_\allowbreak INSTALL or TASK\_\allowbreak COUNT\_\allowbreak V\_\allowbreak INSTALL. Continue the distinct retained-defect, period-height and normalized spectral producers, and the original compatible second-layer selection on one common arithmetic/geometric family. GLOBAL\_\allowbreak LOCAL\_\allowbreak GAIN remains fully closed at its original v0.49 scope. Complete local comparison theorems do not themselves construct the whole common family. Common design existence, NBUILT compatibility, original true-column records, shared fields/embeddings, Fits and eta agreement, simultaneous rows and eventual-family thresholds remain explicit. A stronger historical WPt/TSel construction may serve other producers but is not needed to reopen the proved selected-class local-gain estimate.

\emph{Source.} \href{https://github.com/CocoSgt/riemann-hypothesis-lean/blob/7cbf7a0c32fca1badf05931879a37c3f8a2b03b4/src/Formalism/SharpRestriction/GlobalTheoremInputs.lean#L55-L58}{LinkIRR}.

\emph{Source.} \href{https://github.com/CocoSgt/riemann-hypothesis-lean/blob/7cbf7a0c32fca1badf05931879a37c3f8a2b03b4/src/Development/AbelianVarietyTheory/Proofs/GlobalInterpolation/LinkIRROfJoinEx.lean#L17-L38}{linkIRR\_\allowbreak of\_\allowbreak join\_\allowbreak ex}.

\emph{Source.} \href{https://github.com/CocoSgt/riemann-hypothesis-lean/blob/7cbf7a0c32fca1badf05931879a37c3f8a2b03b4/src/Development/AbelianVarietyTheory/Proofs/GlobalInterpolation/LinkIRROfJoinEx.lean#L143-L166}{link\_\allowbreak irr}.



\paragraph{3. GLOBAL\_\allowbreak PERIOD\_\allowbreak HEIGHT — C — Full actual period and height comparison LinkDIVt}

Closed conditional TuSize -> DGamSize -> (LinkEta -> DIVinN -> LinkDIV), with every original LinkDIV\_\allowbreak at clause.

\emph{Shared objects and choice order.} one FamP; one class $\mathcal A$ and its target-compatible subclass; one coefficient field per family datum; one x and i at each application; original selected rows, words and levels; thresholds chosen before the common family; same RR target; same Bu/besS; same Nu and Bk

\emph{Source.} \href{https://github.com/CocoSgt/riemann-hypothesis-lean/blob/7cbf7a0c32fca1badf05931879a37c3f8a2b03b4/src/Formalism/SharpRestriction/GlobalTheoremInputs.lean#L77-L90}{LinkDIV\_\allowbreak at}.

\emph{Source.} \href{https://github.com/CocoSgt/riemann-hypothesis-lean/blob/7cbf7a0c32fca1badf05931879a37c3f8a2b03b4/src/Formalism/SharpRestriction/GlobalTheoremInputs.lean#L94-L96}{LinkDIVt}.

\emph{Source.} \href{https://github.com/CocoSgt/riemann-hypothesis-lean/blob/7cbf7a0c32fca1badf05931879a37c3f8a2b03b4/src/Development/AbelianVarietyTheory/Proofs/GlobalInterpolation/LinkBrest.lean#L7-L25}{GlobalTheory.Top.link\_\allowbreak brest}.

\emph{Source.} \href{https://github.com/CocoSgt/riemann-hypothesis-lean/blob/7cbf7a0c32fca1badf05931879a37c3f8a2b03b4/src/Development/AbelianVarietyTheory/Proofs/GlobalInterpolation/WireSel.lean#L11-L41}{GlobalTheory.Top.wire\_\allowbreak sel}.



\paragraph{4. GLOBAL\_\allowbreak SPECTRAL — C — Original normalized kernel/Gram estimates LinkS}

Closed TuSize -> LinkS with exact Gamma normalization, off-diagonal Hilbert-Schmidt bound, diagonal error and Dall/Dstar estimate, errors exp(-x\textasciicircum{}(epsilon/8)). CURRENT VERSION0.48 QUALIFICATION: Complete the original same-family spectral/Bessel and Siegel-Weil inputs with all measure, character, transport and common-object ties. The complete MExists measured block and prescribed-measure covolume normalization are supplied; these are not the full normalized spectral theorem. Completed current producers or conditional consumers: ASG1-ASG18, AMD1-AMD17, GCF1-GCF18 at their stated scopes. Earlier frontier descriptions are historical; this qualification is the current scope authority. Original whole-contract premises and quantifiers are retained. CURRENT VERSION0.49 QUALIFICATION: Complete the original same-family spectral/Bessel and Siegel-Weil inputs with every measure, character, transport and common-object tie. The normalized packet and decay are now constructed conditional on the exact BesselNVGap, and the analytic setting uses the original supplied Forms0 scope. These do not imply the full normalized spectral theorem. Completed current producers or conditional consumers: ASG1-ASG18, AMD1-AMD17, GCF1-GCF18 at their stated scopes, BWP1-BWP24 at their exact BesselNVGap/Forms0 scope, WRS1-WRS10 original local gain. Earlier frontier descriptions are historical; this qualification is the current scope authority. Original whole-contract premises and quantifiers are retained.

\emph{Shared objects and choice order.} one FamP; one class $\mathcal A$ and its target-compatible subclass; one coefficient field per family datum; one x and i at each application; original selected rows, words and levels; thresholds chosen before the common family; actual KD; same S.Cu spaces; GamOf; Dstar; Dall

\emph{Completed inputs (version 0.57).} PGR full universal conditional KExpansionC and exact period/operator join, RCV1–RCV25 (full original reference unipotent CovolRef with exact discriminant and independent words), SVW1–SVW18 (full conditional stabilizer volume and actual InVolumes join on unchanged GMR19 output), AOS1–AOS18 (full exact archimedean orbital scalar, gauge, compact volume and all-section integration clause), GMR1-GMR19 (full global/local measure suppliers and RescaleStep3 on the actual setting), BQS1-BQS23 (full QuotP/QuotM, StepStatement and BesselNVGap; setting at Forms0), BST1-BST11, BSE1-BSE9, BNV1-BNV23, BSP1-BSP10, BWP1-BWP24 at original Forms0 scope, AMD1-AMD24

\emph{Remaining producer scope.} PGR proves the full original KExpansionC at its literal universal analytic premises: FoldOK and InHH for both words, FundExists, OOIntAll, OOFundIndep, AmpWellDef and AbsConvT. The support-cell and rational-return dictionaries, finite phase, integral Levi splitting, absolute majorant before unfolding and the actual orbital operator sum preserve the original measures, representative systems, character and column normalization. The resulting period/Gram and operator-compression conclusions retain their printed original inputs; existence of an automorphic family and uniform spectral estimates are not supplied by that identity. The next bounded analytic producer is full InHHC: identify the actual scalar convolution representatives with the original Hilbert-space operators and prove the final automorphic-range equation, reusing PGR's convergence, L2 and reproduction results. The KExpTied join must then receive FoldOKC, InHHC, FundExistsC, OOIntAllC, OOFundIndepC, AmpWellDefC and AbsConvC on the same setting and both words. AbsConv retains operator-norm and auxiliary-vector summability and positive Gam1; scalar integrated convergence alone is insufficient. RCV already supplies full CovolRef for independent words and every original domain. The independent measure tasks are broad JacGauge and the original PlaceSplit/ProdExact consumers, producing RealBracket on those same measures. SVW conditional VolStabFromSWS and actual InVolumes remain complete at their original SWa1 and StabilizerFirstChernCompatibility premises; those premises, finite-place factors and other mass/index/column/form/counting fields remain. BQS full BesselNVGap is available. Forms0, all-family Witt/genus/level/intertwiner Transport/LevelMatch, Bessel uniqueness, AutTie/RowDictTie/LabTie, simultaneous arithmetic cylinder counts and normalized GLOBAL\_\allowbreak SPECTRAL remain.

\emph{Source.} \href{https://github.com/CocoSgt/riemann-hypothesis-lean/blob/7cbf7a0c32fca1badf05931879a37c3f8a2b03b4/src/Formalism/SharpRestriction/GlobalTheoremInputs.lean#L145-L149}{LinkS\_\allowbreak at}.

\emph{Source.} \href{https://github.com/CocoSgt/riemann-hypothesis-lean/blob/7cbf7a0c32fca1badf05931879a37c3f8a2b03b4/src/Formalism/SharpRestriction/GlobalTheoremInputs.lean#L152-L157}{LinkS}.

\emph{Source.} \href{https://github.com/CocoSgt/riemann-hypothesis-lean/blob/7cbf7a0c32fca1badf05931879a37c3f8a2b03b4/src/Development/AbelianVarietyTheory/Proofs/GlobalInterpolation/LinkAlt.lean#L16-L50}{GlobalTheory.Top.link\_\allowbreak alt}.



\paragraph{5. AUT\_\allowbreak UNRAM\_\allowbreak COEF — C — Actual unramified-coefficient automorphy}

For each actual Dat with coefficient ramification one and fixed iota, choose one marked g1 supplying FinA, AutB and LemW together. CURRENT VERSION0.48 QUALIFICATION: Construct the actual AutomorphyUnramifiedCoefficients package: for every original Dat with coefficient ramification one and fixed iota, choose one marked g1 that supplies FinA, AutB and LemW together. The original InpZ record, full RigII, CNR3, finite PropGK consumer and now the entire OneChar implication are supplied. Continue RH Cases I/V and the independent additional Dirichlet-GRH CaseIII/promodularity, SeedIV, tensor-factor and rank-four-chain tasks. Completed current producers or conditional consumers: OCH1-OCH30 (the full consumer implication), KIZ original InpZ, KTR1-KTR13, RNC16, PGK1-PGK6. Earlier frontier descriptions are historical; this qualification is the current scope authority. Original whole-contract premises and quantifiers are retained. CURRENT VERSION0.49 QUALIFICATION: Construct one common GdEnd for each original datum and fixed iota: the same marked g1 must supply FinEnd, AutEnd3 and AutEnd2. UCE now proves all ordinary local/curve/coefficient consumers and the exact implication to FinA/AutB/LemW, including arbitrary entry maps and every original outer-corner extension. The producer still needs the actual common abelian-CM GdFrameEx and LinkedEndE with both rank-two and rank-three runs, full ModUAlt/PfromAut and the two actual seeds, followed by the original local/automorphic descent. At this pin the alternative upstream route gives GdEndAlt; its additional local/descent fields must be matched to the actual original points, so no name-only GdEndAlt-to-GdEnd identification is allowed. Keep the supplied InpZ, RigII, CNR3, finite PropGK and whole OneChar implication. Continue RH Cases I/V and the independent additional Dirichlet-GRH CaseIII/promodularity, SeedIV, tensor-factor and rank-four-chain tasks. Completed current producers or conditional consumers: UCE1-UCE19 (complete GdEnd consumer), OCH1-OCH30, KIZ original InpZ, KTR1-KTR13, RNC16, PGK1-PGK6. Earlier frontier descriptions are historical; this qualification is the current scope authority. Original whole-contract premises and quantifiers are retained.

\emph{Completed inputs (version 0.57).} ACC1–ACC27 full rank-two component theorem and general dimension/incidence suppliers at exact original scope, APD1–APD25 (full broad point/universal maps, actual obstruction and tangent construction, exact original pres), ALC1–ALC22 (full actual local existence, completion, flatness, residual types and representation/component charts), RTC1–RTC25 (upstream universal coefficient rings, full preconnected\_\allowbreak rank\_\allowbreak two and FinPtDown), ACP1-ACP19 (full FinPtOf and abstract propagation principles; ACC now supplies the rank-two component inputs), AFE1-AFE21 (full Gen entries, conditional frame joins, NRep descent and FinPt-to-FinEnd), LFC1-LFC23, CAF1-CAF11, MFC1-MFC14, UCE1-UCE19 (full original GdEnd consumer), OCH1-OCH30, KIZ original InpZ, KTR1-KTR13, RNC16, PGK1-PGK6

\emph{Remaining producer scope.} ACC proves full original c1ch from complete-local connectedness and transports the exact APD presentation to component chains at both ranks. It supplies full loc\_\allowbreak dim, fibre, ab\_\allowbreak gen, l52, NicePrimeEx, min\_\allowbreak hom, jac and l41. The original rank-two sharp2 and PropagationTheoremGL2 are now complete at their literal FinLam, P1-prime and SeedIn premises. The rank-two triangular bound, strict bad-locus bounds, Nice-curve application and seedpt are proved as rank-two specializations; the weaker seed inclusion is preserved, and one closed point is chosen before every original local coefficient extension. The original automorphic P1-prime and seed producers remain hypotheses, not consequences of connectedness. Next are rank three's separate EPres parity data and Sharp theorem, the remaining triangular bound, the rank-two-to-three corner and seed incidence, the exceptional/thick chain installation, rank-three propagation and OffFin. Ordinary APD pres and the now-complete rank-two theorem do not replace these extra rank-three contracts. Existence of originally supplied universal spaces, actual automorphic ModIn/ModU, patching/modularity, classicality, arithmetic P and stronger Alt WitA/WitAQ outputs remain separate. Keep the common marked frame, rank-specific rings/ideals, all coefficient maps and both actual O-prime extensions. Host ell=3 remains included. No upstream FinA conclusion is used.

\emph{Source.} \href{https://github.com/CocoSgt/riemann-hypothesis-lean/blob/7cbf7a0c32fca1badf05931879a37c3f8a2b03b4/src/Development/AutomorphicArithmetic/AutomorphicVarieties/AutomorphicCases/AutomorphyUnramifiedCoefficients.lean#L24-L27}{AutomorphyUnramifiedCoefficients}.

\emph{Source.} \href{https://github.com/CocoSgt/riemann-hypothesis-lean/blob/7cbf7a0c32fca1badf05931879a37c3f8a2b03b4/src/Development/AutomorphicArithmetic/AutomorphicVarieties/AutomorphicCases/AutomorphyUnramifiedCoefficients.lean#L12-L13}{FinA}.

\emph{Source.} \href{https://github.com/CocoSgt/riemann-hypothesis-lean/blob/7cbf7a0c32fca1badf05931879a37c3f8a2b03b4/src/Development/AutomorphicArithmetic/AutomorphicVarieties/AutomorphicCases/AutomorphyUnramifiedCoefficients.lean#L15-L17}{AutB}.

\emph{Source.} \href{https://github.com/CocoSgt/riemann-hypothesis-lean/blob/7cbf7a0c32fca1badf05931879a37c3f8a2b03b4/src/Development/AutomorphicArithmetic/AutomorphicVarieties/AutomorphicCases/AutomorphyUnramifiedCoefficients.lean#L19-L22}{LemW}.



\paragraph{6. TASK\_\allowbreak CASE3\_\allowbreak AUT\_\allowbreak JOIN — B — Complete the original Case-III automorphy/image join}

Install the actual Case-III ProModAlt, SeedIVAlt and shared PMExAlt with the inherited CohomologyHyp/Ferrero-Washington/HFit/SeedII facts into the original whole Case-III automorphy/image conclusion.

\emph{Source.} \href{https://github.com/CocoSgt/riemann-hypothesis-lean/blob/7cbf7a0c32fca1badf05931879a37c3f8a2b03b4/src/Development/AutomorphicArithmetic/AutomorphicVarieties/ModularityArgument/ThresholdArguments/DenseThresholdCases/PromodularityCaseThree.lean#L43-L44}{AutomorphyRankFourDense}.

\emph{Source.} \href{https://github.com/CocoSgt/riemann-hypothesis-lean/blob/7cbf7a0c32fca1badf05931879a37c3f8a2b03b4/src/Development/AutomorphicArithmetic/AutomorphicVarieties/ModularityArgument/ThresholdArguments/DenseThresholdCases/CuFinal.lean#L44-L44}{main4\_\allowbreak alt}.



\paragraph{7. TASK\_\allowbreak CASE3\_\allowbreak PROMOD\_\allowbreak JOIN — C — Reconstruct the actual Case-III chain-to-promodularity proof}

Use the original tensor-factor tower, ChainRankFour, ChainDesc4, seeded rank-four chain reduction, actual PMExAlt and SeedIVAlt to construct ProModAlt on every original Ctx3 star point at the same fixed iota, preserving all finite-coefficient, inertia and support maps.

\emph{Source.} \href{https://github.com/CocoSgt/riemann-hypothesis-lean/blob/7cbf7a0c32fca1badf05931879a37c3f8a2b03b4/src/Development/AutomorphicArithmetic/AutomorphicVarieties/ModularityArgument/ThresholdArguments/DenseThresholdCases/PromodularityCaseThreeOfLemmas.lean#L131-L153}{UpperTheory.V.Cu.promodularity\_\allowbreak case\_\allowbreak three\_\allowbreak of\_\allowbreak chain\_\allowbreak reduction}.

\emph{Source.} \href{https://github.com/CocoSgt/riemann-hypothesis-lean/blob/7cbf7a0c32fca1badf05931879a37c3f8a2b03b4/src/Development/AutomorphicArithmetic/AutomorphicVarieties/ModularityArgument/ThresholdArguments/DenseThresholdCases/PromodularityCaseThree.lean#L40-L42}{PromodularityCaseThreeDense}.



\paragraph{8. SEED\_\allowbreak IV\_\allowbreak ALT — C — Actual Case-III seed with prechosen residual datum}

Base Ctx3/Four tower, iota and Amb -> w,s before every later Dies tower -> compatible finite-coefficient seed with exact residual support/root entries.

\emph{Source.} \href{https://github.com/CocoSgt/riemann-hypothesis-lean/blob/7cbf7a0c32fca1badf05931879a37c3f8a2b03b4/src/Development/AutomorphicArithmetic/AutomorphicVarieties/ModularityArgument/ThresholdArguments/DenseThresholdCases/PromodularityCaseThree.lean#L19-L25}{SeedIVAlt}.



\paragraph{9. TENSOR\_\allowbreak FACTOR\_\allowbreak TOWER — B — Actual tensor-factor tower and characters}

Original finite-field, tame, factor and character constraints; source-level construction is sizeable and has not been reconstructed here.

\emph{Source.} \href{https://github.com/CocoSgt/riemann-hypothesis-lean/blob/7cbf7a0c32fca1badf05931879a37c3f8a2b03b4/src/Development/AutomorphicArithmetic/AutomorphicVarieties/ModularityArgument/ThresholdArguments/DenseThresholdCases/PromodularityCaseThreeOfLemmas.lean#L131-L153}{UpperTheory.V.Cu.promodularity\_\allowbreak case\_\allowbreak three\_\allowbreak of\_\allowbreak chain\_\allowbreak reduction}.



\paragraph{10. CHAIN\_\allowbreak RANK\_\allowbreak FOUR — C — Actual SomeIn-to-AllIn propagation}

Same source Space and PM ideal, residual point Good4, original bound 30(C+1)\textasciicircum{}2 <= d; geometry must preserve all identities.

\emph{Source.} \href{https://github.com/CocoSgt/riemann-hypothesis-lean/blob/7cbf7a0c32fca1badf05931879a37c3f8a2b03b4/src/Development/AutomorphicArithmetic/AutomorphicVarieties/ModularityArgument/ThresholdArguments/DenseThresholdCases/PromodularityCaseThree.lean#L26-L32}{ChainRankFour}.



\paragraph{11. CHAIN\_\allowbreak DESC4 — B — Original residual chain descent}

Finite residual descent retaining source Chain4 compatibility.

\emph{Source.} \href{https://github.com/CocoSgt/riemann-hypothesis-lean/blob/7cbf7a0c32fca1badf05931879a37c3f8a2b03b4/src/Development/AutomorphicArithmetic/AutomorphicVarieties/ModularityArgument/ThresholdArguments/DenseThresholdCases/PromodularityCaseThreeOfLemmas.lean#L131-L153}{UpperTheory.V.Cu.promodularity\_\allowbreak case\_\allowbreak three\_\allowbreak of\_\allowbreak chain\_\allowbreak reduction}.



\paragraph{12. RANK4\_\allowbreak CHAIN\_\allowbreak REDUCTION — C — Seeded reduction to the original residual chain}

Select T1 before every later Deep tower; retain one compatible seed and Chain4 for every original residual lift over Q.

\emph{Source.} \href{https://github.com/CocoSgt/riemann-hypothesis-lean/blob/7cbf7a0c32fca1badf05931879a37c3f8a2b03b4/src/Development/AutomorphicArithmetic/AutomorphicVarieties/ModularityArgument/ThresholdArguments/DenseThresholdCases/PromodularityCaseThree.lean#L33-L39}{RankFourChainReductionDense}.

\emph{Source.} \href{https://github.com/CocoSgt/riemann-hypothesis-lean/blob/7cbf7a0c32fca1badf05931879a37c3f8a2b03b4/src/Development/AutomorphicArithmetic/AutomorphicVarieties/ModularityArgument/ThresholdArguments/DenseThresholdCases/CuPChainIV.lean#L16-L82}{UpperTheory.V.Cu.chainIV\_\allowbreak alt}.



\paragraph{13. TASK\_\allowbreak COUNT\_\allowbreak V\_\allowbreak INSTALL — B — Complete the original CountV finite assembly}

From exactly the listed source host/cohomology inputs, complete countV\_\allowbreak of\_\allowbreak finiteness\_\allowbreak duality\_\allowbreak euler\_\allowbreak host, including original host/frame formation, normalized length, the CaseV guard and R1<=2\textasciicircum{}c2.

\emph{Completed inputs (version 0.57).} ACC full original rank-two propagation; automorphic seed and rank-three producers retained, APD full original presentation on the same arithmetic local objects, now used in ACC component chains, ALC actual local objects and charts, now used in ACC dimension and rank-two seed/propagation arguments, RTC full rank-two preconnected core and coefficient descent, used in ACC; actual counting producers retained, ACP full finite-point and abstract propagation implications, used in ACC; rank-three and counting producers retained, AFE full Host entry/auxiliary frame/local witnesses including ell=3, LFC1-LFC23, CAF9-CAF11, MFC1-MFC14, Inherited CountV consumers

\emph{Remaining producer scope.} ACC proves full original c1ch from complete-local connectedness and transports the exact APD presentation to component chains at both ranks. It supplies full loc\_\allowbreak dim, fibre, ab\_\allowbreak gen, l52, NicePrimeEx, min\_\allowbreak hom, jac and l41. The original rank-two sharp2 and PropagationTheoremGL2 are now complete at their literal FinLam, P1-prime and SeedIn premises. The rank-two triangular bound, strict bad-locus bounds, Nice-curve application and seedpt are proved as rank-two specializations; the weaker seed inclusion is preserved, and one closed point is chosen before every original local coefficient extension. The original automorphic P1-prime and seed producers remain hypotheses, not consequences of connectedness. Next are rank three's separate EPres parity data and Sharp theorem, the remaining triangular bound, the rank-two-to-three corner and seed incidence, the exceptional/thick chain installation, rank-three propagation and OffFin. Ordinary APD pres and the now-complete rank-two theorem do not replace these extra rank-three contracts. Existence of originally supplied universal spaces, actual automorphic ModIn/ModU, patching/modularity, classicality, arithmetic P and stronger Alt WitA/WitAQ outputs remain separate. Keep the common marked frame, rank-specific rings/ideals, all coefficient maps and both actual O-prime extensions. Host ell=3 remains included. No upstream FinA conclusion is used. Complete the original F3 local coefficient bridges and F4 common-overfield readings, with valuation/norm/GoodCoef/iq and the entire coefficient field, and actual complex Hecke standard/no-Jordan plus arithmetic character data on the same form datum. Construct the actual base rows, host/frame lengths, nonzero chains and group ties, MultEq support, eigensystem/carrier/parts, all-embedding logLam, HD/d3b and NEb. Keep the stronger second DivRows3 scope distinct from the completed first DivDef3 and row-pair integrability. The full arithmetic entry and conditional common-frame construction do not by themselves install HOST\_\allowbreak NUMBER\_\allowbreak S, HOST\_\allowbreak FRAME\_\allowbreak LENGTH, TASK\_\allowbreak HOST\_\allowbreak VAN\_\allowbreak E\_\allowbreak INSTALL, TASK\_\allowbreak HOST\_\allowbreak XPP\_\allowbreak INSTALL or TASK\_\allowbreak COUNT\_\allowbreak V\_\allowbreak INSTALL.

\emph{Source.} \href{https://github.com/CocoSgt/riemann-hypothesis-lean/blob/7cbf7a0c32fca1badf05931879a37c3f8a2b03b4/src/Development/AutomorphicArithmetic/AutomorphicVarieties/HeightAndVanishing/Counts.lean#L69-L69}{CountV}.

\emph{Source.} \href{https://github.com/CocoSgt/riemann-hypothesis-lean/blob/7cbf7a0c32fca1badf05931879a37c3f8a2b03b4/src/Development/AutomorphicArithmetic/AutomorphicVarieties/ModularityArgument/ThresholdArguments/DenseThresholdCases/CountVOf.lean#L33-L60}{UpperTheory.V.UnitaryThreeInertia.countV\_\allowbreak of\_\allowbreak finiteness\_\allowbreak duality\_\allowbreak euler\_\allowbreak host}.



\paragraph{14. THM\_\allowbreak S6\_\allowbreak S — C — Original unitary-three ordinary support and automorphy}

Original host/frame and fixed iota -> one marked g1 -> every universal R,u -> weight-zero finiteness and every Type-I BPrime\_\allowbreak s point.

\emph{Source.} \href{https://github.com/CocoSgt/riemann-hypothesis-lean/blob/7cbf7a0c32fca1badf05931879a37c3f8a2b03b4/src/Development/AutomorphicArithmetic/AutomorphicVarieties/ModularityArgument/ThresholdArguments/DenseThresholdCases/ThmS6CoreS.lean#L17-L26}{ThmS6\_\allowbreak s}.

\emph{Source.} \href{https://github.com/CocoSgt/riemann-hypothesis-lean/blob/7cbf7a0c32fca1badf05931879a37c3f8a2b03b4/src/Development/AutomorphicArithmetic/AutomorphicVarieties/ModularityArgument/ThresholdArguments/DenseThresholdCases/ThmS6CoreS.lean#L12-L16}{BPrime\_\allowbreak s}.



\paragraph{15. HOST\_\allowbreak FRAME\_\allowbreak LENGTH — B — Original Case-V host/frame and normalized-length installation}

Actual HostData, frame, local-condition equality, abstract count and normalized original Selmer length must all refer to the same objects.

\emph{Completed inputs (version 0.57).} AFE/LFC full entry and conditional common frame, CAF and MFC, DNE weight/host-degree dictionary at supplied tables

\emph{Remaining producer scope.} Complete the original F3 local coefficient bridges and F4 common-overfield readings, with valuation/norm/GoodCoef/iq and the entire coefficient field, and actual complex Hecke standard/no-Jordan plus arithmetic character data on the same form datum. Construct the actual base rows, host/frame lengths, nonzero chains and group ties, MultEq support, eigensystem/carrier/parts, all-embedding logLam, HD/d3b and NEb. Keep the stronger second DivRows3 scope distinct from the completed first DivDef3 and row-pair integrability. The full arithmetic entry and conditional common-frame construction do not by themselves install HOST\_\allowbreak NUMBER\_\allowbreak S, HOST\_\allowbreak FRAME\_\allowbreak LENGTH, TASK\_\allowbreak HOST\_\allowbreak VAN\_\allowbreak E\_\allowbreak INSTALL, TASK\_\allowbreak HOST\_\allowbreak XPP\_\allowbreak INSTALL or TASK\_\allowbreak COUNT\_\allowbreak V\_\allowbreak INSTALL.

\emph{Source.} \href{https://github.com/CocoSgt/riemann-hypothesis-lean/blob/7cbf7a0c32fca1badf05931879a37c3f8a2b03b4/src/Development/AutomorphicArithmetic/AutomorphicVarieties/ModularityArgument/ThresholdArguments/DenseThresholdCases/CountVOf.lean#L33-L60}{UpperTheory.V.UnitaryThreeInertia.countV\_\allowbreak of\_\allowbreak finiteness\_\allowbreak duality\_\allowbreak euler\_\allowbreak host}.



\paragraph{16. HOST\_\allowbreak NUMBER\_\allowbreak S — C — Uniform original-host rank-three catalogue}

One exponent before all closures, coefficient rings, levels, characters, frames, embeddings and marked points; one separating trace function controls all later finite pairwise nonisomorphic collections.

\emph{Completed inputs (version 0.57).} AFE full Host entry at original scope, LFC full conditional common frame, CAF and MFC, Inherited host-count consumers

\emph{Remaining producer scope.} Complete the original F3 local coefficient bridges and F4 common-overfield readings, with valuation/norm/GoodCoef/iq and the entire coefficient field, and actual complex Hecke standard/no-Jordan plus arithmetic character data on the same form datum. Construct the actual base rows, host/frame lengths, nonzero chains and group ties, MultEq support, eigensystem/carrier/parts, all-embedding logLam, HD/d3b and NEb. Keep the stronger second DivRows3 scope distinct from the completed first DivDef3 and row-pair integrability. The full arithmetic entry and conditional common-frame construction do not by themselves install HOST\_\allowbreak NUMBER\_\allowbreak S, HOST\_\allowbreak FRAME\_\allowbreak LENGTH, TASK\_\allowbreak HOST\_\allowbreak VAN\_\allowbreak E\_\allowbreak INSTALL, TASK\_\allowbreak HOST\_\allowbreak XPP\_\allowbreak INSTALL or TASK\_\allowbreak COUNT\_\allowbreak V\_\allowbreak INSTALL.

\emph{Source.} \href{https://github.com/CocoSgt/riemann-hypothesis-lean/blob/7cbf7a0c32fca1badf05931879a37c3f8a2b03b4/src/Development/AutomorphicArithmetic/AutomorphicVarieties/ModularityArgument/ThresholdArguments/DenseThresholdCases/HostNumberS.lean#L18-L28}{HostNumber\_\allowbreak s}.



\paragraph{17. TASK\_\allowbreak HOST\_\allowbreak VAN\_\allowbreak E\_\allowbreak INSTALL — B — Transport K/duality to literal HostVanE}

Use the already supplied LemmaK/local-global duality and finiteness on the actual CountV host/frame and exact original coefficient modules.

\emph{Source.} \href{https://github.com/CocoSgt/riemann-hypothesis-lean/blob/7cbf7a0c32fca1badf05931879a37c3f8a2b03b4/src/Development/AutomorphicArithmetic/AutomorphicVarieties/UnitaryThreeDeformations/HostNumber.lean#L9-L12}{HostVanE}.

\emph{Source.} \href{https://github.com/CocoSgt/riemann-hypothesis-lean/blob/7cbf7a0c32fca1badf05931879a37c3f8a2b03b4/src/Development/AutomorphicArithmetic/AutomorphicVarieties/LevelCohomology/SoHostVanE.lean#L11-L18}{UpperTheory.V.X.So.hostVanE}.



\paragraph{18. TASK\_\allowbreak HOST\_\allowbreak XPP\_\allowbreak INSTALL — B — Transport CorX to literal HostXpp}

Preserve Psi\textasciicircum{}-2 eta, the prime-to-ell quotient, actual local conditions and the original CountV host while installing BR7.

\emph{Source.} \href{https://github.com/CocoSgt/riemann-hypothesis-lean/blob/7cbf7a0c32fca1badf05931879a37c3f8a2b03b4/src/Development/AutomorphicArithmetic/AutomorphicVarieties/UnitaryThreeDeformations/HostNumber.lean#L14-L17}{HostXpp}.

\emph{Source.} \href{https://github.com/CocoSgt/riemann-hypothesis-lean/blob/7cbf7a0c32fca1badf05931879a37c3f8a2b03b4/src/Development/AutomorphicArithmetic/AutomorphicVarieties/LevelCohomology/Xpp.lean#L11-L84}{UpperTheory.V.X.xpp}.



\paragraph{19. DESIGN\_\allowbreak CHOICES\_\allowbreak COMMON — C — Construct the original common design family and bundle}

For the one target HostExW, construct the original compatible DesignChoices on Cdes=WithTw(WithSm Cprime), where Cprime=WithChN(A intersect target class), A is the one Built construction class, and cT(eps,y)=cTmin(y). Retain every original field, tuple, word, metric and common bundle duty. DSZ10 and RMP11 give size/rank properties of class members, not member existence. The original RecExP/RecAtG and common object/row route remains coupled. Construct ONE original common design family R and its recursive/geometric half with original Inputs/Pins/Holds/ExAt/InGeo and required CoreP. PKT now reconstructs full exists\_\allowbreak pkT, rowPk\_\allowbreak nonempty and exists\_\allowbreak moved with Copy/Coh on those original given inputs and a supplied fitting later stage: one catalogue, both ChT, complete LfT, actual integral lattices and all metric/measure/period constants are retained; one numerical record precedes every eta and V/S/N/G. RDR supplies one Side0 before all readers, ReadPk/pins and full ReadL with exact o.eta=q.eta. The initial actual geometry, Fits/division/central conditions, primitive rational tests and eta agreement must still be produced. Inherited CSR supplies finite columns, one ReadC before later twists, period/evaluation and scalar/error transports, RowSideDetermined and same-R Bundle/BTw/BWit/BunU consumers. RPI supplies exact RowPairInt for actual pairs and finite translates with original mu\_\allowbreak b/fd0; first literal DivDef3 at the same updated R,F,Y,Z,C,sEE follows from its retained original readings, Pins/Holds, stored nonzero IsItW chain, five real-group predicates and group tie. Remaining common fields include host/frame, MultEq nonzero-part support, that original nonzero chain/group tie, eigensystem/carrier/parts and all-embedding logLam, HD/d3b, NEb, actual rational tests and base rows. Retain every simultaneous family hypothesis and eventual threshold. Stronger second DivRows3 remains separate; universal quotient-height CuspRapidDecay and actual common member existence are not inferred from the comparison constructors. Version 0.41 supplies the finite phase-correct rational theta route: exact coefficient extension, split representatives, square-law/PhiIs, Pins from CoversAt, class-sum/point-identity transfer to OneMapOverL, actual-good-projector RatG and both-level LiftTestB. Its original higher covers, all-tuple class-sum uniqueness including phase rescaling, one PointIdentCore before O, base ClassSumsSpan and geometric LatFull remain explicit producers. MLat is available from the original LatFull input. This sharpens the rational-test work without claiming actual common geometry, Fits or eta agreement. Version0.46 proves finite generation of the original whole integral cusp lattice by actual integral-frame evaluation. It also proves the exact identity-cover, analytic unfolding, original-field coefficient descent, global clearing-scalar and both-level lattice-transport implications. Full spanning retains the supplied original IsNab connection and the five explicit principal-chart, level-system, raw-Gaussian, short-rule/dictionary and ramified-jet producers. The higher tuple covers, actual common class sums, PointIdentCore and ClassSumsSpan remain separate. CURRENT VERSION0.47 QUALIFICATION: Construct one original common geometric family and the higher CoversAt/all-tuple class sums and uniqueness with original phase rescalings, one PointIdentCore with common lambda,z, and base ClassSumsSpan. CNC now supplies the full original IsNab; PCL supplies the conditional trace-dual principal-lattice and quantitative normalization. Actual uniformization/HasPrinc and the five principal-chart, finite-level-system, Gaussian/reproduction, short-rule/simultaneous-dictionary and ramified-jet producer groups remain. Keep all original both-level, coefficient-field, Fits and eta agreements. Completed new original producers: ICL1-ICL21, CNC1-CNC16, PCL1-PCL12. Any earlier narrative naming these completed producers as outstanding is historical; this qualification is the current scope authority. CURRENT VERSION0.48 QUALIFICATION: Construct the original finite holomorphic FamilyUnif/N1c, normalized theta families and local lifting, then the common principal charts and finite-level theta systems. PUN now supplies the exact fixed-space models, all-family center lattices, A=1 exponent matching and individual primitive points with the full point-equality criterion; CNC supplies the full canonical connection and PCL the quantitative trace-dual normalization. Complete higher CoversAt and all-tuple class sums/uniqueness with the original phase rescalings, one PointIdentCore with common lambda,z, and base ClassSumsSpan. The Gaussian/reproduction, short-rule/simultaneous-dictionary and ramified-jet producers, and all original both-level, coefficient-field, Fits and eta agreements remain explicit. Completed current producers or conditional consumers: ICL1-ICL21, CNC1-CNC16, PCL1-PCL12, PUN1-PUN24, GCF1-GCF18 at their stated scopes. Earlier frontier descriptions are historical; this qualification is the current scope authority. Original whole-contract premises and quantifiers are retained. CURRENT VERSION0.49 QUALIFICATION: Construct the adjoining simultaneous principal framed points and test-ring dictionary. Full finite holomorphic FamilyUnif/N1c, normalized theta families and local lifts are now supplied by HUF. The next bounded target is allTest\_\allowbreak of\_\allowbreak secBaseChange: represent FrPt by the Lie-eigenframe bundles and theta-section frame line, use the full n=1 SecBaseChangeSt on every nonflat coefficient base change, and prove finite-generated subalgebra descent with the exact TRing/FromTest transport. Then supply the full FramesOK/CoverTri/FibTri and IsUnf/hasPrinc package, including adapted mod-eight symplectic basis, normalized theta multiplier/weight and arithmetic invariance. Continue higher CoversAt, all-tuple class sums and uniqueness with original phase rescalings, one PointIdentCore with common lambda,z, base ClassSumsSpan, Gaussian/reproduction, short-rule dictionaries and ramified jets, retaining both-level, coefficient-field, Fits and eta agreements. Common-family nonemptiness remains distinct from compatibility of an already supplied member. Completed current producers or conditional consumers: ICL1-ICL21, CNC1-CNC16, PCL1-PCL12, PUN1-PUN24, HUF1-HUF21 (full original FamilyUnif/N1c), GCF1-GCF18 and WRS1-WRS10 at their stated scopes. Earlier frontier descriptions are historical; this qualification is the current scope authority. Original whole-contract premises and quantifiers are retained.

\emph{Completed inputs (version 0.57).} SGV full original section growth, DJW full determinant window and actual integral row construction at original normal-form input, PBC1–PBC27 (full arbitrary-base polarized pullback, actual very-ampleness and section finiteness), SHV1–SHV21 (both original Serre conclusions, actual finite section model and full all-cover return), ASL1–ASL21 (actual section localization and full original arbitrary-base Perf producer), FCD1-FCD18 (full homogeneous-chart descent, graded presentation and Noetherian cube model), GSL1-GSL26 (full conditional section laws, universal basis and projective-model algebra), NFL1-NFL22, TDR1-TDR13, NFR1-NFR13, HUF1-HUF21, PUN1-PUN24, ICL1-ICL21, CNC1-CNC16, PCL1-PCL12, GCF/WRS at original supplied-member scope

\emph{Remaining producer scope.} PBC supplies full arbitrary-base polarized pullback, actual very-ampleness and every-integer-power SecM finiteness. SGV now supplies the full original secM\_\allowbreak finite\_\allowbreak and\_\allowbreak growth, hence GSL\_\allowbreak FINITE\_\allowbreak GROWTH: on the literal identity-base product there are fixed integers N>=1 and d>=1 with dim SecM(N*t)=d*t\textasciicircum{}g for every t>=1. Its explicit classical inputs are Milne Proposition 9.1 and Theorem 13.3(b),(c); the section-coordinate and scalar dictionaries, ampleness and geometric-integrality/dimension hypotheses are proved on the original object in every characteristic. DJW supplies full ThmP, arbitrary-index BiValue, BiContent/BiWindow, all four MtoHM fields and full block\_\allowbreak of\_\allowbreak normal\_\allowbreak form for every supplied restriction identity. It constructs actual integral germs with the exact unit determinant factor, and obtains the original row window from the original MConcl. The original RowData.LemmaM existence producer remains: adapted Cartier chains, residue ranks, pairing-kernel estimate and weighted remainder must construct that normal form on the given row data. Remaining for the geometric section laws are multiplication-pullback generated comparison, finite/flat/surjective and etale multiplication geometry, invariant-function descent on every relevant open, the later every-positive-power GSL\_\allowbreak VANISH\_\allowbreak CLOSED conclusion, the required arbitrary-test integrality, principal rank one, torsion-indexed section decomposition, and square/translation generators. SHV17 all-cover eventual Serre vanishing and SHV21 homogeneous-polynomial generation remain complete. The completed identity-base integrality and growth do not themselves supply these different contracts. Original Moduli and HUF/PUN foundations remain. Later finite-level systems, Gaussian/reproduction, simultaneous dictionaries, ramified jets, all-tuple class sums, phase-sensitive uniqueness, common PointIdentCore and base ClassSumsSpan remain, retaining both levels, whole fields, Fits and eta agreement.



\paragraph{20. LOCAL\_\allowbreak NBUILT\_\allowbreak COMMON — C — Actual local specialization on the common selected family}

Construct, with the required closed-class property on the same Des(C,3), one eventual family of actual NBuilt certificates for the original Fm.D(x,i), retaining the certificate before forgetting to TSel. HLP now reconstructs full HeckeLift0, all-index commutation, HeckeStd0 and HeckeIntTw from the original MtInput and complex Hecke facts, in the full ambient GenAt range k<=nV. Every maximal ideal is an actual supported residual character class; literal MtSub is the O-span of projections of the original MtLat generators, is finite free, and localizes onto the exact NCov.part. Literal Psi has its exact range and off-space projection error, both original span assertions hold, and ActEx/ActHK are constructed on the same lattice under original semisimplicity. Original MjLat relation quotients, universal coefficient-map relation fields, selected tame generators, coefficient embeddings and all identity transports give the complete nonempty\_\allowbreak spec constructor on its full printed input list. A single compositum combines the four original coefficient requirements with big fixed beforehand. Whole ndSpec\_\allowbreak exists still needs actual common geometric/choice MtInput, complex HeckeStdPlain and the no-Jordan producer, original F3 bridges with valuation/GoodCoef/iq, original F4 common-overfield reading of the entire coefficient field, and exact arithmetic character/reciprocity/admissibility suppliers. Original AutSpec, Galois specialization and the remaining NBuilt fields on that same installation also remain. Do not replace the actual projected lattice or maps by abstract objects with the same dimensions.

\emph{Completed inputs (version 0.57).} ACC/PGR/DJW/SGV/FET named original suppliers and joins at their explicitly retained inputs, PBC full polarized pullback; APD exact actual presentation; RCV reference covolume; WPR original conditional WachEx, SHV both Serre conclusions; ALC actual arithmetic charts; SVW same-setting volume join, ASL full actual perfect-complex model; RTC common-frame arithmetic inputs; AOS exact unchanged infinity normalization, FCD full algebraic model, ACP abstract arithmetic suppliers, GMR same-setting measure normalization, GSL/NFL full conditional section laws and uniformization, AFE/LFC full entry and conditional common frame, TDR/NFR and CAF/MFC, WRS selected built-class extraction at membership

\emph{Remaining producer scope.} Continue the distinct retained-defect, period-height and normalized spectral producers, and the original compatible second-layer selection on one common arithmetic/geometric family. GLOBAL\_\allowbreak LOCAL\_\allowbreak GAIN remains fully closed at its original v0.49 scope. Complete local comparison theorems do not themselves construct the whole common family. Common design existence, NBUILT compatibility, original true-column records, shared fields/embeddings, Fits and eta agreement, simultaneous rows and eventual-family thresholds remain explicit. A stronger historical WPt/TSel construction may serve other producers but is not needed to reopen the proved selected-class local-gain estimate. PBC supplies full arbitrary-base polarized pullback, actual very-ampleness and every-integer-power SecM finiteness. SGV now supplies the full original secM\_\allowbreak finite\_\allowbreak and\_\allowbreak growth, hence GSL\_\allowbreak FINITE\_\allowbreak GROWTH: on the literal identity-base product there are fixed integers N>=1 and d>=1 with dim SecM(N*t)=d*t\textasciicircum{}g for every t>=1. Its explicit classical inputs are Milne Proposition 9.1 and Theorem 13.3(b),(c); the section-coordinate and scalar dictionaries, ampleness and geometric-integrality/dimension hypotheses are proved on the original object in every characteristic. DJW supplies full ThmP, arbitrary-index BiValue, BiContent/BiWindow, all four MtoHM fields and full block\_\allowbreak of\_\allowbreak normal\_\allowbreak form for every supplied restriction identity. It constructs actual integral germs with the exact unit determinant factor, and obtains the original row window from the original MConcl. The original RowData.LemmaM existence producer remains: adapted Cartier chains, residue ranks, pairing-kernel estimate and weighted remainder must construct that normal form on the given row data. Remaining for the geometric section laws are multiplication-pullback generated comparison, finite/flat/surjective and etale multiplication geometry, invariant-function descent on every relevant open, the later every-positive-power GSL\_\allowbreak VANISH\_\allowbreak CLOSED conclusion, the required arbitrary-test integrality, principal rank one, torsion-indexed section decomposition, and square/translation generators. SHV17 all-cover eventual Serre vanishing and SHV21 homogeneous-polynomial generation remain complete. The completed identity-base integrality and growth do not themselves supply these different contracts. Original Moduli and HUF/PUN foundations remain. Later finite-level systems, Gaussian/reproduction, simultaneous dictionaries, ramified jets, all-tuple class sums, phase-sensitive uniqueness, common PointIdentCore and base ClassSumsSpan remain, retaining both levels, whole fields, Fits and eta agreement.



\paragraph{21. SECOND\_\allowbreak LAYER\_\allowbreak COMPATIBLE\_\allowbreak SELECTION — C — Original second-layer selector and class duties}

For the same A, supply Duty on A and a Closed common original TSel family whose selected records are the original true-column rec7 records (or have the entire dependent tuple transported), with their real witness and original coefficient-field root data. Existence/compatibility remains arithmetic; merely abstractly isomorphic fields or independently selected columns do not suffice. CSR21–22 now construct the same-R Bundle/BTw/BWit/BunU outputs from the exact selected realization/row-side packet and NEb, with the original EvP/ParEv quantifier order. Actual common selection, the five stated realization packets and NEb remain upstream; none is replaced by an independently chosen compatible record. Version 0.38 closes the selected analytic extension and reconstructs the first updated DivDef3 consumer; it leaves the original common selection, geometric readings, rational tests, base rows, host/parts/NEb and exact retained DI premises upstream, with unchanged EvP/ParEv thresholds. Version 0.39 supplies the original given-data reader constructors and exact RowPairInt/first DivDef3 on original premises. It leaves one actual common geometric/comparison/test member, eta agreement, compatible selection, rational tests, base rows, host/parts/NEb and original EvP/ParEv thresholds upstream. Version 0.40 also supplies original comparison/RowPk/fitting-stage Moved algorithms and exact common data transport. The actual initial geometric member, fitting tests, eta agreement and compatible second-layer selector with original eventual thresholds remain upstream. Version 0.41 supplies the finite phase-correct rational theta route: exact coefficient extension, split representatives, square-law/PhiIs, Pins from CoversAt, class-sum/point-identity transfer to OneMapOverL, actual-good-projector RatG and both-level LiftTestB. Its original higher covers, all-tuple class-sum uniqueness including phase rescaling, one PointIdentCore before O, base ClassSumsSpan and geometric LatFull remain explicit producers. MLat is available from the original LatFull input. This sharpens the rational-test work without claiming actual common geometry, Fits or eta agreement. CURRENT VERSION0.49 QUALIFICATION: Continue the distinct global denominator, period/height, normalized spectral and second-layer compatible-selection producers, together with existence of the original common arithmetic/geometric family. GLOBAL\_\allowbreak LOCAL\_\allowbreak GAIN and its actual GroupCompatibilityFamily supplier are now complete on the unchanged selected built class; no additional WPt/TSel realization is needed for that estimate. Do not reopen local first-growth, mean, auxiliary admission, selected D3 tops/Adm2F, the single GIs2 stage or the exact comparison/kernel transport. A full RecT or TSel may remain useful for other tasks, but is not claimed by this local-gain closure. Completed current producers or conditional consumers: GCF1-GCF18 at their stated scopes, WRS1-WRS10 original selected-class compatibility/local gain. Earlier frontier descriptions are historical; this qualification is the current scope authority. Original whole-contract premises and quantifiers are retained.

\emph{Completed inputs (version 0.57).} ACC/PGR/DJW/SGV/FET named original suppliers on unchanged actual objects, PBC/APD/RCV/WPR full named original producers on unchanged actual objects, IFD full invariant matrix descent; SHV/ALC/SVW full named nested producers, WIS actual inclusion and invariant completeness; AOS exact scalar; ASL/RTC full named nested producers, WLD full generic lattice constructors, GMR actual-setting normalization, FCD/ACP exact nested suppliers, WRS1-WRS10 original selected-class compatibility/local gain, DNE full denominator consumer, DFD/SSC canonical-data and signature comparisons, GSL/NFL and AFE/LFC conditional constructors, Inherited exact same-record realization/transport consumers

\emph{Remaining producer scope.} Continue the distinct retained-defect, period-height and normalized spectral producers, and the original compatible second-layer selection on one common arithmetic/geometric family. GLOBAL\_\allowbreak LOCAL\_\allowbreak GAIN remains fully closed at its original v0.49 scope. Complete local comparison theorems do not themselves construct the whole common family. Common design existence, NBUILT compatibility, original true-column records, shared fields/embeddings, Fits and eta agreement, simultaneous rows and eventual-family thresholds remain explicit. A stronger historical WPt/TSel construction may serve other producers but is not needed to reopen the proved selected-class local-gain estimate.



\paragraph{22. BASE\_\allowbreak THETA\_\allowbreak READINGS — C — Actual base reading on the common bundle}

For every family in the original class, eventually on the SAME selected Bundle and for each i<m, supply the actual R,F,Read,ReadL, weak models/Weil data and phase-one pins. Identify their LiftTie.theta, mtRead, mjRead and ThDef3 basis with that Bundle.P.A3, P.Zs, P.S0 and its stored maps. Exact field, embeddings, metrics, measures, row family and weak choices are retained. This remains one part of the coupled BTH12 producer. PKT now supplies the complete original comparison packet, RowPk and fitting-stage Moved with Copy/Coh on original given geometry, Inputs/Pins/Holds/ExAt/InGeo and Fits. Its ChT/LfT/Reps, actual lattice carriers, original Haar and norm constants, period maps, representatives and all-place RrT are constructed together; its one numerical record precedes every eta/V/S/N/G. RDR gives the original level-zero side, moved-Read and ReadPk/pin constructors on those objects, with one Side0 before every reader and all twelve Read ties. Full ReadL retains the exact o.eta(j,v)=q.eta(j,v) guard for every 1<=j<=m and place v; arbitrary Over does not imply it. Inherited CSR fixes one ReadC before all TwArgs, transports both metrics and all thirteen reading maps, and gives actual period/evaluation with kappa=(c*nP)\textasciicircum{}(-1/2). RPI supplies exact row-pair integrability, so first updated DivDef3 follows from the retained original Pins/Holds, nonzero chain, five real-group predicates and group tie. Initial theta extension remains complete. Remaining work is actual initial common geometry and its original predicates, Fits/primitive rational tests, eta agreement and outer family/eventual/same-Bundle duties, reFace compatibility, MultEq support, original chain/group data, actual base rows and any stronger DivRows3 theta square. No new compatible reading, family threshold or Haar normalization is substituted. Version 0.41 supplies the finite phase-correct rational theta route: exact coefficient extension, split representatives, square-law/PhiIs, Pins from CoversAt, class-sum/point-identity transfer to OneMapOverL, actual-good-projector RatG and both-level LiftTestB. Its original higher covers, all-tuple class-sum uniqueness including phase rescaling, one PointIdentCore before O, base ClassSumsSpan and geometric LatFull remain explicit producers. MLat is available from the original LatFull input. This sharpens the rational-test work without claiming actual common geometry, Fits or eta agreement.

\emph{Completed inputs (version 0.57).} SGV full original section growth, DJW full determinant window and actual integral row construction at original normal-form input, PBC1–PBC27 (full arbitrary-base polarized pullback, actual very-ampleness and section finiteness), SHV1–SHV21 (both original Serre conclusions, actual finite section model and full all-cover return), ASL1–ASL21 (actual section localization and full original arbitrary-base Perf producer), FCD1-FCD18 (full homogeneous-chart descent, graded presentation and Noetherian cube model), GSL1-GSL26, NFL1-NFL22, TDR1-TDR13, NFR1-NFR13, HUF1-HUF21, PUN1-PUN24, CNC1-CNC16, Inherited actual-reading consumers at original inputs

\emph{Remaining producer scope.} PBC supplies full arbitrary-base polarized pullback, actual very-ampleness and every-integer-power SecM finiteness. SGV now supplies the full original secM\_\allowbreak finite\_\allowbreak and\_\allowbreak growth, hence GSL\_\allowbreak FINITE\_\allowbreak GROWTH: on the literal identity-base product there are fixed integers N>=1 and d>=1 with dim SecM(N*t)=d*t\textasciicircum{}g for every t>=1. Its explicit classical inputs are Milne Proposition 9.1 and Theorem 13.3(b),(c); the section-coordinate and scalar dictionaries, ampleness and geometric-integrality/dimension hypotheses are proved on the original object in every characteristic. DJW supplies full ThmP, arbitrary-index BiValue, BiContent/BiWindow, all four MtoHM fields and full block\_\allowbreak of\_\allowbreak normal\_\allowbreak form for every supplied restriction identity. It constructs actual integral germs with the exact unit determinant factor, and obtains the original row window from the original MConcl. The original RowData.LemmaM existence producer remains: adapted Cartier chains, residue ranks, pairing-kernel estimate and weighted remainder must construct that normal form on the given row data. Remaining for the geometric section laws are multiplication-pullback generated comparison, finite/flat/surjective and etale multiplication geometry, invariant-function descent on every relevant open, the later every-positive-power GSL\_\allowbreak VANISH\_\allowbreak CLOSED conclusion, the required arbitrary-test integrality, principal rank one, torsion-indexed section decomposition, and square/translation generators. SHV17 all-cover eventual Serre vanishing and SHV21 homogeneous-polynomial generation remain complete. The completed identity-base integrality and growth do not themselves supply these different contracts. Original Moduli and HUF/PUN foundations remain. Later finite-level systems, Gaussian/reproduction, simultaneous dictionaries, ramified jets, all-tuple class sums, phase-sensitive uniqueness, common PointIdentCore and base ClassSumsSpan remain, retaining both levels, whole fields, Fits and eta agreement. Complete the original F3 local coefficient bridges and F4 common-overfield readings, with valuation/norm/GoodCoef/iq and the entire coefficient field, and actual complex Hecke standard/no-Jordan plus arithmetic character data on the same form datum. Construct the actual base rows, host/frame lengths, nonzero chains and group ties, MultEq support, eigensystem/carrier/parts, all-embedding logLam, HD/d3b and NEb. Keep the stronger second DivRows3 scope distinct from the completed first DivDef3 and row-pair integrability. The full arithmetic entry and conditional common-frame construction do not by themselves install HOST\_\allowbreak NUMBER\_\allowbreak S, HOST\_\allowbreak FRAME\_\allowbreak LENGTH, TASK\_\allowbreak HOST\_\allowbreak VAN\_\allowbreak E\_\allowbreak INSTALL, TASK\_\allowbreak HOST\_\allowbreak XPP\_\allowbreak INSTALL or TASK\_\allowbreak COUNT\_\allowbreak V\_\allowbreak INSTALL.



\paragraph{23. BASE\_\allowbreak THETA\_\allowbreak RATIONAL\_\allowbreak TESTS — C — Simultaneous original base rational tests}

On the SAME original Bundle/reading/theta map, for every i<m supply one finite separating complex-linear test family whose values lie in the fixed embedded coefficient field on theta of EVERY original rational source vector and on EVERY stored target-coordinate basis vector. BTH8/BTH11 use this same LiftTestB for rational adjoint image and matrix coordinates. RTG now reconstructs prescribed-field coefficient extension, original split class representatives/RepsOK, the self-dual split-lattice square law and finite residue PhiIs; every Pins field follows from the actual CoversAt. With the base coefficient record literally fixed, exact class-sum existence/uniqueness including phase-rescaled O, PointIdentCore with common lambda,z before every O, and base ClassSumsSpan imply all-tuple finite lift span and OneMapOverL. The phase is exactly cj(twAt embLL A3.A.v0)(R.phs i). RatG uses the actual good projector eOp(y), and full LiftTestB follows at BOTH original levels. These finite implication proofs are complete. Remaining original producers are the prescribed higher covers, those actual class sums/uniqueness, PointIdentCore, base ClassSumsSpan and the base/common geometric datum; exact MLat is supplied only from original LatFull, whose geometric production remains separate. Finite dimension or an unrelated rational theta model cannot replace these same-object identities. Original outer eventual-family/Bundle quantifiers remain. Version0.46 proves finite generation of the original whole integral cusp lattice by actual integral-frame evaluation. It also proves the exact identity-cover, analytic unfolding, original-field coefficient descent, global clearing-scalar and both-level lattice-transport implications. Full spanning retains the supplied original IsNab connection and the five explicit principal-chart, level-system, raw-Gaussian, short-rule/dictionary and ramified-jet producers. The higher tuple covers, actual common class sums, PointIdentCore and ClassSumsSpan remain separate. CURRENT VERSION0.47 QUALIFICATION: Construct one original common geometric family and the higher CoversAt/all-tuple class sums and uniqueness with original phase rescalings, one PointIdentCore with common lambda,z, and base ClassSumsSpan. CNC now supplies the full original IsNab; PCL supplies the conditional trace-dual principal-lattice and quantitative normalization. Actual uniformization/HasPrinc and the five principal-chart, finite-level-system, Gaussian/reproduction, short-rule/simultaneous-dictionary and ramified-jet producer groups remain. Keep all original both-level, coefficient-field, Fits and eta agreements. Completed new original producers: ICL1-ICL21, CNC1-CNC16, PCL1-PCL12. Any earlier narrative naming these completed producers as outstanding is historical; this qualification is the current scope authority. CURRENT VERSION0.48 QUALIFICATION: Construct the original finite holomorphic FamilyUnif/N1c, normalized theta families and local lifting, then the common principal charts and finite-level theta systems. PUN now supplies the exact fixed-space models, all-family center lattices, A=1 exponent matching and individual primitive points with the full point-equality criterion; CNC supplies the full canonical connection and PCL the quantitative trace-dual normalization. Complete higher CoversAt and all-tuple class sums/uniqueness with the original phase rescalings, one PointIdentCore with common lambda,z, and base ClassSumsSpan. The Gaussian/reproduction, short-rule/simultaneous-dictionary and ramified-jet producers, and all original both-level, coefficient-field, Fits and eta agreements remain explicit. Completed current producers or conditional consumers: ICL1-ICL21, CNC1-CNC16, PCL1-PCL12, PUN1-PUN24. Earlier frontier descriptions are historical; this qualification is the current scope authority. Original whole-contract premises and quantifiers are retained. CURRENT VERSION0.49 QUALIFICATION: Construct the adjoining simultaneous principal framed points and test-ring dictionary. Full finite holomorphic FamilyUnif/N1c, normalized theta families and local lifts are now supplied by HUF. The next bounded target is allTest\_\allowbreak of\_\allowbreak secBaseChange: represent FrPt by the Lie-eigenframe bundles and theta-section frame line, use the full n=1 SecBaseChangeSt on every nonflat coefficient base change, and prove finite-generated subalgebra descent with the exact TRing/FromTest transport. Then supply the full FramesOK/CoverTri/FibTri and IsUnf/hasPrinc package, including adapted mod-eight symplectic basis, normalized theta multiplier/weight and arithmetic invariance. Continue higher CoversAt, all-tuple class sums and uniqueness with original phase rescalings, one PointIdentCore with common lambda,z, base ClassSumsSpan, Gaussian/reproduction, short-rule dictionaries and ramified jets, retaining both-level, coefficient-field, Fits and eta agreements. Completed current producers or conditional consumers: ICL1-ICL21, CNC1-CNC16, PCL1-PCL12, PUN1-PUN24, HUF1-HUF21 (full original FamilyUnif/N1c and theta family). Earlier frontier descriptions are historical; this qualification is the current scope authority. Original whole-contract premises and quantifiers are retained.

\emph{Completed inputs (version 0.57).} SGV full original section growth, DJW actual integral germs and full every-h1 block\_\allowbreak of\_\allowbreak normal\_\allowbreak form, PBC1–PBC27 (full arbitrary-base polarized pullback, actual very-ampleness and section finiteness), SHV1–SHV21 (both original Serre conclusions, actual finite section model and full all-cover return), ASL1–ASL21 (actual section localization and full original arbitrary-base Perf producer), FCD1-FCD18 (full homogeneous-chart descent, graded presentation and Noetherian cube model), GSL1-GSL26, NFL1-NFL22, TDR1-TDR13, NFR1-NFR13, HUF1-HUF21, PUN1-PUN24, ICL1-ICL21, CNC1-CNC16, PCL1-PCL12, Inherited both-level RatG/LiftTestB consumers

\emph{Remaining producer scope.} PBC supplies full arbitrary-base polarized pullback, actual very-ampleness and every-integer-power SecM finiteness. SGV now supplies the full original secM\_\allowbreak finite\_\allowbreak and\_\allowbreak growth, hence GSL\_\allowbreak FINITE\_\allowbreak GROWTH: on the literal identity-base product there are fixed integers N>=1 and d>=1 with dim SecM(N*t)=d*t\textasciicircum{}g for every t>=1. Its explicit classical inputs are Milne Proposition 9.1 and Theorem 13.3(b),(c); the section-coordinate and scalar dictionaries, ampleness and geometric-integrality/dimension hypotheses are proved on the original object in every characteristic. DJW supplies full ThmP, arbitrary-index BiValue, BiContent/BiWindow, all four MtoHM fields and full block\_\allowbreak of\_\allowbreak normal\_\allowbreak form for every supplied restriction identity. It constructs actual integral germs with the exact unit determinant factor, and obtains the original row window from the original MConcl. The original RowData.LemmaM existence producer remains: adapted Cartier chains, residue ranks, pairing-kernel estimate and weighted remainder must construct that normal form on the given row data. Remaining for the geometric section laws are multiplication-pullback generated comparison, finite/flat/surjective and etale multiplication geometry, invariant-function descent on every relevant open, the later every-positive-power GSL\_\allowbreak VANISH\_\allowbreak CLOSED conclusion, the required arbitrary-test integrality, principal rank one, torsion-indexed section decomposition, and square/translation generators. SHV17 all-cover eventual Serre vanishing and SHV21 homogeneous-polynomial generation remain complete. The completed identity-base integrality and growth do not themselves supply these different contracts. Original Moduli and HUF/PUN foundations remain. Later finite-level systems, Gaussian/reproduction, simultaneous dictionaries, ramified jets, all-tuple class sums, phase-sensitive uniqueness, common PointIdentCore and base ClassSumsSpan remain, retaining both levels, whole fields, Fits and eta agreement.



\paragraph{24. SLR\_\allowbreak BASE\_\allowbreak ROWS\_\allowbreak PRODUCER — C — Actual base row identities on that same bundle}

For the SAME selected Bundle, common original fields and readings used by the other three BTH12 inputs, supply actual RowsHold0\_\allowbreak 3 uniformly eventually for every family. Its original primitive row realization, spectrum/type and basis identifications must be retained. BTH supplies a theta consumer only; it does not construct the base row identity. No separate witness may be chosen merely to satisfy this row requirement.



\clearpage
\section{Connected local components and propagation in rank two}
\label{sec:arithmetic-connectedness}

\paragraph{The conclusion and the rings on which it holds.}
This chapter proves the original rank-two component theorem.  A finite
weight ideal which satisfies the original curve-propagation condition
and has the original arithmetic seed lies on every type-I component.
The proof constructs the local component chains, bounds the bad loci,
finds the required characteristic-prime curves, and carries the seed
to the connected closed fibre.  Its automorphic hypotheses remain the
original hypotheses; they are not consequences of connectedness.

Fix a complete coefficient DVR \(O\), with finite residue field
\(k\) of odd characteristic \(\ell\), and a frame
\(F.\mathrm{Ok}(C_0)\).  For either permitted rank \(n=2,3\),
let \(S\) be the supplied universal space.  Write
\[
 A=S.A,\quad B=S.B,\quad \mathfrak m=\mathfrak m_B,
 \quad u=S.u,\quad d=F.d=\sum_v d_v.
                                                        \tag{ACC.1}
\]
The ring \(A\) is the actual algebra generated over \(B\) by the
matrix entries.  RTC7 gives its finite generation, Noetherianity and
direct \(\mathbb Z\)-grading, with \(A_0=B\), finite \(B\)-modules
\(A_j\), and entry weight \(w_i-w_j\).  In rank two the weights
are \(1,-1\), and in rank three they are \(1,0,-1\).
Every local character lies in \(B\), and every balanced entry monomial
lies in \(B\).  These are the upstream results from the supplied
space; no \(\mathrm{FinA}\) assumption is used.

For every original local object \(L:S.\mathrm{Loc}(x)\), set
\[
 \begin{gathered}
 O'=L.O',\quad R=L.R,\quad \iota:A\longrightarrow R,
 \quad b=\iota|_B,\quad \rho=\iota_*u,\\
 \pi_L(P)=b^{-1}(P),\qquad \operatorname{im}_L(P)=\iota^{-1}(P).
 \end{gathered}
                                                        \tag{ACC.2}
\]
ALC4 and ALC11 identify this very \(R\) with the completion of the
selected localization of \(O'\otimes_O A\), and give flatness and
\(x\cap B=\mathfrak m\).  The finite unramified \(O'/O\) and
its residue map are those of \(L\).  ALC17 gives continuity, the
residual representation, and the exact contraction of all named
entry ideals.  In particular, if the lower-entry ideal
\(\mathfrak t_+=u.\mathrm{triP}\) is not contained in \(x\),
its extension to \(R\) is the unit ideal.  Write
\(\mathfrak t_-=u.\mathrm{triM}\) for the upper-entry ideal.
In rank two type I means avoiding both ideals.

We keep the distinction between the universal complete ring \(S.R\),
its finite-type entry algebra \(A\), and the local completion \(R\).
Transcendence degrees below concern \(A/\eta\), never the fraction
field of an arbitrary completed chart quotient.  The argument allows
equal-characteristic coefficient DVRs and includes \(\ell=3\).

\paragraph{Classical inputs.}
We use the height theorem, the dimension formula for finite-type
domains over a universally catenary domain, and Grothendieck
connectedness in the following precise form.  If \(T\) is complete
Noetherian local with irreducible spectrum of dimension \(r\), and
\(J\) is proper and generated up to radical by \(h\) elements, then
\(V(J)\setminus Z\) is connected whenever
\(\dim Z<r-1-h\).  This follows from the cohomological-dimension
form in Stacks Project, Lemma 52.14.1, tag
\href{https://stacks.math.columbia.edu/tag/0ECR}{0ECR}:
an \(h\)-generator \v Cech complex has no local cohomology above
degree \(h\), and an irreducible \(r\)-dimensional spectrum has
connectedness dimension \(r\).  Complete Noetherian local rings are
universally catenary by the Cohen structure theorem; see
\href{https://stacks.math.columbia.edu/tag/032C}{032C}.
The dimension formula is
\href{https://stacks.math.columbia.edu/tag/02JU}{02JU}.
These are stated classical inputs, not unproved deformation-specific
predicates.  All their applications, including the numerical thresholds
and changes of rings, are proved below.  The simplified rank-two route
was suggested by
\href{https://github.com/JoseBrox/GRH_distillation/blob/fd38313e5c7018c3c988ba20510e6b0b3bca0519/paper/shared/selmer-support-continuation-rank-two.tex}
{Brox's rank-two supplement}, at its pinned November 2026 revision.
The needed arguments are given here with their literal source scopes
and their connections to our completed chapters; no status declaration
in that project is used as a proof.

\subsection{The complete local connectedness theorem}

% BEGIN ACC1
\paragraph{ACC1. Cutting dimensions without a regular-sequence assumption.}
\label{claim:ACC1}
Let \(T\) be Noetherian local and let \(f_1,\ldots,f_h\)
belong to its maximal ideal.  Then
\[
 \dim T\le \dim T/(f_1,\ldots,f_h)+h.                    \tag{ACC.3}
\]
If \(T\) is complete and its nilradical is prime, every minimal
prime \(p\) over these elements satisfies
\[
                        \dim T/p\ge \dim T-h.            \tag{ACC.4}
\]
The same assertion holds when their ideal is replaced by any ideal
with the same radical.  Nilpotents in \(T\) are permitted.

\begin{proof}
A Noetherian local ring has finite dimension: its finitely generated
maximal ideal and the height theorem give a finite bound.  Choose a
system of parameters of \(T/(f_1,\ldots,f_h)\).  Lifting its
\(s=\dim T/(f_1,\ldots,f_h)\) elements gives an ideal generated by
\(s+h\) elements whose radical is the maximal ideal of \(T\).
The height theorem gives \(\dim T\le s+h\), proving (ACC.3).

For (ACC.4), pass to \(D=T/\sqrt{0}\).  This is a complete
Noetherian local domain with the same spectrum and dimension as
\(T\).  It is universally catenary.  A minimal prime \(\bar p\)
over the images of the \(h\) elements has height at most \(h\).
In a catenary local domain, concatenating saturated chains from
\(0\) to \(\bar p\) and from \(\bar p\) to the maximal ideal
gives
\(\dim D=\operatorname{ht}(\bar p)+\dim D/\bar p\).
This proves (ACC.4).  Passing to the reduction and replacing an
ideal by its radical preserve all the primes and quotient dimensions
in this argument.  No independence of the equations was assumed.
\end{proof}
% END ACC1

% BEGIN ACC2
\paragraph{ACC2. Full original complete-local component chain.}
\label{claim:ACC2}
Let \(T\) be complete Noetherian local with prime nilradical.
Let \(J=(f_1,\ldots,f_h)\), with all \(f_i\) in its maximal
ideal, and suppose
\[
                          \dim T\ge D+1+h.                \tag{ACC.5}
\]
Every minimal prime \(p\) over \(J\) has \(\dim T/p\ge D+1\).
Any two such primes can be connected by a finite sequence of minimal
primes over \(J\), with successive members satisfying
\[
                          \dim T/(p+q)\ge D.              \tag{ACC.6}
\]
This is the full original \(\mathrm{c1ch}\), including its general
prime-nilradical scope.  The original additional antecedent that
\(T\) is an \(\mathrm{IsCO}\)-algebra over some complete DVR is
retained when applying the source theorem, but is not needed for
this stronger statement.

\begin{proof}
ACC1 proves the component dimension assertion.  The ideal \(J\)
is proper, so it has a nonempty finite set of minimal primes.
Form the graph on this set with edges (ACC.6).  If the graph were
disconnected, partition its vertices into two nonempty classes and
let \(Y_1,Y_2\) be the unions of the corresponding irreducible
components of \(V(J)\).  Their intersection is the finite union
of the \(V(p+q)\) for pairs in different classes.  Consequently
\(Z=Y_1\cap Y_2\) has dimension less than \(D\).
Neither \(Y_i\) lies in \(Z\): its component generic points
belong to only their own minimal components.  Thus
\(V(J)\setminus Z\) is separated into the two nonempty disjoint
closed, and hence open, subsets \(Y_i\setminus Z\).

Put \(r=\dim T\).  By (ACC.5), \(r-1-h\ge D\).
The stated connectedness theorem therefore forbids this separation.
The graph is connected, so any two vertices admit a finite edge
path, with the constant path allowed.  The proof uses the whole
irreducible spectrum of \(T\), not regularity or reducedness of
\(T\).  It also covers \(h=0\) and \(D=0\); in these cases
the same dimension assertions and graph argument remain valid.
\end{proof}
% END ACC2

% BEGIN ACC3
\paragraph{ACC3. Transporting the exact component graph through a quotient.}
\label{claim:ACC3}
If \(p:T\twoheadrightarrow R\) is a surjective ring map and
the minimal primes of \(\ker p\) have dimensions at least
\(D+1\) and are connected by (ACC.6), then the minimal primes of
\(R\) have those same bounds and chains.  In the original notation,
\[
             \mathrm{Chained}(\ker p,D)
                 \Longrightarrow \mathrm{Chained}(0_R,D).
                                                               \tag{ACC.7}
\]

\begin{proof}
The inverse-image correspondence for quotient rings identifies the
prime ideals of \(R\) with the primes of \(T\) containing the
kernel.  It preserves and reflects inclusion, hence identifies the
minimal primes in question.  For every ideal \(Q\subset R\),
\(R/Q\simeq T/p^{-1}(Q)\).  Moreover
\(p^{-1}(Q+Q')=p^{-1}(Q)+p^{-1}(Q')\), since both inverse images
contain \(\ker p\).  These isomorphisms preserve dimensions of
the components and of their pairwise intersections.  Transport each
member of a finite path by the quotient correspondence.  A constant
path remains a constant path.  This proves every clause of (ACC.7).
\end{proof}
% END ACC3

% BEGIN ACC4
\paragraph{ACC4. The APD presentation gives actual chains in both ranks.}
\label{claim:ACC4}
For every original local object in (ACC.2), at either rank \(2\)
or \(3\),
\[
                  \mathrm{Chained}(0_R,nd+1).              \tag{ACC.8}
\]
In particular every minimal component has dimension at least
\(nd+2\), and the minimal-component graph with intersections of
dimension at least \(nd+1\) is connected.  This needs neither
\(F.\mathrm{Big}\) nor a nontriangularity assumption on \(x\).

\begin{proof}
APD25 gives the actual surjective \(O'\)-algebra map
\[
 O'[[X_1,\ldots,X_N]]\twoheadrightarrow R,
 \qquad N=N'+nd+1,
 \qquad \ker p=(f_1,\ldots,f_{N'}).                       \tag{ACC.9}
\]
The source is a complete Noetherian local domain of dimension
\(N+1\).  The kernel is proper, so each listed relation lies in
its maximal ideal, including any zero relations used for padding.
Apply ACC2 with \(h=N'\) and \(D=nd+1\).  Its required dimension
is exactly
\(D+1+h=N'+nd+2=N+1\).  ACC3 then proves (ACC.8).
The coefficient ring \(O'\), point \(x\), map \(\iota\) and
all ordinary characters are unchanged.  The assertion does not
provide rank three's separate even/odd \(\mathrm{EPres}\) or
its \(\mathrm{Sharp}\) record.
\end{proof}
% END ACC4

\subsection{Dimension comparisons before and after completion}

% BEGIN ACC5
\paragraph{ACC5. Full dimension comparison for the original local chart.}
\label{claim:ACC5}
For both original ranks define
\[
 t(\eta)=\operatorname{trdeg}_{\operatorname{Frac}(B/(\eta\cap B))}
                              \operatorname{Frac}(A/\eta)
\]
for a prime \(\eta\subset A\).  For every prime \(P\subset R\),
\[
 \dim R/P\le \dim B/\pi_L(P)+t(\operatorname{im}_L(P)).   \tag{ACC.10}
\]
Conversely, for every prime \(\eta\subset x\) there is a prime
\(P\subset R\) with
\[
 \operatorname{im}_L(P)=\eta,
 \qquad \dim R/P=\dim B/(\eta\cap B)+t(\eta).             \tag{ACC.11}
\]
These assertions give both clauses of the full original
\(\mathrm{loc\_\allowbreak dim}\), with any supplied numerical upper bound
or exact value for the transcendence degree.

\begin{proof}
Put \(E=O'\otimes_O A\), and let \(y\) be the prescribed
maximal ideal at which it is completed.  The map \(O\to O'\)
is injective: a nonzero prime kernel would be the maximal ideal of
the DVR \(O\), contradicting
\(\mathfrak m_O O'=\mathfrak m_{O'}\ne0\).  Thus \(O'\) is
torsion-free, hence flat, over \(O\).  It is finite by the original
local-object data.  Consequently \(E\) is finite and flat over
\(A\).  ALC4 identifies \(R=\widehat{E_y}\) with the actual
maps, and \(y\cap A=x\).

Let \(q\subset y\) be a prime of \(E\), with contraction
\(\eta\) in \(A\) and \(\mathfrak p\) in \(B\).
The extension \(A/\eta\subset E/q\) is integral, so their
fraction fields have the same transcendence degree over
\(\operatorname{Frac}(B/\mathfrak p)\).  The residue field at
\(y\) is finite over \(B/\mathfrak m\): it is a field of
finite type over that field and algebraic by the maximal-ideal
theorem.  The complete ring \(B\), and its domain quotient
\(B/\mathfrak p\), are universally catenary.  The finite-type
dimension formula therefore gives
\[
               \dim (E/q)_y=\dim B/\mathfrak p+t(\eta).
                                                               \tag{ACC.12}
\]

For a prime \(P\subset R\), take \(q=P\cap E\).
The ring \(R/qR\) is the completion of \((E/q)_y\).
Completion preserves the dimension of a Noetherian local ring:
a system of parameters remains primary to the maximal ideal after
completion, giving the upper bound by its number of generators,
and faithful flatness lifts a prime chain, giving the reverse bound.
Hence (ACC.12) is its dimension.  The quotient map to
\(R/P\) proves (ACC.10).

For the converse, flat going down from \(y\) gives a prime
\(q\subset y\) of \(E\) over the prescribed \(\eta\subset x\).
Choose a minimal prime \(P\) over \(qR\) whose quotient has the
largest dimension.  There are finitely many such minimal primes;
their maximum dimension equals \(\dim R/qR\).  Flat going down
for \(E_y\to R\) shows that every minimal prime over \(qR\)
contracts to \(qE_y\): a strictly larger contraction would, by
flat going down to \(qE_y\), yield a smaller prime over \(qR\),
contradicting minimality.  Thus this chosen \(P\)
contracts to \(\eta\) in \(A\), and (ACC.12) proves (ACC.11).
No algebraicity statement about \(\operatorname{Frac}(R/P)\) is
used.  Quotients and completion retain their possible nilpotents
until the explicitly chosen prime quotients are taken.
\end{proof}
% END ACC5

% BEGIN ACC6
\paragraph{ACC6. Opposite weights leave at most one parameter.}
\label{claim:ACC6}
For either original rank, every off-triangular prime \(\eta\) of
\(A\) satisfies \(t(\eta)\le1\).  This is the full original
\(\mathrm{fibre}\) theorem.  Hence every off-triangular prime
\(P\subset R\) satisfies
\[
                         \dim R/P\le\dim B/\pi_L(P)+1.
                                                               \tag{ACC.13}
\]

\begin{proof}
Work in the domain \(D=A/\eta\), and let
\(K=\operatorname{Frac}(B/(\eta\cap B))\).  Some positive-weight
entry \(u\) and negative-weight entry \(v\) survive; write their
weights as \(a>0\) and \(-b<0\).  The balanced product
\(u^b v^a\) belongs to \(K^\times\).  Thus \(u\) is algebraic
over \(K(v)\).  For an entry \(z\) of positive weight \(c\),
the product \(z^b v^c\) belongs to \(K\), so \(z\) is also
algebraic over \(K(v)\).  For an entry of negative weight
\(-c\), use \(z^a u^c\in K\).  Weight-zero entries and all
ordinary characters already lie in \(K\).  These statements allow
zero entries; their algebraic equations still hold.

Since the entries generate \(A\) over \(B\), its fraction field
modulo \(\eta\) is algebraic over the field generated by the one
element \(v\).  This proves \(t(\eta)\le1\), without requiring
\(\eta\) homogeneous.  ALC17 identifies off-triangularity of
\(P\) with that of \(\operatorname{im}_L(P)\).  Apply ACC5
to get (ACC.13).  The argument uses integer weights and never
divides by \(\ell\), by a weight, or by three.
\end{proof}
% END ACC6

\subsection{The rank-two triangular locus}

% BEGIN ACC7
\paragraph{ACC7. A polynomial description of the strict triangular fibre.}
\label{claim:ACC7}
In rank two let \(W=F.\mathrm{strS}(k)\).  Thus \(a\in W\)
is locally constant and satisfies
\[
 \begin{split}
 a(gh)&=a(g)\Psi(h)+\Psi(g)^{-1}a(h),\qquad a(g^*)=a(g),\\
 a(\operatorname{loc}_{0,v}s)
   &=\gamma_v\bigl(\Psi(\operatorname{loc}_v s)
                         -\Psi(\operatorname{loc}_v s)^{-1}\bigr)
                    \quad\text{for some }\gamma_v\in k.
 \end{split}                                                \tag{ACC.14}
\]
Then \(r=\dim_k W\le C_0\), and the actual fibre algebra is
\[
                   A/(\mathfrak m A+\mathfrak t_-)
                       \simeq k[T_1,\ldots,T_r].           \tag{ACC.15}
\]
The isomorphism keeps every diagonal and local character equal to
its prescribed residual value.

\begin{proof}
Set \(\chi=\Psi^{-2}\delta\) on \(\Gamma\), where
\(\delta\) is the sign character.  For \(a\in W\), the function
\(z(g)=a(g)/\Psi(g)\) on \(G_0\), extended with \(z(c)=0\),
is a continuous \(\chi\)-cocycle.  Indeed the first identity in
(ACC.14) is its cocycle identity on \(G_0\), and the star identity
gives \(z(cgc^{-1})=-z(g)\), the required rule for the semidirect
product.  The last identity gives the local coboundaries, with
primitive \(-\gamma_v\).  The local separation condition makes
each primitive unique.  If its strict cohomology class is zero,
then \(z(s)=(\chi(s)-1)t\); evaluation at \(c\) gives
\(-2t=0\), so \(t=0\) and \(a=0\).  This embeds \(W\)
linearly in the original strict cohomology group bounded by
\(F.\mathrm{Ok}.b1p\), proving \(r\le C_0\).

Choose a basis \(a_1,\ldots,a_r\) of \(W\).  RTC13 proves,
for every \(k\)-vector space \(M\), that all strict cocycles
with values in \(M\) are finite sums of scalar-valued strict
cocycles.  Its coefficient proof applies to this fixed basis and
gives uniquely \(a(g)=\sum_j a_j(g)m_j\).  This includes
infinite-dimensional \(M\).  The coordinate evaluations on
\(W\) span \(W^*\): a vector annihilated by every evaluation
is the zero function.  Thus finitely many evaluations express
each \(m_j\) as a \(k\)-linear combination of values of \(a\).

Apply RTC15 to \(\sum_j T_j a_j\) over the polynomial algebra
\(k[T_1,\ldots,T_r]\).  It gives a map from the left side of
(ACC.15), sending the lower entry to this universal function.
In the opposite direction, the lower entry of the quotient's own
representation is a strict cocycle: RTC10 and RTC14 apply because
\(\mathfrak m\) and the upper entries have been killed.  Write
its coefficients as \(m_j\) and send \(T_j\mapsto m_j\).
Both composites agree on every lower entry.  They agree on \(B\),
the diagonals and local characters, which are the prescribed
constants, and these together generate the quotient.  Hence the
composites are identity maps.  Equivalently, RTC16's uniqueness
checks the same assertion.  This proves (ACC.15) on the full
possibly nonreduced fibre algebra, without taking its reduction.
\end{proof}
% END ACC7

% BEGIN ACC8
\paragraph{ACC8. The actual closed coefficient fibre has dimension at most \(C_0\).}
\label{claim:ACC8}
Let \(n=2\), let \(L\) be any original local object, and assume
\(\mathfrak t_+\not\subset x\).  Then
\[
                         \dim R/\mathfrak mR\le C_0.
                                                               \tag{ACC.16}
\]
This is a strengthened rank-two version of the original fibre
dimension bound.

\begin{proof}
Some lower entry \(v\) is outside \(x\), hence \(\iota(v)\)
is a unit of \(R\).  For every upper entry \(u'\), the product
\(u'v\) is balanced and therefore belongs to \(B\).  It belongs
to \(\mathfrak m\): in the universal residual representation
all off-diagonal entries are zero.  Dividing its image by the unit
\(\iota(v)\) gives \(\iota(u')\in\mathfrak mR\).
Thus \(\rho.\mathrm{triM}\subset\mathfrak mR\), as an actual
ideal inclusion, not just after reduction.

For a prime \(P\supset\mathfrak mR\), its contraction
\(\eta\) to \(A\) contains \(\mathfrak mA+\mathfrak t_-\).
By ACC7, \(A/\eta\) is a domain quotient of a polynomial ring
in at most \(C_0\) variables over \(k\).  Its transcendence
degree is at most \(C_0\), while \(\eta\cap B=\mathfrak m\).
ACC5 gives \(\dim R/P\le C_0\).  The dimension of a Noetherian
quotient is the maximum of the dimensions of its finitely many
minimal prime quotients.  Apply this to \(R/\mathfrak mR\)
to obtain (ACC.16).  The specified residue extension in \(L\)
has not been replaced; ACC5 accounts for it and the completion.
\end{proof}
% END ACC8

% BEGIN ACC9
\paragraph{ACC9. The full finite-character detector supplied by the frame.}
\label{claim:ACC9}
For any field \(k\) of characteristic \(\ell\) and any frame
with \(F.\mathrm{Ok}(C_0)\), there are
\[
 s_1,\ldots,s_{1+C_0+d}\in G_0,\qquad
 s'_1,\ldots,s'_d\in G_0,\qquad H>0                       \tag{ACC.17}
\]
with \(\Psi(s_a)=1\), and with each \(s'_a\) the image of
an original inertia generator on which \(\Psi=1\), such that
every open-kernel character \(\xi:G_0\to C\) to a finite
abelian \(\ell\)-group satisfies both implications:
\[
 \begin{split}
 \xi(s_a)=1\ \text{for all }a
       &\Longrightarrow \xi(g)^H=1\ \text{for all }g,\\
 \xi(g^*)=\xi(g),\quad \xi(s'_a)=1\ \text{for all }a
       &\Longrightarrow \xi(g)^H=1\ \text{for all }g.
 \end{split}                                                \tag{ACC.18}
\]
This is the full original \(\mathrm{ab\_\allowbreak gen}\); no chosen
coefficient DVR is required.  The finiteness of \(k\) is already
one of the \(F.\mathrm{Ok}\) fields.

\begin{proof}
The space of continuous additive homomorphisms \(\Gamma\to k\)
embeds in the trivial-character \(H^1_C\) with full local
subspaces: send a homomorphism to its cocycle with zero local
primitives.  A boundary has zero global component.  The frame's
\(h1\) bound therefore makes this space have dimension at most
\(1+C_0\).  Choose \(1+C_0\) evaluations detecting it, padding
with the identity if necessary.  Call the corresponding group
elements \(t_j\).  They detect every finite abelian
\(\ell\)-group character of \(\Gamma\): a nontrivial image has
a nontrivial quotient of order \(\ell\), and composing with
\(\mathbb F_\ell\hookrightarrow k\) would give an undetected
nonzero additive character.

Choose the original \(d_v\) inertia generators at each place
from \(F.\mathrm{Ok}.gen\), and enumerate their global images
as \(s'_1,\ldots,s'_d\).  If \(\xi\) kills these generators,
the same generator property kills its restriction to every original
inertia group.  If also \(\xi(g^*)=\xi(g)\), then
\(\xi(cgc^{-1})=\xi(g)^{-1}\); it kills the conjugate inertia
groups as well.  The exact \(cl\) field supplies one nonzero
integer \(H\) killing every such unramified character.  This
proves the second implication of (ACC.18).

Write \(t_j=a_jc^{\epsilon_j}\), with \(a_j\in G_0\), and
put \(q=|k^\times|\).  As the first \(1+C_0\) elements of
\(s\), take \((a_jca_jc^{-1})^q\), and append the \(s'_a\).
All have \(\Psi=1\).  If \(\xi\) kills them, coprimality
of \(q\) and \(\ell\) shows that
\(\xi(a_j)\xi(ca_jc^{-1})=1\).  The character
\(g\mapsto\xi(g)\xi(cgc^{-1})\) is conjugation invariant,
so extends to \(\Gamma\) by sending \(c\) to one.  Its
kernel is open, since \(G_0\) and the original kernels are open.
It kills all \(t_j\), so is trivial by their detecting property.
Thus \(\xi(g^*)=\xi(g)\).  The appended inertia generators
and the second implication give the first.  The same \(H\)
works in both assertions and is chosen before every test character.
\end{proof}
% END ACC9

% BEGIN ACC10
\paragraph{ACC10. Passing finite-character tests to an adic character.}
\label{claim:ACC10}
Let \(D\) be Noetherian local of characteristic \(\ell\),
let \(G\) be a compact topological group, and let
\(\psi:G\to D^\times\) be adically continuous.  Suppose
\(q\) is prime to \(\ell\), \(\psi(g)^q\equiv1\pmod{
\mathfrak m_D}\) for all \(g\), and specified elements \(u_a\)
are sent to one.  If the original finite abelian \(\ell\)-group
generator test at these \(u_a\) forces every such character to
kill an element \(g\), and \(\psi(g)\equiv1\pmod{
\mathfrak m_D}\), then \(\psi(g)=1\).

Also, if \(D\) is a domain and a residually trivial adic character
\(\xi:G\to D^\times\) has exponent dividing one fixed \(H>0\)
in every finite adic image, then \(\xi=1\).

\begin{proof}
For each \(N\ge1\), the image of \(\psi^q\) in
\((D/\mathfrak m_D^N)^\times\) is finite: its map is locally
constant and its domain compact.  The image lies in \(1+
\mathfrak m_D/\mathfrak m_D^N\), an \(\ell\)-group on any
finite subgroup, since
\((1+z)^{\ell^a}=1+z^{\ell^a}\) and \(\ell^a\ge N\)
annihilates the second term.  The kernel is open.  Apply the
stated finite generator property to this actual finite image.
It gives \(\psi(g)^q=1\) modulo \(\mathfrak m_D^N\).
The residually trivial \(\psi(g)\) also has \(\ell\)-power
order there, so Bezout's identity gives \(\psi(g)=1\) in that
quotient.  Krull intersection gives equality in \(D\).
Neither completeness nor finiteness of the residue field was
needed in this first argument.

For the second assertion the same finite-image argument and
separatedness give \(\xi(g)^H=1\) in \(D\).  Write
\(H=\ell^a H_0\), with \(H_0\) prime to \(\ell\).
Since \(D\) is a domain,
\((\xi(g)^{H_0}-1)^{\ell^a}=0\) implies
\(\xi(g)^{H_0}=1\).  The factor
\(1+\xi(g)+\cdots+\xi(g)^{H_0-1}\) has residue \(H_0\)
and is a unit.  Its product with \(\xi(g)-1\) is zero, so
\(\xi(g)=1\).  This proves the asserted passage from bounded
finite exponents to the exact character identity.
\end{proof}
% END ACC10

% BEGIN ACC11
\paragraph{ACC11. The triangular cuts force the full coefficient ideal.}
\label{claim:ACC11}
Return to rank two and a local object with
\(\mathfrak t_+\not\subset x\).  Choose \(s_a\) from ACC9
and a uniformizer \(\varpi\) of \(O\).  Put
\[
 J=\rho.\mathrm{triM}+
       (\varpi,\ \rho(s_a)_{00}-1\ (1\le a\le1+C_0+d))
                          \subset R.                     \tag{ACC.19}
\]
Then \(\mathfrak mR\subset\sqrt J\).  All listed cuts are
in \(\mathfrak m_R\), and their number, beyond the triangular
ideal, is \(d+C_0+2\).

\begin{proof}
The coefficient map is local, and the diagonal reduction at each
\(s_a\) equals \(\Psi(s_a)=1\), so the cuts belong to
\(\mathfrak m_R\).  Let \(P\) be any prime containing \(J\)
and put \(D=R/P\).  It is a separated complete local domain of
characteristic \(\ell\), with its specified embedding of \(k\).
The triangular representation has diagonal characters \(\alpha,
\beta\).  Dividing \(\alpha\) by the embedded \(\Psi\)
gives a residually trivial adic character which kills all \(s_a\).
ACC9 applied to each finite adic image, followed by ACC10, gives
\(\alpha=\Psi\).  Polarization gives
\(\beta(g)=\alpha(g^*)=\Psi(g)^{-1}\).

The local characters also equal their prescribed values, including
their order.  To see this without choosing a new flag, use ACP4's
actual characteristic-polynomial identity.  At a local separating
element \(h\), its two diagonal values have distinct residues.
The two roots supplied by the ordered local characters are therefore
matched to these diagonal values by their residues.  For arbitrary
\(g\), compare the traces at \(g\) and \(hg\).  The resulting
two linear equations have determinant
\(\alpha(h)-\beta(h)\), a unit, and force the same ordered
matching at \(g\).  Thus both full local characters are residual.

A balanced entry monomial containing a strictly lower entry also
contains an upper entry, so its image is zero.  The other balanced
monomials, and every local character value, are the prescribed
constants.  Hence the dense generating subalgebra \(B_0\) of
\(B\) maps to its residue in \(D\).  RTC7's topology comparison
makes \(B_0\) dense for the actual maximal-ideal topology on
\(B\), and the local map \(B\to D\) is continuous.  For
\(b\in B\), its approximants have eventually the fixed residue
of \(b\); separatedness then gives the same image for \(b\).
Consequently \(\mathfrak m\) maps to zero in \(D\).
This holds for every prime \(P\supset J\), proving the radical
inclusion.  The argument has retained all character values, not
only their semisimplified characteristic polynomials.
\end{proof}
% END ACC11

% BEGIN ACC12
\paragraph{ACC12. The full rank-two triangular dimension bound.}
\label{claim:ACC12}
Under the rank-two hypotheses of ACC11,
\[
 \dim R/\rho.\mathrm{triM}\le d+2C_0+2,
 \qquad \rho.\mathrm{triP}=R.                            \tag{ACC.20}
\]
In particular the original rank-two instance of \(\mathrm{l51}\)
holds, whose numerical upper bound is \(d+2+(3C_0+5)\).
If \(F.\mathrm{Big}(C_0)\), then
\(\dim R/\rho.\mathrm{triM}<2d\).

\begin{proof}
ACC11 gives \(\dim R/J\le\dim R/\mathfrak mR\), and ACC8
bounds the latter by \(C_0\).  Apply ACC1 to the local quotient
\(R/\rho.\mathrm{triM}\) and its \(d+C_0+2\) additional
cuts.  This proves the first bound in (ACC.20).  The unit lower
entry proves the second, as in ALC17.  The displayed bound is no
larger than \(d+3C_0+7\), so gives exactly the required source
instance.  Finally \(\mathrm{Big}\) says \(d>3C_0+7\);
in particular \(d>2C_0+2\), which proves the strict inequality.
There is no assertion here that the rank-three instance of
\(\mathrm{l51}\) has been supplied.
\end{proof}
% END ACC12

% BEGIN ACC13
\paragraph{ACC13. Full original rank-two sharpness.}
\label{claim:ACC13}
For every original rank-two space and local object, if
\(F.\mathrm{Big}(C_0)\) and \(\mathfrak t_+\not\subset x\),
then
\[
 \mathrm{Chained}(0_R,2d+1),\qquad
       \rho.\mathrm{TI}(P)\quad\text{for every minimal prime }P
       \text{ of }R.                                     \tag{ACC.21}
\]
This is the full original \(\mathrm{sharp2}\).

\begin{proof}
The chain assertion is ACC4.  Let \(P\) be a minimal prime.
Its quotient has dimension at least \(2d+2\).  The lower-entry
ideal is the unit ideal, so it is not contained in \(P\).
If the upper-entry ideal lay in \(P\), the quotient map would
give
\(2d+2\le\dim R/P\le\dim R/\rho.\mathrm{triM}<2d\),
contradiction.  Thus \(P\) avoids both triangular ideals.
In rank two this is precisely type I; the extra exceptional-type
clause of the rank-three definition has false antecedent.  Every
original quantifier over \(O,F,S,x,L\) has been retained.
\end{proof}
% END ACC13

\subsection{Coincident characters and characteristic-prime curves}

% BEGIN ACC14
\paragraph{ACC14. Counting the exact ordinary-character cuts.}
\label{claim:ACC14}
This claim applies at both ranks \(n=2,3\).  Let \(I\subset B\)
satisfy the original finite-weight condition
\[
 \mathfrak m\subset\sqrt{I+\lambda},\qquad
 \lambda=\mathfrak m_O B+
  (\psi_{v,i}(g)-1:g\in I_v,\ \Psi(\operatorname{loc}_v g)=1).
                                                               \tag{ACC.22}
\]
Let \(\mathfrak p\) be any ideal of \(B\), containing \(I\),
such that for specified \(v\) and \(i\ne j\) the two ordered
characters agree on \(N_v\) modulo \(\mathfrak p\).  Then
\[
                        \dim B/\mathfrak p\le nd-d_v+2.
                                                               \tag{ACC.23}
\]
No primality or characteristic assumption on \(\mathfrak p\)
is required.  This supplies the original coefficient-ring
dimension argument \(\mathrm{dim\_\allowbreak B}\).

\begin{proof}
Use the original inertia generators \(u_{w,a}\), with one exceptional
index \(a_0(w)\) at each place; all nonexceptional generators lie
in \(N_w\).  Impose these equations in \(B/\mathfrak p\):
the image of \(\varpi\), every
\(\psi_{w,r}(u_{w,a})-1\) except the \(d_v\) equations for
\((w,r)=(v,j)\), and the one omitted exceptional equation
\(\psi_{v,j}(u_{v,a_0(v)})-1\).  There are exactly
\(1+(nd-d_v)+1=nd-d_v+2\) listed elements, all in the maximal
ideal.  A zero or duplicate element does not affect this bound.

In any prime quotient of these cuts, every generator equation holds.
Indeed the missing nonexceptional \(j\)-equations follow by
subtracting the \(i\)-equation and using coincidence on \(N_v\).
The quotient is a local domain of characteristic \(\ell\).
For every original local character, its \(q=|k^\times|\)-th
power is residually one.  Apply ACC10 with the exact local generator
property of \(F.\mathrm{Ok}.gen\).  It makes the character one
on every inertia element on which its prescribed residual value is
one.  Hence this prime contains \(\lambda\), as well as \(I\).
Condition (ACC.22) makes it the maximal ideal of \(B\).

The cut quotient therefore has dimension zero.  ACC1 gives (ACC.23).
If \(\mathfrak p=B\), its quotient is zero and the assertion is
automatic with the convention \(\dim0=-\infty\).  Counting the
uniformizer remains necessary in equal characteristic: the equation
\(\ell=0\) there does not impose \(\varpi=0\).
\end{proof}
% END ACC14

% BEGIN ACC15
\paragraph{ACC15. Full original coincident-prime dimension bound in both ranks.}
\label{claim:ACC15}
For either original rank, let \(P\subset R\) be prime and
off triangular.  Suppose \(I\subset\pi_L(P)\), where \(I\)
satisfies (ACC.22), and suppose its two specified distinct ordinary
characters coincide on \(N_v\) modulo \(P\).  Then
\[
                            \dim R/P+d_v\le nd+3.         \tag{ACC.24}
\]
In particular this proves both conclusions of the full original
\(\mathrm{l52}\): the bound \(nd+6\), and the bound \(nd+5\)
when \(\mathfrak m_O R\subset P\).  No extra characteristic
condition is needed for the stronger bound (ACC.24).

\begin{proof}
The full ordered character differences contract to
\(\mathfrak p=\pi_L(P)\), since the map \(b\) is the
restriction of the original \(\iota\).  Apply ACC14 to get
\(\dim B/\mathfrak p\le nd-d_v+2\).  ACC6 and ACC5 give
\(\dim R/P\le\dim B/\mathfrak p+1\).  Combining these
inequalities proves (ACC.24).  It implies the two stated weaker
source bounds, retaining the antecedent of the second rather than
changing it.  There has been no identification of the rank-two and
rank-three rings and no use of rank-two triangularity in this proof.
\end{proof}
% END ACC15

% BEGIN ACC16
\paragraph{ACC16. Every bad rank-two locus has the required strict dimension.}
\label{claim:ACC16}
Let \(n=2\), assume \(F.\mathrm{Big}(C_0)\), and suppose
\(\mathfrak t_+\not\subset x\).  Besides the triangular
bound \(\dim R/\rho.\mathrm{triM}<2d\), for every \(I\)
satisfying (ACC.22), every \(v\), and every \(i\ne j\),
\[
 \dim R/\bigl(IR+\mathfrak m_O R+C_{v,i,j}\bigr)<2d,
 \quad
 C_{v,i,j}=(\rho_{v,i}(g)-\rho_{v,j}(g):g\in N_v).
                                                               \tag{ACC.25}
\]
This is the original strict-dimension theorem at rank two.

\begin{proof}
The first assertion is ACC12.  Let \(Q\) be any prime containing
the ideal in (ACC.25).  If \(\rho.\mathrm{triM}\subset Q\),
its quotient dimension is bounded by the triangular quotient.
Otherwise it is off triangular, because the lower-entry ideal is
the unit ideal.  All hypotheses of ACC15 hold for \(Q\), so
\(\dim R/Q\le2d+3-d_v<2d\), using \(d_v>5\).
There are finitely many minimal primes over the displayed ideal;
their dimensions are integers, and the dimension of the quotient is
their maximum.  This proves the same strict inequality for the whole
quotient, including its nonreduced structure.  If that quotient is
zero the convention \(-\infty<2d\) gives the assertion directly.
\end{proof}
% END ACC16

% BEGIN ACC17
\paragraph{ACC17. Full original prime-avoidance curve lemma.}
\label{claim:ACC17}
Let \(T\) be any Noetherian local ring, let \(a\in\mathfrak m_T\),
let \(N\ge1\), and let \((J_i)\) be a finite family of ideals
with \(\dim T/J_i<N\).  For every ideal \(K\) with
\(\dim T/K\ge N+1\), there is a prime \(p\) such that
\[
 K\subset p,\quad a\in p,\quad \dim T/p=1,
                     \quad J_i\not\subset p\quad\text{for all }i.
                                                               \tag{ACC.26}
\]
This is the full original \(\mathrm{NicePrimeEx}\), without
completeness, catenarity, infinitude of a residue field, or primality
of \(K\) added.

\begin{proof}
By ACC1, \(\dim T/(K+(a))\ge N\).  Select a minimal prime
\(Q\) of \(K+(a)\) whose quotient has dimension at least \(N\).
No \(J_i\) is contained in \(Q\), since that would contradict
its dimension bound.  Choose \(b_i\in J_i\setminus Q\).
Also choose \(g\in\mathfrak m_T\setminus Q\), possible because
\(\dim T/Q\ge1\).  The product \(f=g\prod_i b_i\) belongs
to \(\mathfrak m_T\setminus Q\).  If the family is empty, use
the product one.

By the ascending chain condition choose a maximal prime \(p\)
among primes containing \(Q\) and avoiding \(f\).  In the local
domain \(D=T/p\), the image of \(f\) is nonzero and nonunit,
and every nonzero prime contains it.  We prove \(\dim D=1\).
If \(\dim D\ge2\), all minimal primes over \((f)\) have
height one, by the principal ideal theorem.  There are finitely
many, and none is the maximal ideal.  Prime avoidance supplies
\(z\) in the maximal ideal outside their union.  A minimal prime
\(q\) over \((z)\) is nonzero and has height one.  By the
defining property of \(D\), it contains \(f\), hence contains
a minimal prime over \((f)\).  Two such height-one primes in
a domain cannot be strictly nested, so \(q\) is that prime,
contradicting the choice of \(z\).  Thus \(\dim D\le1\).
The nonzero nonunit \(f\) shows its maximal ideal is nonzero,
so \(\dim D\ge1\).  This proves equality.

The chosen \(p\) contains \(K+(a)\).  If it contained any
\(J_i\), it would contain \(b_i\) and hence \(f\), a
contradiction.  This proves all clauses of (ACC.26).  Only finite
prime avoidance was used, so finite residue fields cause no issue.
\end{proof}
% END ACC17

% BEGIN ACC18
\paragraph{ACC18. Actual Nice curves above every sufficiently large rank-two ideal.}
\label{claim:ACC18}
At rank two assume the hypotheses of ACC16, and let \(I\subset B\)
satisfy (ACC.22).  For every ideal \(K\subset R\), not necessarily
prime, with
\[
                       IR\subset K,\qquad \dim R/K\ge2d+1,
\]
there is an original Nice prime \(p\supset K\).  Explicitly,
\[
 \dim R/p=1,\quad \ell\in p,\quad
 \rho.\mathrm{OffTri}(p),\quad
 \neg\rho.\mathrm{Coin}(p,v,i,j)\quad(v,\ i\ne j).
                                                               \tag{ACC.27}
\]
This is the full rank-two instance of
\(\mathrm{exists\_\allowbreak nice\_\allowbreak prime\_\allowbreak above\_\allowbreak of\_\allowbreak large\_\allowbreak dim}\).

\begin{proof}
Apply ACC17 with \(N=2d\), \(a\) the image of the actual
uniformizer \(\varpi\) of \(O\), and the following finite
family of bad ideals: \(\rho.\mathrm{triM}\) and all
\(IR+\mathfrak m_O R+C_{v,i,j}\) for \(i\ne j\).
Their strict dimension bounds are ACC16.  The frame's \(\mathrm{Big}\)
condition ensures \(N\ge1\), and the coefficient map sends
\(\varpi\) into \(\mathfrak m_R\).

The resulting prime contains \(\varpi\), hence contains
\(\mathfrak m_O R\) and \(\ell\).  It avoids the upper-entry
ideal; it also avoids the lower-entry ideal, which is the unit
ideal.  If its two local characters coincided on \(N_v\), it
would contain \(C_{v,i,j}\); together with \(IR\subset K\subset p\)
and \(\mathfrak m_OR\subset p\), this would contradict avoidance
of the corresponding bad ideal.  This proves (ACC.27) with the
literal ordered local tests.  In rank two every such prime is
type I, and ALC15 identifies its field-valued representation with
an absolutely irreducible one.  The prime itself, not a change to
a more favorable coefficient chart, is the required output.
\end{proof}
% END ACC18

% BEGIN ACC19
\paragraph{ACC19. Actual local rank-two propagation data.}
\label{claim:ACC19}
Fix the same rank-two space, an ideal \(I\) satisfying (ACC.22),
and the original condition \(P1'(S,I)\): for every original
\(x,L\), every Nice type-I prime \(p\) with
\(I\subset\pi_L(p)\), and every prime \(P\subset p\),
one has \(I\subset\pi_L(P)\).  At any chart with
\(\mathfrak t_+\not\subset x\), the literal rank-two
\(\mathrm{LocData}\) with threshold \(2d+1\), sharp ideal
zero, and empty exceptional and thick sets satisfies all ten
original \(\mathrm{Ok}\) clauses.

\begin{proof}
Good primes are those containing \(IR\), hence are upward closed.
Type-I primes avoid the two triangular ideals, hence are downward
closed.  Every Nice prime is type I.  ACC18 supplies the required
Nice prime above each good prime of dimension at least \(2d+1\).
The thick-prime and exceptional-chain clauses are vacuous.  A type-I
minimal prime is a minimal prime over zero; ACC13 supplies their
adjacency chains and shows that every minimal prime is type I.
The remaining downward-goodness clause is exactly the supplied
\(P1'\), on this chart.  These verify the ten fields individually.

For clarity, the resulting propagation also has the following
direct description.  A good prime of dimension at least \(2d+1\)
lies in a Nice prime by ACC18.  A minimal prime below it is good
by \(P1'\).  If a good minimal prime \(P\) is adjacent to
\(Q\), the ideal \(P+Q\) has dimension at least \(2d+1\)
and contains \(IR\).  ACC18 gives a Nice prime above it, and
\(P1'\) makes \(Q\) good.  ACC13's finite chain reaches every
minimal prime.  This is the actual local use of ACP14, with no
new local propagation predicate left to assume.
\end{proof}
% END ACC19

\subsection{Closed-fibre incidence and the original seed}

% BEGIN ACC20
\paragraph{ACC20. Minimal primes of the actual entry algebra are homogeneous.}
\label{claim:ACC20}
At both original ranks, every minimal prime \(\eta\) of \(A\)
is homogeneous for the actual grading.  This is the full original
\(\mathrm{min\_\allowbreak hom}\).

\begin{proof}
For any prime \(p\) in a directly \(\mathbb Z\)-graded ring,
let \(p^h\) consist of the elements all of whose homogeneous
components belong to \(p\).  It is a homogeneous ideal contained
in \(p\).  It is prime.  Indeed, for elements \(a,b\notin p^h\),
choose the largest indices \(r,s\) for which their components
\(a_r,b_s\) do not belong to \(p\).  Such largest indices exist
because each element has only finitely many homogeneous components.
In degree \(r+s\), the product has the term \(a_r b_s\),
which is outside \(p\).  Every other term lies in \(p\), since
its first index exceeds \(r\) or its second exceeds \(s\).
Thus the degree \(r+s\) component of \(ab\) is outside \(p\),
proving \(ab\notin p^h\).

For a minimal prime \(\eta\) of \(A\), the prime
\(\eta^h\subset\eta\) must equal \(\eta\).  Consequently
\(\eta\) is generated by its homogeneous elements, precisely
the original \(S.\mathrm{IsHom}(\eta)\) condition.  No action
of the finite residue field's multiplicative group is used to
detect degrees; the direct integer grading is essential here.
\end{proof}
% END ACC20

% BEGIN ACC21
\paragraph{ACC21. Closed points in every nonempty part of the core.}
\label{claim:ACC21}
At both original ranks put
\[
 \mathcal C=S.\mathrm{BB}
   =V(\mathfrak m A)\cap D(\mathfrak t_+),\qquad
 \mathcal Z=\{x\in\mathcal C:x\text{ is maximal in }A\}.
                                                               \tag{ACC.28}
\]
If a closed subset \(Y\subset\operatorname{Spec}A\) and an
open subset \(U\) satisfy \(Y\cap U\cap\mathcal C\ne\varnothing\),
their intersection contains a point of \(\mathcal Z\).
This is the full original \(\mathrm{jac}\).  In particular
\(\mathcal Z\) is dense in \(\mathcal C\), and every closed
set meeting \(\mathcal C\) meets \(\mathcal Z\).

\begin{proof}
The fibre \(A/\mathfrak m A\) is a finite-type algebra over
\(k\).  At a given point of the stated intersection, choose a
basic open \(D(f)\subset U\) and a lower entry \(v\) which
does not vanish.  The ring obtained by imposing the closed equations
of \(Y\) on the fibre and inverting \(fv\) is nonzero and
of finite type over \(k\).  It has a maximal ideal, and its
residue field is finite over \(k\) by the maximal-ideal theorem.
Contract that ideal to the fibre ring.  Its quotient embeds in this
finite field extension; being a \(k\)-subalgebra of a finite
dimensional field, it is itself a field.  Thus the contracted ideal
is maximal, avoids \(fv\), and contains all equations of \(Y\).
Its inverse image in \(A\) is the required maximal point.
Taking \(Y\) to be the whole spectrum proves density; taking
\(U\) to be the whole spectrum proves the final assertion.
The residue field may be a finite extension of \(k\); there is
no assertion of a \(k\)-rational point.
\end{proof}
% END ACC21

% BEGIN ACC22
\paragraph{ACC22. Every homogeneous component with a lower entry meets the core.}
\label{claim:ACC22}
At both ranks, if \(\eta\subset A\) is a homogeneous prime
and \(\mathfrak t_+\not\subset\eta\), there exists a maximal
point \(x\in\mathcal C\) with \(\eta\subset x\).
This is the full original \(\mathrm{l41}\).

\begin{proof}
Suppose, to the contrary, that
\(\mathfrak t_+\subset\sqrt{\eta+\mathfrak m A}\).
Noetherianity gives \(M\ge1\) such that
\(\mathfrak t_+^M\subset\eta+\mathfrak m A\).
Every entry has weight between \(-2\) and \(2\), and the
negative-weight entries are precisely lower entries.  A monomial
of weight at most \(-2M\) must therefore contain at least
\(M\) lower factors.  RTC7's monomial spanning description
implies
\[
                   A_j\subset\mathfrak t_+^M
                           \quad(j\le-2M).               \tag{ACC.29}
\]

Choose a lower entry \(v\notin\eta\), of weight \(e\le-1\).
Set \(j=2Me\), so \(j\le-2M\).  By (ACC.29), the finite
\(B\)-module \(Q_j=A_j/(A_j\cap\eta)\) is equal to
\(\mathfrak mQ_j\).  To see equality with this exact module,
project the inclusion \(A_j\subset\eta+\mathfrak m A\)
to degree \(j\): \(\eta\) is homogeneous, and
\(\mathfrak m\) has degree zero.  Nakayama gives \(Q_j=0\).
But \(v^{2M}\in A_j\), so \(v^{2M}\in\eta\), contradicting
primality and \(v\notin\eta\).

Thus the asserted radical containment fails.  There is a prime
containing \(\eta+\mathfrak m A\) and avoiding
\(\mathfrak t_+\).  Apply ACC21 with \(Y=V(\eta)\) to
obtain a maximal such point \(x\).  This proves the required
incidence by finite generation and the actual graded pieces, without
assuming that the whole entry algebra is finite over \(B\).
\end{proof}
% END ACC22

% BEGIN ACC23
\paragraph{ACC23. Preserving the weaker seed inclusion in rank two.}
\label{claim:ACC23}
Let \(n=2\).  Suppose \(\xi\subset A\) is a type-I prime,
\(Z\subset B\) is prime, and \(Z\subset\xi\cap B\).
There is a homogeneous type-I prime \(\eta\subset A\) with
\[
                         \eta\cap B=Z,\qquad t(\eta)=1.
                                                               \tag{ACC.30}
\]
This is the rank-two instance of the original homogeneous seed-prime
supplier.  It does not strengthen \(Z\subset\xi\cap B\) to
an equality as an input.

\begin{proof}
Choose an upper entry \(u\) and a lower entry \(v\) outside
\(\xi\).  Their balanced product \(a=uv\) belongs to \(B\)
and is outside \(\xi\cap B\), hence outside \(Z\).
In \(A_a\), both \(u\) and \(v\) are units.  Every upper
entry \(u'\) and lower entry \(v'\) satisfies
\[
                   u'=\frac{u'v}{a}u,\qquad
                   v'=(uv')u^{-1}.                        \tag{ACC.31}
\]
The coefficients in these expressions belong to \(B_a\), while
diagonal entries and all ordinary characters already belong to
\(B\).  Consequently the map
\(B_a[U,U^{-1}]\to A_a\), \(U\mapsto u\), is surjective.
It is injective as well: a finite Laurent relation is a sum of
distinct homogeneous degrees, so each coefficient times the
corresponding unit power of \(u\) vanishes separately.  Each
coefficient is therefore zero.  This proves the actual graded
isomorphism
\[
                         A_a\simeq B_a[U,U^{-1}].         \tag{ACC.32}
\]

Extend \(ZB_a\) to the Laurent ring and contract its prime
ideal to \(A\), calling the result \(\eta\).  This prime
is homogeneous, contracts to \(Z\), and avoids both \(u\)
and \(v\).  It is therefore type I.  Localizing its quotient
at \(a\) gives \((B/Z)_a[U,U^{-1}]\); hence the fraction-field
extension has transcendence degree exactly one.  All choices were
made on the same \(A,B\).  No coefficient extension or local
object has yet been chosen.
\end{proof}
% END ACC23

% BEGIN ACC24
\paragraph{ACC24. Full rank-two seed incidence with the original order of choices.}
\label{claim:ACC24}
Under the rank-two assumptions of ACC23, there exists a maximal
\(x\in\mathcal C\) such that, for every supplied original
\(L:S.\mathrm{Loc}(x)\), there exists a prime \(P\subset R\)
satisfying
\[
 \rho.\mathrm{TI}(P),\qquad \pi_L(P)=Z,\qquad
                      \dim R/P=\dim B/Z+1.               \tag{ACC.33}
\]
This is the complete rank-two instance of \(\mathrm{seedpt}\).
The point \(x\) is chosen before the universally quantified
local object and its coefficient extension.

\begin{proof}
Take the homogeneous prime \(\eta\) of ACC23.  ACC22 gives a
maximal \(x\in\mathcal C\) above it; fix this \(x\) once.
For any later \(L\) at this point, the second clause of ACC5,
with \(t(\eta)=1\), supplies a prime \(P\subset R\)
contracting to \(\eta\), with the dimension in (ACC.33).
Its coefficient contraction is exactly \(Z\) by (ACC.30).
ALC17's exact ideal comparisons make it type I because \(\eta\)
is type I.  The same reasoning applies to every allowed \(L\),
with its own specified unramified extension, without rechoosing
\(x\) or \(\eta\).  In the original seed hypothesis
\(\dim B/Z=1+2d\), this produces a good type-I prime of
dimension \(2d+2\), exceeding the required threshold \(2d+1\).
\end{proof}
% END ACC24

\subsection{Return to the actual component theorem}

% BEGIN ACC25
\paragraph{ACC25. All three actual local-to-global chart conclusions.}
\label{claim:ACC25}
At rank two retain \(F.\mathrm{Big}\), \(I\)'s original
finite-weight condition and \(P1'(S,I)\).  Define
\[
 T=D(\mathfrak t_+)\cap D(\mathfrak t_-),\qquad
 G=V(IA)\subset X=\operatorname{Spec}A.                   \tag{ACC.34}
\]
For every \(x\in\mathcal Z\), goodness propagates between
all type-I irreducible components through \(x\), and at least
one type-I component passes through \(x\).  If an original
chart at \(x\) has a good type-I prime of dimension at least
\(2d+1\), some type-I component of \(X\) is wholly good.
These are the three literal conclusions of the original
rank-two \(\mathrm{at\_\allowbreak pt}\) application.

\begin{proof}
ALC21 supplies an original local object at the specified maximal
point \(x\) over \(\mathfrak m\).  Its map \(A\to R\)
is flat by ALC11 and has closed-point contraction \(x\).
ACP16 therefore gives the component chart: minimal primes of
\(R\) contract to minimal primes of \(A\), and every global
component through \(x\) is obtained in this way.  ALC17 identifies
the inverse image of \(T\) with the actual local type-I locus,
and ideal contraction identifies the inverse image of \(G\)
with the local good locus.  The first set is open and the second
closed by (ACC.34).

ACC19 supplies the full local propagation conclusion at this very
chart.  ACC13 gives dimension at least \(2d+2\), hence at least
\(2d+1\), and type I for every local minimal prime.  All
hypotheses of ACP17 are now supplied.  Its first, second and third
parts give exactly the three asserted conclusions.  More concretely,
a good global component through \(x\) gives a good local minimal
prime; ACC19 makes every local minimal prime good, and contraction
gives every other global component through \(x\).  A local
minimal prime always exists and is type I, while any good local
prime above the numerical threshold starts the same argument.
This uses the canonical local map, not an abstract substitute chart.
\end{proof}
% END ACC25

% BEGIN ACC26
\paragraph{ACC26. Full original propagation theorem in rank two.}
\label{claim:ACC26}
For every original coefficient DVR, frame \(F.\mathrm{Ok}(C_0)\)
with \(F.\mathrm{Big}(C_0)\), rank-two space \(S\), and
ideal \(I\subset B\), suppose
\[
 \mathrm{FinLam}(I),\qquad P1'(S,I),\qquad
 \begin{gathered}
 \xi\subset A\text{ is a type-I prime},\quad Z\subset B\text{ is prime},\\
 I\subset Z\subset\xi\cap B,\qquad \dim B/Z=1+2d
 \end{gathered}                                             \tag{ACC.35}
\]
for some \(\xi,Z\).  Then
\[
             I\subset p\cap B
       \quad\text{for every type-I prime }p\subset A.     \tag{ACC.36}
\]
This is the full original \(\mathrm{PropagationTheoremGL2}\),
or \(\mathrm{components\_\allowbreak good\_\allowbreak rank\_\allowbreak two}\), at all its
original premises.

\begin{proof}
We verify the actual global propagation data on \(X\), with
type-I set \(T\) and good set \(G\) from (ACC.34), core
\(\mathcal C\), selected closed points \(\mathcal Z\), and
empty exceptional set.  Noetherianity of \(X\) is RTC7.
The inclusion \(\mathcal Z\subset\mathcal C\) is its definition.
RTC19 supplies preconnectedness of this exact core.  ACC21 shows
that every closed set meeting the core meets \(\mathcal Z\),
and that the core lies in the closure of \(\mathcal Z\).

Every relevant irreducible component meets the core.  Indeed let
its minimal prime be \(\eta\).  Relevance means that it contains
a type-I point, so \(\eta\) itself is type I by generization.
ACC20 makes \(\eta\) homogeneous, and ACC22 supplies a core
point above it.  ACC25 supplies propagation between relevant
components through a selected closed point, and shows that every
such closed point belongs to a relevant component.  These verify
the first seven original global fields individually.

For the final field use the supplied seed \(\xi,Z\), retaining
its weaker inclusion.  ACC24 chooses a closed core point before
any local chart and gives, for a chart there, a type-I prime \(P\)
with \(\pi_L(P)=Z\) and dimension \(2d+2\).  Since
\(I\subset Z\), this prime is good.  ACC25 therefore supplies
one relevant component wholly contained in \(G\).

All eight fields of the actual global data have now been proved.
ACP15 makes every relevant component good and hence gives
\(T\subset G\).  For a prime \(p\), membership in \(G\)
is exactly \(IA\subset p\), or \(I\subset p\cap B\).
This proves (ACC.36).  The proof has not constructed an automorphic
seed or \(P1'\); it has used precisely the seed and propagation
premises in the original theorem.  The same frame and marked
element are retained throughout.
\end{proof}
% END ACC26

% BEGIN ACC27
\paragraph{ACC27. The exact two-rank consumer and what remains upstream.}
\label{claim:ACC27}
In the original \(\mathrm{AbsMain}\) construction, once its
rank-two \(\mathrm{Setup}\) and \(\mathrm{ModIn}\) data have
supplied their actual \(\mathrm{FinLam}\), \(P1'\) and seed
transfers, ACC26 returns the required rank-two \(\mathrm{AllGood}\)
on that same space and ideal.  It can then be supplied to the
original rank-three component theorem.  The unconditional local
chains (ACC.8), full \(\mathrm{loc\_\allowbreak dim}\), \(\mathrm{fibre}\),
\(\mathrm{l52}\), \(\mathrm{min\_\allowbreak hom}\), \(\mathrm{jac}\)
and \(\mathrm{l41}\) are also available at rank three with
their stated original scopes.

\begin{proof}
The actual consumer uses separate maps \(\Phi_2,\Phi_3\) from
its determinant rings to the two weight-zero rings and separate
ideals \(J_2,J_3\).  Its rank-two inputs are the transferred
ideal \(\Phi_2(J_2)\), the corresponding \(\mathrm{FinLam}\),
\(P1'\), and \(\mathrm{SeedIn}\).  These match (ACC.35)
literally, so ACC26 returns
\(\mathrm{AllGood}(S_2,\Phi_2(J_2))\), the next argument
passed by \(\mathrm{AbsMain}\) to the rank-three theorem.
No equality between the two spaces, rings, coefficient extensions
or ideals is needed or asserted.  Each of the additional listed
suppliers was explicitly proved above for both ranks, so is
available on \(S_3\) independently of this rank-two return.

Rank three still requires its actual even/odd presentation and
\(\mathrm{Sharp}\) argument, its remaining triangular bound,
corner/seed incidence and the thick-chain installation.  The ordinary
APD presentation does not supply those additional parity equations.
The actual automorphic \(\mathrm{ModIn}\) and original
\(\mathrm{ModU}\), \(P\)-from-automorphy and seed producers
remain separate arithmetic work.  Once the rank-three \(\mathrm{OffFin}\)
producer is supplied, ACP11, RTC24 and AFE20 give the already
proved finite-point, coefficient-descent and finite-end returns;
UCE's downstream integral finiteness need not be reproved here.
The original \(P\) premise, Host scope \(\ell=3\), common
marking, and extra Alt witnesses \(\mathrm{WitA}\) and
\(\mathrm{WitAQ}\) are unchanged.  Thus this chapter completes
the exact rank-two component theorem and its listed subtheorems,
not a whole terminal arithmetic obligation.
\end{proof}
% END ACC27

\paragraph{Reading the new boundary.}
The substantive new rank-two construction is the component theorem
ACC26: its local chains, curve existence, homogeneous incidence and
seed localization have now been proved on the actual rings.  Only
the original finite-weight, propagation and seed assumptions remain
at its entrance.  The two-rank and terminal arithmetic joins retain
their own supplied structures.  No whole selected obligation, full
\(\mathrm{GdEnd}\), full \(\mathrm{HostEnd}\), or Alt endpoint
is claimed by the local and rank-two completion above.


\section{The period Gram identity on every supplied original setting}
\label{sec:period-gram-v57}

The purpose of this chapter is to identify the actual Poincar\'e periods
with the original sum of orbital operators.  The identity is proved for
every supplied original setting and every pair of words, under exactly
the analytic premises of \texttt{KExpansion\allowbreak C}.  No arithmetic gap
inequality is added.  This is possible because that statement already
supplies the folding, square-integrability, orbital-integrability and
convergence hypotheses.  We prove the intervening measure decompositions
and finite-character identities; we do not replace them by another
period-matching assumption.

Brox's \emph{actual Hilbert-valued period Gram expansion}, at revision
\texttt{fd38313e}, suggested the useful
organization into support cells, rational returns and an antilinear
column map.  Its installed arithmetic gap is replaced here by the
literal universal analytic premises.  In particular, a summable series
of integrated operators is never used as a substitute for the absolute
integrability needed to unfold before integration.  That integrability
is proved separately from the supplied folding data.

\subsection{The fixed objects and the exact endpoint}

Fix an original datum \(D\), an arbitrary supplied
\(X:\texttt{Hida\allowbreak Setting}(D)\), and an arbitrary
\(Y:\texttt{Full}(X)\).  Write \(E/F\) for its CM extension,
\(G\) for its rational unitary group, and
\[
 G_{\mathbb A}=G_\infty\times G_f,\qquad
 \Gamma=G(F),\qquad d\mu_{\mathbb A}=d\mu_\infty\,d\mu_f.
\]
These are the original groups and measures, transported through the
supplied coefficient coordinates.  The rational embedding is the
original \texttt{Rat\allowbreak Link}; we suppress that embedding in formulas.
Write \(r\) for the length of the isotropic frame and \(k\) for the
positive complementary rank, so the Hermitian dimension is \(2r+k+1\).
Let \(K_0\) be the stabilizer of the completed reference lattice and
let \(K=K_2\), with its given character \(\chi=\chi_2\), be the
original smaller level.  Set \(v_K=\mu_f(K)\).  The level modulus is
the original product \(g=g_0g_1\) of two distinct rational primes.

For a word \(u\), retain its negative vector \(z_u\), its
hyperbolic frame \(e_i,f_i\), its finite transport \(x_u\), its
compact-type space \(V_u\), and its wave function
\(W_u:G_\infty\to V_u\).  The original stabilizer is
\(S_u=N_u\rtimes H_u\).  Here \(H_u\) fixes every
\(e_i,f_i,z_u\), while \(N_u\) acts upper unitriangularly on the
\(e\)-flag and trivially on \(e^\perp/\langle e\rangle\).
The wave-function law, with inner products conjugate linear in their
first variable, is
\[
 W_u(sx)=\psi_{u,\infty}(s)\pi_u(s)W_u(x),
 \qquad \pi_u|_{N_u}=1.
\tag{PGR.1}
\]
The infinite representation is unitary on \(S_u\); \(H_{u,\infty}\)
is a product of compact unitary groups.

The finite set \(I_u\) and representatives \(h_{u,\rho}\) are
those of \(Y.cs(u)\).  They represent
\[
 H_u(F)\backslash H_u(\mathbb A_f)/K_{H,u},\qquad
 K_{H,u}=H_u(\mathbb A_f)\cap x_uKx_u^{-1}.
\]
Put
\[
 a_{u,\rho}=h_{u,\rho}x_u,\qquad
 \widehat\Lambda_{u,\rho}=a_{u,\rho}\widehat L_0,
 \qquad N_{u,\rho}=N_u(F)\cap\operatorname{Aut}(\Lambda_{u,\rho}).
\tag{PGR.2}
\]
The equality with the actual row lattice is a supplied field of
\(Y.cs(u)\).  All the original row conditions, splitting conditions,
finite reduction data and local models of \(Y\) are retained.
They are not newly constructed by this chapter.

The column space \(C_u\) is the original space of linear maps
\(T:V_u\to\{\hbox{functions on }H_{u,\infty}\times H_u(\mathbb A_f)\}\)
with rational left invariance, right \(K_{H,u}\)-character \(\chi\),
and compact-type covariance.  Its scalar seed is exactly
\[
 \Phi_u(T)(x,n h x_u\kappa)
   =\psi_{u,f}(n)\chi(\kappa)\,T(W_u(x))(1,h),
 \quad n\in N_u(\mathbb A_f),\ h\in H_u(\mathbb A_f),\ \kappa\in K;
\tag{PGR.3}
\]
it is zero outside this finite support.  Independence of the
decomposition in (PGR.3) is proved below.  The original Poincar\'e
series, for the arbitrary representatives \(D_u=Y.Dt(u)\), is
\[
 P_uT(x)=\sum_{\delta\in S_u(F)\backslash\Gamma}
                         \Phi_u(T)(\delta x).
\tag{PGR.4}
\]
The quotient notation in this formula means that the representatives
are multiplied on the left of \(x\), as in the source.

Keep the original strict domains \(F_{N,u}=Y.FN(u)\) and set
\[
 c_u=\frac{v_K\,\mu_{N,u}(F_{N,u})\,\mu_{H,u}(H_{u,\infty})}
                 {\dim_{\mathbb C}V_u},\qquad
 s_{u,u'}=\bigl(\sqrt{c_uc_{u'}}\bigr)^{-1}.
\tag{PGR.5}
\]
At the exact formal interface, each mass in this display means its
original real-valued mass, and inverse and division are total, with
\(0^{-1}=0\).  Thus (PGR.5) means the literal
\texttt{c0X}/\texttt{scl}, even before positivity is supplied.
No argument below divides by \(c_u\) to prove an unnormalized identity.
The positive finite number \(v_K\) will be established independently.

For two words, the precise retained analytic assumptions are these.
Both seeds have the original \texttt{Fold\allowbreak OK} data, expanded in
PGR12 below.  Both words satisfy \texttt{In\allowbreak HH}: their scalar series
converge absolutely at every point, are independent of the supplied
coset representatives, lie in the original quotient \(L^2\), and
belong to its specified automorphic range.  The original term set has
\texttt{Fund\allowbreak Exists}, \texttt{OOInt\allowbreak All}, \texttt{OOFund\allowbreak Indep}, and
\texttt{Amp\allowbreak Well\allowbreak Def}: its orbital domains exist, its prescribed
operator integrands are integrable, and the integrals and
phase-weighted amplitudes have the stated independence of choices.
Finally the original typed operator datum satisfies \texttt{Abs\allowbreak Conv\allowbreak T}.
All quantifiers in these premises are retained, including every column
vector and every permitted domain or representative.  The last premise
includes summability of the Hilbert--Schmidt norms of the original term
operators.  Nothing here presupposes a spectral estimate for their sum.

\subsection{The finite phase, the Levi factor and the columns}

% BEGIN PGR1
\paragraph{PGR1. The supplied additive character annihilates the inverse different.}
\label{claim:PGR1}
Let \(\mathfrak d=\mathfrak d_{F/\mathbb Q}\mathcal O_E\), the
original ideal \(\texttt{dd\allowbreak S}(D)\).  The supplied character
\(\psi_{E,f}\) is unitary and locally constant, and
\[
       \psi_{E,f}(t y)=1
       \quad(y\in\mathfrak d^{-1},\ t\in\widehat{\mathcal O}_E).
\tag{PGR.6}
\]
It therefore annihilates the entire completed fractional ideal
\(\widehat{\mathfrak d^{-1}}\).  This includes the full original
\texttt{Psi\allowbreak Dd} statement for every allowed supplied character.
\begin{proof}
The defining identities are multiplicative additivity, triviality on
\(\widehat{\mathcal O}_E\), and
\[
 \psi_{E,f}(b)\exp(2\pi i\operatorname{Tr}_{E/\mathbb Q}b)=1
                         \quad(b\in E).
\tag{PGR.7}
\]
The trace in (PGR.7) is exactly the sum over the chosen CM embeddings
and their conjugates in the original infinite character.  For an
arbitrary finite adele \(a\), some positive integer \(m\) makes
\(ma\) integral.  Additivity gives \(\psi_{E,f}(a)^m=1\).
Thus every value has modulus one.  Triviality on the open integral
subring makes the character locally constant; no continuity assumption
has been inserted.

If \(y\in\mathfrak d^{-1}\) and \(o\in\mathcal O_E\), then
\(\operatorname{Tr}_{E/F}(oy)\in\mathfrak d_{F/\mathbb Q}^{-1}\).
Indeed multiplication by each element of \(\mathfrak d_{F/\mathbb Q}\)
makes \(oy\) integral over \(\mathcal O_F\), and hence makes its
relative trace integral.  The defining trace-dual property of the
different now gives
\(\operatorname{Tr}_{E/\mathbb Q}(oy)\in\mathbb Z\).

Choose a nonzero integer \(m\) with \(my\in\mathcal O_E\).
Finite Chinese remaindering gives
\(t=o+mt_1\), with \(o\in\mathcal O_E\) and
\(t_1\in\widehat{\mathcal O}_E\).  Consequently
\(ty=oy+t_1(my)\).  The second term is integral, and the first
has integral absolute trace.  Equation (PGR.7) proves (PGR.6).
A finite set of integral-module generators for
\(\mathfrak d^{-1}\) generates its completed module as well;
apply (PGR.6) to each generator and multiply the character values.
This proves the last assertion, including at ramified finite places.
\end{proof}
% END PGR1

% BEGIN PGR2
\paragraph{PGR2. Rational splitting and the actual unipotent characters.}
\label{claim:PGR2}
Over \(F\), every completion of \(F\), and the finite adeles,
the actual group \(S_u\) is \(N_u\rtimes H_u\).  The splitting
is defined over \(F\).  The characters \(\psi_{u,f}\) and
\(\psi_{u,\infty}\) are unitary characters of \(S_u\), trivial
on \(H_u\), invariant under its conjugation on \(N_u\), and
\[
       \psi_{u,f}(n_f)\psi_{u,\infty}(n_\infty)=1
                           \quad(n\in N_u(F)).
\tag{PGR.8}
\]
In particular this proves the original rational Levi and rational
character-product assertions.
\begin{proof}
Let \(U=\langle e_0,\ldots,e_{r-1}\rangle\), and let \(W\)
be the complement orthogonal to all \(e_i,f_i\).  The hyperbolic
pairings give the rational projection
\[
 p_W(y)=y-\sum_i h(f_i,y)e_i-\sum_i h(e_i,y)f_i.
\tag{PGR.9}
\]
An element of \(S_u\) induces an isometry of \(U^\perp/U\),
identified with \(W\), fixing \(z_u\).  Extend this isometry
by the identity on every \(e_i,f_i\); the result is its unique
Levi factor in \(H_u\).  The residual factor is exactly in
\(N_u\).  This construction and its inverse use the displayed
projection and the unitary equations, so commute with all the stated
scalar extensions.  In particular a rational element has rational
factors.

In the original character put
\[
 B_u(s)=\sum_{i+1<r}h(f_i,se_{i+1})+
             \begin{cases}h(f_{r-1},sz_u),&r>0,\\0,&r=0.\end{cases}
\tag{PGR.10}
\]
The finite character is \(\psi_{E,f}(B_u(s)/g)\); the infinite
one is the product of the corresponding real exponential characters.
The upper unitriangular action on the \(e\)-flag gives
\(B_u(st)=B_u(s)+B_u(t)\): in an adjacent entry there is no
intermediate index, and in the last row of the \(z_u\)-column
all earlier-row contributions vanish.  A Levi factor fixes all the
vectors occurring in (PGR.10).  Unitarity then gives
\(B_u(hnh^{-1})=B_u(n)\) and \(B_u(h)=0\).
This proves the character and conjugacy statements.  Their values
are unitary by PGR1 and the explicit real exponentials.
For rational \(n\), (PGR.10) is a rational element of \(E\),
and (PGR.8) is precisely (PGR.7) applied to \(B_u(n)/g\).
All sums are empty, with value zero, when \(r=0\).
\end{proof}
% END PGR2

% BEGIN PGR3
\paragraph{PGR3. The exact finite-level phase, including the different.}
\label{claim:PGR3}
For every original \(X,u\) and every
\(n\in N_u(\mathbb A_f)\cap K_0\),
\[
 x_u^{-1}nx_u\in K,\qquad
             \chi(x_u^{-1}nx_u)=\psi_{u,f}(n).
\tag{PGR.11}
\]
This is the complete original \texttt{Unipotent\allowbreak Conjugate\allowbreak In\allowbreak Level}.
\begin{proof}
The membership assertion is the already proved RCV12, whose scope is
every supplied finite reference frame and transport.  It imposes no
character comparison.  We prove that remaining comparison here.
Let \(A\) be the actual numerator defining
\(\chi(x_u^{-1}nx_u)=\psi_{E,f}(A/g)\), including both its
\(e,z\) chain and its \(bw\) chain.  Let \(B=B_u(n)\).
All vectors in the numerator \(A\) are integral, and
\(x_u^{-1}nx_u\) preserves the integral self-dual lattice.
Thus \(A\in\widehat{\mathcal O}_E\).

We first locate \(B\) at every finite place.  The actual lattice
splits as
\[
 L_0=\bigoplus_{i<r}\mathfrak a_i e_i\ \oplus L_b\ \oplus
               \bigoplus_{i<r}\mathfrak a_i^{-1} f_i,
             \qquad \mathfrak a_i=\mathfrak d^{i-r}.
\tag{PGR.12}
\]
Projection onto its \(e_i\)-coordinate sends it into
\(\mathfrak a_i\).  Since \(n\) preserves the completed
lattice, its coefficient of \(e_i\) in \(ne_{i+1}\) belongs
to \(\mathfrak a_i\mathfrak a_{i+1}^{-1}=\mathfrak d^{-1}\).
Also \(z_u\in L_b\), so the coefficient of \(e_{r-1}\)
in \(nz_u\) belongs to \(\mathfrak a_{r-1}=\mathfrak d^{-1}\).
These are exactly the summands of \(B\).  Therefore
\(B\in\widehat{\mathfrak d^{-1}}\).

At each place above \(g\), the supplied reference vectors agree
modulo \(g\) with the rational vectors, after the actual transport:
\(x_ue_i^0\equiv e_i\), \(x_uf_i^0\equiv f_i\), and
\(x_uz^0\equiv z_u\).  These congruences have precisely the
localized meaning proved in ASG11; their clearing integers are units
at these places.  Applying the lattice isometries \(n,x_u\) and
pairing preserves the congruences.  The \(e,z\) part of \(A\)
is thus congruent to \(B\) modulo \(g\).  The strengthened
\(bw\)-congruence in the proof of RCV12 says that
\((x_u^{-1}nx_u-1)bw_j\) is in the span of the \(e_i^0\)
modulo \(g\).  Pairing with \(bwd_i\) annihilates this span
modulo \(g\), and the adjacent pairing
\(h(bwd_i,bw_{i+1})\) is zero modulo \(g\).  Hence the
entire additional \(bw\) numerator is zero modulo \(g\).
We obtain \((A-B)/g\in\widehat{\mathcal O}_{E,v}\) there.

Away from \(g\), division by \(g\) is an integral unit.
Since \(\mathcal O_E\subset\mathfrak d^{-1}\), the preceding
locations of \(A\) and \(B\) show that \((A-B)/g\) belongs
to the local completed inverse different there too.  Thus globally
\((A-B)/g\in\widehat{\mathfrak d^{-1}}\).  PGR1 gives
\(\psi_{E,f}((A-B)/g)=1\), proving (PGR.11).  This argument
explains why integrality of the rational dual vectors at every place
is neither assumed nor needed.
\end{proof}
% END PGR3

% BEGIN PGR4
\paragraph{PGR4. Integral Levi extraction and decomposition independence.}
\label{claim:PGR4}
If \(n\in N_u(\mathbb A_f)\), \(h\in H_u(\mathbb A_f)\),
and \(nh\in K_0\), then both \(h,n\in K_0\).
Moreover
\[
 S_u(\mathbb A_f)\cap x_uKx_u^{-1}
   =(N_u(\mathbb A_f)\cap K_0)\rtimes K_{H,u}.
\tag{PGR.13}
\]
The displayed factorization is unique.  These assertions supply the
original \texttt{Kset\allowbreak Split} and its smaller-level consequence.
\begin{proof}
The projection (PGR.9) preserves the completed lattice by (PGR.12):
it is zero on its two isotropic summands and the identity on \(L_b\).
For an integral vector \(y\), put \(p=p_W(y)\).  The Levi
factor preserves \(W\), and \(n\) is the identity on
\(U^\perp/U\), so
\[
                  p_W(nhp)=hp.
\]
The left side is integral because \(p\) is integral and \(nh\)
preserves the lattice.  The vector \(y-p\) lies in the hyperbolic
summand and is fixed by \(h\).  Hence
\(hy=(y-p)+hp\) is integral.  Apply the same argument to
\((nh)^{-1}=(h^{-1}n^{-1}h)h^{-1}\); PGR2 ensures that
its first factor lies in \(N_u\).  It follows that \(h^{-1}\)
also preserves the lattice.  Thus \(h\in K_0\), and then
\(n=(nh)h^{-1}\in K_0\).

If \(nh\in x_uKx_u^{-1}\), it is in \(K_0\) because
\(x_u\in K_0\).  PGR3 puts \(n\) in \(x_uKx_u^{-1}\).
Dividing gives \(h\in K_{H,u}\).  Conversely both factors on
the right of (PGR.13) belong to the left by PGR3.  Uniqueness is
the uniqueness of the splitting in PGR2.  Notice that the proof uses
the actual fractional-ideal splitting, not a globally integral
hyperbolic basis that the data do not provide.
\end{proof}
% END PGR4

% BEGIN PGR5
\paragraph{PGR5. Triviality of every original column stabilizer.}
\label{claim:PGR5}
For every \(h\in H_u(\mathbb A_f)\),
\[
          H_u(F)\cap hK_{H,u}h^{-1}=\{1\}.
\tag{PGR.14}
\]
This is the full original \texttt{Neat\allowbreak Ku}, for every word and
every permitted level-character witness.
\begin{proof}
Let \(\gamma\) belong to the intersection.  It preserves the
rational lattice whose completion is \(hx_u\widehat L_0\).
This is a full lattice: clearing the finitely many finite-place
denominators bounds it between \(mL_0\) and \(m^{-1}L_0\)
for some positive integer \(m\), and intersection with the rational
space recovers the stated completed lattice.  The element \(\gamma\)
fixes the hyperbolic pairs and \(z_u\).  Their further orthogonal
complement is positive definite at every real place.  On that
complement, the absolute trace of the Hermitian form is a positive
definite rational quadratic form.  Project the full lattice onto
this rational complement.  The projection is rational, so its image
is a full integral lattice after clearing one denominator.
The induced automorphism preserves that lattice and its positive
definite quadratic form.

Such a lattice has only finitely many isometries: the image of each
fixed integral basis vector lies in a bounded set of lattice points,
and there are finitely many possibilities for all basis images.
The restriction is faithful because the other summands are fixed.
Consequently \(\gamma\) has finite order.

The integral matrix
\(\kappa=x_u^{-1}h^{-1}\gamma_fhx_u\) belongs to \(K\).
The original residue-frame calculation in ASG14 makes its reduction
unipotent at both \(g_0\) and \(g_1\).  To recall why no rank
restriction is hidden here, the \(e\)-flag is upper unitriangular;
the \(z\)-column is fixed modulo that flag; and the \(bw\)-flag
is upper unitriangular after projection off those fixed summands.
The integral idempotent and its conjugate split the remaining
residue module.  Their dual pairings make the action on the conjugate
summand the inverse dual.  It too is unipotent.  The supplied spanning
condition accounts for the whole residue module.

If an integral finite-order matrix has unipotent reduction in residue
characteristic \(p\), its order is a power of \(p\).  Indeed write
its order as \(p^ab\), with \((p,b)=1\).  Its \(p^a\)-th
power \(M\) has order dividing \(b\); its reduction is unipotent
and has order dividing \(b\), hence is the identity.  In
\((M-I)(I+M+\cdots+M^{b-1})=0\), the second factor reduces
to the invertible matrix \(bI\), and is invertible over the local
ring.  Thus \(M=I\).  Applying this at the two distinct primes
forces the order of \(\gamma\) to be one.  This works at ramified
places as well and never compares a prime with the rank.
\end{proof}
% END PGR5

% BEGIN PGR6
\paragraph{PGR6. The original column coordinates are antilinear Hilbert coordinates.}
\label{claim:PGR6}
For \(T\in C_u\), let \(v_{u,\rho}(T)\in V_u\) be determined by
\[
       T(\xi)(1,h_{u,\rho})
                    =\langle v_{u,\rho}(T),\xi\rangle.
\tag{PGR.15}
\]
Then \(J_uT=(v_{u,\rho}(T))_\rho\) is exactly the original
\texttt{sharp}; it is an antilinear bijection onto
\(\mathcal V_u=\bigoplus_{\rho\in I_u}V_u\).  The original
\texttt{Col\allowbreak T}/\texttt{col\allowbreak Core} inner product is
\[
       \langle T,T'\rangle_{C_u}=\langle J_uT',J_uT\rangle_{\mathcal V_u}.
\tag{PGR.16}
\]
Thus \(J_u\) is an antilinear isometry; in particular the original
column space is finite dimensional with this prescribed norm.
\begin{proof}
Riesz representation on the finite-dimensional space \(V_u\)
gives (PGR.15), additivity and conjugate scalar multiplication.
Every \(h_f\) has a unique column class and a presentation
\(h_f=\gamma_fh_{u,\rho}\kappa\), with
\(\gamma\in H_u(F)\), \(\kappa\in K_{H,u}\).
For a fixed class, PGR5 makes \(\gamma\), and therefore
\(\kappa\), unique.  An arbitrary list of linear functionals
\(\lambda_\rho\in V_u^*\) consequently defines
\[
 T(\xi)(h_\infty,h_f)
   =\chi(x_u^{-1}\kappa x_u)\,
              \lambda_\rho\bigl(\pi_u(\gamma_\infty^{-1}h_\infty)\xi\bigr).
\tag{PGR.17}
\]
Uniqueness of the presentation proves this is well defined.  The
representation and character laws verify rational left invariance,
finite right covariance and compact-type covariance directly.
Conversely those three laws force (PGR.17) from the evaluations
(PGR.15).  Hence evaluation is a linear bijection with the list of
duals, and Riesz representation makes \(J_u\) an antilinear
bijection.  Equation (PGR.16) is the literal transported inner product,
not a new assertion about an earlier norm on a function space.
It is positive definite by injectivity, and its norm is
\(\|T\|^2=\sum_\rho\|v_{u,\rho}(T)\|^2\), as claimed.
\end{proof}
% END PGR6

% BEGIN PGR7
\paragraph{PGR7. The original scalar seed is well defined and has the required covariance.}
\label{claim:PGR7}
Formula (PGR.3) is independent of its decomposition.  It defines a
measurable function, invariant under \(S_u(F)\) on the left and
transforming by \(\chi\) under \(K\) on the right.  It also
satisfies the original convolution reproduction formula.  Thus both
\texttt{Ph\allowbreak Well\allowbreak Def} and \texttt{Ph\allowbreak Props} hold on every supplied \(X\).
\begin{proof}
Suppose \(nhx_u\kappa=n'h'x_u\kappa'\).  The element
\((n'h')^{-1}nh\) belongs to \(x_uKx_u^{-1}\).
Its Levi factor is \(h'^{-1}h\), and PGR4 says this lies in
\(K_{H,u}\); its unipotent factor lies in \(N_u\cap K_0\).
The character equality in PGR3, the conjugation invariance in PGR2,
and the right covariance of \(T\) then give equality of the two
values in (PGR.3).  More explicitly, write
\(n'^{-1}n=h' n_0 h'^{-1}\) and \(h=h'h_0\).  Then
\(n_0h_0=x_u\kappa'\kappa^{-1}x_u^{-1}\), and
\[
 \psi_{u,f}(n)=\psi_{u,f}(n')\psi_{u,f}(n_0),\qquad
 T(\xi)(1,h)=\chi(x_u^{-1}h_0x_u)T(\xi)(1,h').
\]
Their product with \(\chi(\kappa)\) is the expression using
\(n',h',\kappa'\), by (PGR.11).

Right \(K\)-covariance is now immediate.  For left rational
covariance split \(s=n_sh_s\in S_u(F)\) by PGR2.
At infinity (PGR.1) contributes
\(\psi_{u,\infty}(n_s)\pi_u(h_s)\).  At finite places the
unipotent factor contributes \(\psi_{u,f}(n_s)\), while
rational left invariance and compact-type covariance of \(T\)
remove the common \(h_s\).  Equation (PGR.8) cancels the two
remaining factors.  The support is stable under these operations,
so the same assertions hold at points where the seed is zero.

The finite support is a union of right \(K\)-cosets.  Each coset
has a chosen decomposition, and on it (PGR.3) is a continuous
function of \(x\) times a locally constant finite character.
The group \(G_f/K\) is countable: a second-countable space has
at most countably many disjoint open cosets.  Thus the piecewise
formula is measurable; its value outside the union is zero.
Here \(K\) is compact open, with positive finite mass, by the
universal residue-frame argument ASG14.  That argument applies
because PGR1 supplies exactly the original character properties;
it is independent of a choice of construction for \(X\).

For a fixed finite component, the seed is a scalar multiple of a
continuous linear functional applied to \(W_u(x)\).
Apply that functional to the Bochner-integrable reproduction formula
in \(X.Wu\_\allowbreak spec\).  A continuous linear functional commutes with
the integral and preserves integrability.  This gives precisely the
scalar reproduction clause of \texttt{Ph\allowbreak Props}, including zero
outside the finite support.  No assertion that the reproducing
kernel is compactly supported is needed here.
\end{proof}
% END PGR7

% BEGIN PGR8
\paragraph{PGR8. Rational unipotent approximation and strict column domains.}
\label{claim:PGR8}
The subgroup \(N_u(F)\) is dense in \(N_u(\mathbb A_f)\).
Consequently \(N_u(\mathbb A_f)=N_u(F)U_f\) for every open
subgroup \(U_f\) of \(N_u(\mathbb A_f)\).  Each
\(N_{u,\rho}\), in its original infinite embedding, is countable
and discrete and admits a measurable strict fundamental domain
\(\Omega_{u,\rho}\) in \(G_\infty\).
\begin{proof}
In a basis adapted to the hyperbolic flag, the unipotent Lie algebra
is a finite-dimensional \(F\)-vector space of strictly triangular
matrices satisfying linear skew-Hermitian equations.  Exponential
and logarithm are finite polynomials, inverse to each other on this
space and \(N_u\).  To check this, logarithm respects the flag and
is zero on the successive diagonal quotients.  The unitary relation
\(g^*=g^{-1}\) gives \((\log g)^*=-\log g\); conversely the
exponential of a skew-Hermitian nilpotent matrix is unitary and has
the required flag and quotient action.  The finite polynomials are
defined over \(F\) and are continuous at every place.
Equivalently, the complete polynomial reconstruction
in RCV14 gives these same rational and adelic points.  Finite adelic
approximation in a rational basis of the Lie algebra therefore proves
density.  An open coset \(nU_f\) meets \(N_u(F)\), proving the
factorization assertion.

The group \(N_{u,\rho}\) preserves the full row lattice.
Its matrix entries in an integral basis lie in a fixed discrete
integer lattice inside the product of the real matrix spaces.
They are therefore discrete at infinity, and countable.  More
generally any discrete subgroup \(L\) of a second-countable locally
compact group has a measurable strict domain for left translation.
To see this, choose a countable cover by open sets so small that
distinct \(L\)-translates of each are disjoint.  From the first
open set take all its points; from each later set remove all translates
of the points already taken.  The remaining pieces are Borel, meet
each orbit once, and their countable union is the required domain.
Apply this construction in \(G_\infty\).

The same elementary reasoning supplies fundamental domains for the
countable rational subgroups needed in adelic unfolding.  The
diagonal rational group is discrete in \(G_{\mathbb A}\): bound
finite denominators in a compact finite neighborhood, then use the
discreteness of its integral coordinates in all infinite embeddings.
The supplied quotient domain \(X.B.fund\) already has the original
almost-everywhere fundamental-domain property, which is sufficient
for every integral below.  We neither replace it nor normalize its
measure anew.
\end{proof}
% END PGR8

% BEGIN PGR9
\paragraph{PGR9. The exact column unipotents and their phases.}
\label{claim:PGR9}
For every rational \(n\in N_u(F)\) and every column \(\rho\),
\[
 n\in N_{u,\rho}\quad\Longleftrightarrow\quad
             a_{u,\rho}^{-1}n_fa_{u,\rho}\in K.
\tag{PGR.18}
\]
When these conditions hold,
\[
 \chi(a_{u,\rho}^{-1}n_fa_{u,\rho})=\psi_{u,f}(n_f)
                              =\psi_{u,\infty}(n_\infty)^{-1}.
\tag{PGR.19}
\]
Thus the subgroup here is exactly the source's \texttt{NOcol},
with its unique rational lift, not a finite-index replacement.
\begin{proof}
Preservation of \(\Lambda_{u,\rho}\) by a rational isometry is
equivalent to preservation of its completion.  This follows by
expressing a rational vector in an integer basis: being integral
in every finite completion is equivalent to having integer rational
coordinates.  The converse for the completion follows by scalar
extension.  Since \(x_u\) preserves \(\widehat L_0\), the
left side of (PGR.18) says
\(h_{u,\rho}^{-1}n_fh_{u,\rho}\in K_0\).
This is an element of \(N_u\) by PGR2.  PGR3 proves the required
membership in \(K\) and the first equality in (PGR.19), using
invariance of the unipotent character under \(h_{u,\rho}\).
Conversely \(K\subset K_0\) proves lattice preservation.
The last equality is (PGR.8).  The rational embedding into either
the finite adeles or all infinite places is injective, so there is
exactly one rational lift of the infinite element.  The original
definition of \(N\!\operatorname{Lam}\) is unipotent flag
preservation together with lattice preservation, giving the claimed
literal identification.
\end{proof}
% END PGR9

% BEGIN PGR10
\paragraph{PGR10. The support cells have exactly the original product measure.}
\label{claim:PGR10}
For arbitrary strict domains \(\Omega_{u,\rho}\) as in PGR8,
the sets
\[
              \Omega_{u,\rho}\times a_{u,\rho}K
\tag{PGR.20}
\]
form a disjoint set of representatives for the support of the seed in
\(S_u(F)\backslash G_{\mathbb A}\).  If a measurable
\(S_u(F)\)-invariant function \(f\) is supported there, its
quotient integral is the sum of its integrals over (PGR.20), first
for nonnegative \(f\), and then for every integrable complex \(f\).
If \(f\) is also right \(K\)-invariant, each summand is
\[
                  v_K\int_{\Omega_{u,\rho}}f(x,a_{u,\rho})\,d\mu_\infty(x).
\tag{PGR.21}
\]
\begin{proof}
For coverage, start with a finite component \(nhx_u\kappa\).
Reduce \(h\) by its supplied rational Levi double class, then
use PGR8 with
\(U_f=N_u(\mathbb A_f)\cap a_{u,\rho}Ka_{u,\rho}^{-1}\)
to remove its finite unipotent component by a rational unipotent.
The finite component is now in \(a_{u,\rho}K\).  Finally move
its infinite component into \(\Omega_{u,\rho}\) using
\(N_{u,\rho}\); (PGR.18) keeps the finite component in that
same coset.

For uniqueness, suppose a rational \(s\in S_u(F)\) carries
a point of the \(\rho\)-cell to a point of the \(\sigma\)-cell.
Write \(s=nh\) by PGR2.  The equality of finite components,
conjugated by the two \(h\)-representatives, puts the resulting
unipotent--Levi product in \(x_uKx_u^{-1}\).
PGR4 shows that its Levi factor lies in \(K_{H,u}\).
Uniqueness of the original column classes forces \(\rho=\sigma\).
The remaining rational Levi element stabilizes that column, hence
is the identity by PGR5.  The residual \(s\) is in
\(N_{u,\rho}\), by (PGR.18); strictness of the infinite domain
then gives \(s=1\).

To prove the integral statement, intersect (PGR.20) with the
countably many translates of any supplied adelic fundamental domain.
The pieces are null measurable and partition both systems of
representatives, up to the allowed null sets.  For integration replace
the supplied adelic domain by a measurable representative; its difference
and the countably many Haar translates of that difference are null.
This is the measurable-representative convention detailed in PGR11.
Left Haar invariance and invariance of
\(f\) identify their integrals.  Tonelli proves the assertion for
nonnegative functions without a prior finiteness assumption; apply
it to \(|f|\) before making the complex-valued change of domain.
Finally \(\mu_{\mathbb A}=\mu_\infty\times\mu_f\), and
\(\mu_f(a_{u,\rho}K)=v_K\).  Product integration and finite
right invariance give (PGR.21).  A rational unipotent also translates
the point of \(K\); this causes no extra factor because \(\mu_f\)
is left invariant.  There is no column-stabilizer denominator, precisely
because PGR5 proved it is one.
\end{proof}
% END PGR10

\subsection{Absolute unfolding before the orbital series is summed}

% BEGIN PGR11
\paragraph{PGR11. The full folding inequality and its compact finiteness bound.}
\label{claim:PGR11}
Let \(G\) be a second countable locally compact topological group,
\(\mu\) a left Haar measure, \(\Gamma\) a discrete subgroup,
and \(S\leq\Gamma\).  Let \(D\) be any left-coset
representatives for \(S\backslash\Gamma\), and \(F_S\) any
original almost-everywhere fundamental domain for the left action
of \(S\).  Thus \(F_S\) need only be null measurable; covering
and disjointness of its translates are allowed to hold almost everywhere.
For a measurable set \(\Omega\), put
\[
 m_\Omega(x)=\sup_{y\in G}
       \#\{\gamma\in\Gamma:x^{-1}\gamma y\in\Omega\}.
\tag{PGR.22}
\]
Counts and the supremum may be infinite.  If a measurable,
left \(S\)-invariant scalar function \(\phi\) satisfies
\[
 |\phi(y)|\leq a\int_\Omega|\phi(y\omega)|\,d\mu(\omega),
 \qquad a\geq0,
\tag{PGR.23}
\]
then, as an inequality of extended nonnegative integrals,
\[
 \sum_{\delta\in D}|\phi(\delta x)|
       \leq a\,m_\Omega(x)\int_{F_S}|\phi(y)|\,d\mu(y).
\tag{PGR.24}
\]
This is the full original \texttt{Fold\allowbreak Bound}, including arbitrary
measurable \(\Omega\), without a Hausdorff assumption.
If in addition \(G\) is Hausdorff and \(\Omega\) is compact
and contains \(1\), then \(m_\Omega\) is finite, Borel measurable and at
least one.  In particular, with
\[
 A_\phi=\int_{F_S}|\phi|,\qquad
 J_\phi=\int_{F_S}m_\Omega|\phi|<\infty,
\tag{PGR.25}
\]
we have \(A_\phi\leq J_\phi\) and the sum in (PGR.24)
is finite at every \(x\).
\begin{proof}
A discrete subgroup of a second countable group is countable.
The domain convention causes no extra regularity assumption.
Replace each null-measurable domain by a measurable representative
modulo \(\mu\).  Every Haar translate of the difference is null.
The union of these differences, the exceptional set for coverage,
and the intersections witnessing failure of disjointness is still
null, because the acting groups are countable.  The translated
representatives therefore give the required partitions off one null
set, with all original set integrals unchanged.  The domain partitions
below, and in PGR12, are understood in this sense.  This does not
weaken the separate strict \texttt{Is\allowbreak SFund} statements in this chapter.
Apply (PGR.23) at \(\delta x\), and change variable
\(y=\delta x\omega\), which uses only left invariance of
\(\mu\).  Sum the resulting nonnegative integrals.  Partition
\(G\) by \(sF_S\), \(s\in S\), and use the invariance
of \(|\phi|\).  The result is at most
\[
 a\int_{F_S}\sum_{\delta\in D,\ s\in S}
    1_\Omega(x^{-1}\delta^{-1}s y)\,|\phi(y)|\,d\mu(y).
\tag{PGR.26}
\]
Every \(\gamma\in\Gamma\) has exactly one expression
\(\delta^{-1}s\): apply the given coset decomposition to
\(\gamma^{-1}\).  Thus the sum of indicators in (PGR.26)
is the count occurring in (PGR.22), for this particular \(y\),
and is bounded by \(m_\Omega(x)\).  This proves (PGR.24).
All interchanges so far are Tonelli interchanges of nonnegative
measurable functions.  No integrability or right invariance of Haar
measure was required.

For the remaining assertions assume \(G\) is Hausdorff and
\(\Omega\) compact.  The discrete subgroup \(\Gamma\) is
then closed and meets compact sets in finite sets.
If one element \(\gamma_0\) is
counted at \((x,y)\), every other counted element satisfies
\[
 \gamma\gamma_0^{-1}\in
                     \Gamma\cap x\Omega\Omega^{-1}x^{-1}.
\]
The latter set is finite and is independent of \(y\).  Its
cardinality bounds \(m_\Omega(x)\).  Choosing \(y=x\)
and \(\gamma=1\) proves \(m_\Omega(x)\geq1\).
For completeness, its measurability can be seen without taking an
uncontrolled supremum of Borel functions.  The condition
\(m_\Omega(x)\geq q\), for a positive integer \(q\), is the
union over the countably many distinct tuples
\(\gamma_1,\ldots,\gamma_q\) of the conditions that some
\(\omega\in\Omega\) satisfies
\[
       x^{-1}\gamma_i\gamma_1^{-1}x\omega\in\Omega
                                  \quad(1\leq i\leq q).
\]
For each tuple this is closed: it is the projection of a closed subset
of \(G\times\Omega\), with compact second factor.  The union
is Borel.  This proves the stated measurability.  The final assertions
follow from (PGR.24), (PGR.25), and \(m_\Omega\geq1\).
\end{proof}
% END PGR11

% BEGIN PGR12
\paragraph{PGR12. The original mixed Poincar\'e identity, with absolute control.}
\label{claim:PGR12}
In the compact Hausdorff situation of PGR11, let \(\phi'\) be a
left \(S'\)-invariant seed satisfying the same conditions with
constant \(a'\), domain \(F_{S'}\), and finite quantities
\(A_{\phi'},J_{\phi'}\).  Use any respective representatives
\(D,D'\), and write
\[
 P\phi(x)=\sum_{\delta\in D}\phi(\delta x),\qquad
 P^{\rm abs}\phi(x)=\sum_{\delta\in D}|\phi(\delta x)|.
\]
Both series converge at every point and are independent of the
representatives.  For every \(\Gamma\)-fundamental domain
\(F_\Gamma\),
\[
 \begin{split}
 \int_{F_S}|\phi|P^{\rm abs}\phi'
       &\leq a'A_{\phi'}J_\phi<\infty,\\
 \int_{F_{S'}}|\phi'|P^{\rm abs}\phi
       &\leq aA_\phi J_{\phi'}<\infty,\\
 \int_{F_\Gamma}|P\phi|^2
       &\leq a A_\phi J_\phi\leq aJ_\phi^2.
 \end{split}
\tag{PGR.27}
\]
The functions \(P\phi,P\phi'\) are measurable and belong to
the indicated quotient \(L^2\).  If \(R\) is any double-coset
representative set for \(S\backslash\Gamma/S'\), put
\(S_r=S\cap rS'r^{-1}\), and choose arbitrary fundamental
domains \(F_r\) for these groups, with the same null-measurable,
almost-everywhere convention as in PGR11.  Then
\[
 \sum_{r\in R}\int_{F_r}|\phi(x)|\,|\phi'(r^{-1}x)|\,d\mu(x)
                                                   <\infty,
\tag{PGR.28}
\]
and
\[
 \begin{split}
 \int_{F_\Gamma}\overline{P\phi(x)}P\phi'(x)\,d\mu(x)
   &=\int_{F_S}\overline{\phi(x)}P\phi'(x)\,d\mu(x)\\
   &=\sum_{r\in R}\int_{F_r}
          \overline{\phi(x)}\phi'(r^{-1}x)\,d\mu(x).
 \end{split}
\tag{PGR.29}
\]
Taking complex conjugates gives the source's other inner-product
orientation.  Thus this proves the full original
\texttt{Poincare\allowbreak Series\allowbreak Finite}, \texttt{Poincare\allowbreak Series\allowbreak Mem\allowbreak Lp}, and
\texttt{Poincare\allowbreak Series\allowbreak Inner} at their stated general group scopes.

For the actual setting, \texttt{Fold\allowbreak OK} supplies a domain \(F_S\)
and a compact \(\Omega_0\ni1\), with one constant \(a\)
before all column vectors \(T\).  It supplies (PGR.23) for every
compact \(\Omega\supset\Omega_0\), and the finiteness of
\(J_{\Phi_u(T)}\) for every compact \(\Omega\) and every
\(S_u(F)\)-fundamental domain.  For two words we take the one
common compact set \(\Omega_0\cup\Omega'_0\).  All the
conclusions above therefore apply to their seeds for every \(T,T'\).
\begin{proof}
PGR11 gives pointwise absolute convergence.  Changing a representative
multiplies it on the left by an element of the seed's stabilizer,
which changes no summand.  Right multiplication by \(\Gamma\)
permutes the left cosets.  Consequently both periodized functions,
including their absolute versions, are left \(\Gamma\)-invariant.
They are measurable as countable sums of measurable functions.

Multiply the pointwise bound
\(P^{\rm abs}\phi'\leq a'm_\Omega A_{\phi'}\)
by \(|\phi|\) and integrate over \(F_S\).  This is the
first inequality of (PGR.27); interchange the seeds for the second.
To unfold, partition each translate of \(F_\Gamma\) by its
intersection with \(F_S\), or equivalently tile an
\(S\)-domain by the translates \(\delta F_\Gamma\),
\(\delta\in D\).  For every nonnegative left
\(\Gamma\)-invariant measurable \(f\), this gives
\[
 \int_{F_\Gamma}P^{\rm abs}\phi(x) f(x)\,d\mu(x)
                 =\int_{F_S}|\phi(x)|f(x)\,d\mu(x).
\tag{PGR.30}
\]
This is a nonnegative identity before it is an integrability assertion.
Apply it with \(f=P^{\rm abs}\phi'\).  The first bound of
(PGR.27) proves that all scalar products to be unfolded are dominated
by an integrable function.  Applying the same reasoning to a seed
paired with itself proves the third bound, hence its \(L^2\)
membership.  Dominated integration of the now absolutely convergent
scalar sums proves the first identity in (PGR.29).

For the second identity, replace the left-coset index \(\delta'\)
by its inverse, a right-coset index in \(\Gamma/S'\).  The left
action of \(S\) on these right cosets has orbits indexed by
\(R\), and the stabilizer of \(rS'\) is \(S_r\).
Choose representatives \(s\) of \(S/S_r\).  The corresponding
summand is \(\phi'((sr)^{-1}x)\).  In its integral over
\(F_S\), set \(y=s^{-1}x\); left \(S\)-invariance of
\(\phi\) gives the integrand
\(\overline{\phi(y)}\phi'(r^{-1}y)\).
The sets \(s^{-1}F_S\), for these representatives, form an
\(S_r\)-domain.  Partitioning intersections with the arbitrary
given \(F_r\) changes neither the nonnegative integral nor the
integrable scalar integral.  First doing this with absolute values
identifies the sum in (PGR.28) with the first integral in (PGR.27).
It is finite.  Every term is therefore integrable, and the same
partition proves the complex identity.  This also explains why an
arbitrary choice of the domains \(F_r\) is permitted.

Finally, the definition of \texttt{Fold\allowbreak OK} has precisely the
uniformity stated above.  The union of the two supplied compact sets
is compact, contains one, and permits both averaging inequalities and
both weighted integrability bounds.  No pointwise decay exponent or
arithmetic restriction on the words has been used.
\end{proof}
% END PGR12

% BEGIN PGR13
\paragraph{PGR13. Absolute slab integrability and single unfolding in the actual measure.}
\label{claim:PGR13}
For every pair of column vectors \(T,T'\) and every family of
strict domains \(\Omega_{u,\rho}\) for \(N_{u,\rho}\),
\[
 \int_{\Omega_{u,\rho}}|\Phi_u(T)(x,a_{u,\rho})|
       \sum_{\delta\in D_{u'}}
          |\Phi_{u'}(T')(\delta(x,a_{u,\rho}))|
                     \,d\mu_\infty(x)<\infty.
\tag{PGR.31}
\]
Thus the full original \texttt{Slab\allowbreak Abs} holds on the supplied
setting.  In the actual Hilbert space selected by \texttt{In\allowbreak HH},
\[
 \langle P_uT,P_{u'}T'\rangle
   =v_K\sum_{\rho\in I_u}\int_{\Omega_{u,\rho}}
       \overline{\langle J_uT(\rho),W_u(x)\rangle}
          P_{u'}T'(x,a_{u,\rho})\,d\mu_\infty(x).
\tag{PGR.32}
\]
There is exactly one finite-level factor \(v_K\).
\begin{proof}
Take \(\phi=\Phi_u(T)\), \(\phi'=\Phi_{u'}(T')\)
in PGR12.  The nonnegative function
\(|\phi|P^{\rm abs}\phi'\) is left \(S_u(F)\)-invariant,
measurable, and supported on the first seed's finite support.  Both
factors acquire unit-modulus characters under right multiplication by
\(K\), so their absolute values are right \(K\)-invariant.
Its integral over an \(S_u(F)\)-domain is finite by (PGR.27).
PGR10 identifies that integral with \(v_K\) times the sum of
the nonnegative integrals in (PGR.31).  There are finitely many columns,
and \(0<v_K<\infty\) by ASG14.  Each integral is finite.

Now apply the same change of domain to
\(\overline{\phi}P\phi'\).  Its integrability has just been
proved; its right \(K\)-characters cancel.  PGR6 identifies the
first seed on the cell with
\(\langle J_uT(\rho),W_u(x)\rangle\).
The first identity of (PGR.29) gives (PGR.32).  The quotient measure
here is the original measure obtained from \(X.B.fund\), which is
an allowed \(\Gamma\)-fundamental domain.  The supplied
\texttt{In\allowbreak HH} identifies its measurable scalar periods with the
classes used to define \(P_u\) and \(P_{u'}\) in that Hilbert
space.  Equality is consequently equality of the original inner
products, rather than of an independently normalized quotient model.
\end{proof}
% END PGR13

\subsection{Rational returns and the literal orbital terms}

% BEGIN PGR14
\paragraph{PGR14. The return dictionary includes the full original flag predicate.}
\label{claim:PGR14}
Fix columns \(\rho\in I_u\), \(\rho'\in I_{u'}\), and put
\[
 \Gamma_{\rho,\rho'}=
 \{r\in\Gamma:k_r=a_{u,\rho}^{-1}r_fa_{u',\rho'}\in K\}.
\tag{PGR.33}
\]
Every such \(r\) is an original \texttt{GCar}: its rational
isometry \(\sigma_r\) carries \(\Lambda_{u',\rho'}\)
onto \(\Lambda_{u,\rho}\), and its transported flag satisfies
the full original \texttt{PD}, namely \texttt{Is\allowbreak Lift},
\texttt{Is\allowbreak Admissible}, and \texttt{Frame\allowbreak OK}.  For these columns,
the active original terms are exactly
\[
       N_{u,\rho}\backslash
                 \Gamma_{\rho,\rho'}/N_{u',\rho'}.
\tag{PGR.34}
\]
Here ``active'' has its literal meaning: both row lattices match the
columns and the prescribed representative has \(k_r\in K\).
No term is restricted by a nondegeneracy condition on its Gram matrix.
\begin{proof}
The original rational coordinates identify \(r\) with a unique
\(E\)-linear isometry \(\sigma_r\).  The identity
\(r_f a_{u',\rho'}=a_{u,\rho}k_r\), with
\(k_r\widehat L_0=\widehat L_0\), gives equality of the
completed lattices.  The integral-coordinate argument of PGR9 turns
this into equality of the rational lattices themselves.  Thus the
lattice part of \texttt{GCar} holds in both directions.

We spell out the flag conditions, since a bare lattice isometry is
not the source's return carrier.  The supplied row condition for
\(\Lambda_{u',\rho'}\) includes that its negative vector belongs
to that lattice, that its chosen lift is this vector, the prescribed
local convention for the isotropic frame, the integral primitivity
condition on the negative vector in the quotient, and the equality
\(X\Lambda=XS\) with the prescribed different-weighted span.
Applying \(\sigma_r\) preserves all these conditions.  For the
local convention this follows by transporting each localized spanning
identity with the same nonzero coefficients; coefficients required
to be one at primes above \(g\) remain one.  Intersections with
the frame span and its perpendicular, sums with that span, and the
quotient condition defining \(L_{\rm sub}\) are carried to their
corresponding objects because \(\sigma_r\) is an isometry.
Consequently any scalar whose multiple of the transported negative
vector belongs to \(L_{\rm sub}\) is integral exactly when it was
integral before transport.  The different-weighted span transports
with its same coefficient ideals.  These observations prove both
\texttt{Is\allowbreak Lift} and every field of \texttt{Frame\allowbreak OK}.

For \texttt{Is\allowbreak Admissible}, write \(\mathcal F_u\) for the
ordered frame \((e_0,\ldots,e_{r-1},z_u)\).  Above every prime
dividing \(g\), the supplied congruences for \(x_u\) and the
fact that \(h_{u,\rho}\) fixes \(\mathcal F_u\) identify
\(a_{u,\rho}^{-1}\mathcal F_u\) with the reference frame
modulo \(g\widehat L_0\).  The same holds for \(u',\rho'\).
The defining conditions of \(K_2\) say that \(k_r\) acts
upper unitriangularly on the reference \((e,z)\)-frame modulo
\(g\widehat L_0\).  Hence it supplies coefficients, integral
in the finite adeles, with zero entries below the diagonal and one
on the diagonal, for
\[
 a_{u,\rho}^{-1}
       \bigl(\sigma_r\mathcal F_{u'}-\mathcal F_u\xi\bigr)
                 \in g\widehat L_0
                        \quad\hbox{at all primes above }g.
\tag{PGR.35}
\]
The finitely many off-diagonal residue coefficients can be lifted
simultaneously to \(\mathcal O_E\) by Chinese remaindering.
Set the lower entries to zero and the diagonal entries to one
exactly; call the resulting matrix \(\xi\) as well.  Replacing
the finite coefficients by these lifts does not change (PGR.35).

For a rational vector \(v\), the condition
\(a_{u,\rho}^{-1}v\in g\widehat L_0\) at the primes
above \(g\) says precisely that
\[
 v\in(g\Lambda_{u,\rho})_{\mathfrak p}
                         \quad\hbox{for every }\mathfrak p\mid g.
\tag{PGR.36}
\]
This follows by writing coordinates in a rational lattice basis and
localizing; equivalently use the localized congruence interpretation
of ASG11.  The statement (PGR.36), not a stronger global
\(g\Lambda\)-membership assertion, is the source's
\texttt{Mod\allowbreak G}.  Thus (PGR.35) with this one integral upper
unitriangular \(\xi\) proves every field of
\texttt{Is\allowbreak Admissible}.  The entire \texttt{PD} predicate is proved.

Finally, the source relation on two carriers with the same rows is
exactly \(r' = n r n'\), with
\(n\in N\!\operatorname{Lam}(\Lambda_{u,\rho})\) and
\(n'\in N\!\operatorname{Lam}(\Lambda_{u',\rho'})\).
PGR9 identifies these groups, including their unique rational lifts,
with the two groups in (PGR.34).  It also proves that both left and
right multiplication preserve \(k_r\in K\).  Hence activity
does not depend on the chosen carrier representative, and the
quotient is exactly (PGR.34).  Conversely every active original
carrier has its specified rational return in (PGR.33), so nothing
has been lost from this quotient dictionary.
\end{proof}
% END PGR14

% BEGIN PGR15
\paragraph{PGR15. Pointwise support reindexing, with the exact finite phase.}
\label{claim:PGR15}
For every \(T'\), \(\rho\), and \(x\in G_\infty\),
\[
 P_{u'}T'(x,a_{u,\rho})
  =\sum_{\rho'\in I_{u'}}
       \sum_{[r]\in\Gamma_{\rho,\rho'}/N_{u',\rho'}}
         \chi(k_r)^{-1}
           \langle J_{u'}T'(\rho'),W_{u'}(r_\infty^{-1}x)\rangle.
\tag{PGR.37}
\]
The same reindexing holds with absolute values, and both sums are
absolutely convergent.  The right-coset index is also exactly the
original \texttt{Col\allowbreak Cl}, originally defined using \(S_{u'}(F)\).
Thus this is the full original \texttt{Support\allowbreak Reindex}, with its
norm-sum assertion as well as its scalar-sum assertion.
\begin{proof}
Invert a left-coset index \(\delta\) in the original period, so
its summand is \(\Phi_{u'}(T')(r^{-1}(x,a_{u,\rho}))\).
The support of this seed depends only on its finite coordinate.
Whenever this coordinate is in the support, the finite part of the
coverage argument in PGR10 supplies a unique column \(\rho'\)
and a rational element of \(S_{u'}(F)\) by which to modify \(r\)
on the right so that
\(r_f^{-1}a_{u,\rho}\in a_{u',\rho'}K\).
Equivalently \(r\in\Gamma_{\rho,\rho'}\).
Left rational invariance of the seed, proved in PGR7, means this
modification does not alter its value at any infinite coordinate.
Cosets outside the support contribute zero and are simply omitted.

Suppose two returns in \(\Gamma_{\rho,\rho'}\) differ on the
right by \(s'\in S_{u'}(F)\).  Their finite relation implies
\(a_{u',\rho'}^{-1}s'_f a_{u',\rho'}\in K\).
Split \(s'=n'h'\) rationally.  Integral Levi extraction, PGR4,
puts the Levi factor in the conjugate compact level of that column.
PGR5 forces its rational Levi part to be one.  PGR9 then says
\(s'\in N_{u',\rho'}\).  The converse follows immediately
from PGR9.  Distinct columns cannot be identified, by uniqueness of
the original Levi column class.  This proves the claimed bijection
of support indices, including the exact \texttt{Col\allowbreak Cl} relation.

At a return in (PGR.33) one has the literal finite equality
\(r_f^{-1}a_{u,\rho}=a_{u',\rho'}k_r^{-1}\).
Evaluate the seed using \(n'=1\), \(h'=h_{u',\rho'}\),
and \(\kappa'=k_r^{-1}\).  Formula (PGR.3) gives precisely
the summand on the right of (PGR.37).
Changing \(r\) to \(rn'\), \(n'\in N_{u',\rho'}\),
multiplies \(\chi(k_r)^{-1}\) by
\(\psi_{u',\infty}(n')\), by (PGR.19), and multiplies
the wave-function value by its inverse.  The value is unchanged.
The bijection therefore preserves the scalar summands and their
absolute values.  Pointwise absolute convergence follows from PGR12
(and is also part of the retained \texttt{In\allowbreak HH}).
\end{proof}
% END PGR15

% BEGIN PGR16
\paragraph{PGR16. Orbit unfolding with arbitrary original sections.}
\label{claim:PGR16}
The left action of \(N_{u,\rho}\) on the right-coset set in
(PGR.37) has the active original terms as its orbit set.  If \(r\)
represents a term \(t\), its stabilizer is exactly
\[
 R_r=N_{u,\rho}\cap rN_{u',\rho'}r^{-1},
\tag{PGR.38}
\]
viewed in \(G_\infty\); this is the original
\texttt{ROinf}.  Define
\[
 f_r(x)=\overline{\langle J_uT(\rho),W_u(x)\rangle}
       \chi(k_r)^{-1}
       \langle J_{u'}T'(\rho'),W_{u'}(r^{-1}x)\rangle.
\tag{PGR.39}
\]
Whenever the corresponding absolute slab integral is finite,
\[
 \int_{\Omega_{u,\rho}}
       \sum_{[r]\in\Gamma_{\rho,\rho'}/N_{u',\rho'}}
                 f_r(x)\,d\mu_\infty(x)
   =\sum_{t\ {\rm active}}\int_{F_t} f_{r_t}(x)\,d\mu_\infty(x),
\tag{PGR.40}
\]
where each \(r_t\) is the original representative and each
\(F_t\) is any prescribed strict fundamental domain for its
\(R_{r_t}\).  The integrand on the left and every term on the
right are integrable.  The same equality holds for the nonnegative
sum of their absolute values.  In particular the sum of the absolute
orbital integrals is finite.  This proves the full original
\texttt{Orbit\allowbreak Unfold} in this setting, using its stated absolute
slab premise; PGR13 supplies that premise in the application.
\begin{proof}
The orbit and stabilizer identifications are now algebraic.
Two right cosets lie in the same left orbit if and only if
their returns differ by \(r\mapsto n r n'\), precisely the
relation in PGR14.  A left multiplier fixes \(rN_{u',\rho'}\)
if and only if it belongs to (PGR.38).  Unique rational lifting,
PGR9, identifies this intersection with the source's definition in
infinite coordinates.

We check the scalar covariance that makes the change of domains valid.
The right multiplication calculation in PGR15 gives
\(f_{rn'}=f_r\).  If \(n\in N_{u,\rho}\), (PGR.19)
and (PGR.1) give
\[
                  f_{nr}(x)=f_r(n^{-1}x).
\tag{PGR.41}
\]
Indeed \(\chi(k_{nr})^{-1}=\psi_{u,\infty}(n)
\chi(k_r)^{-1}\), whereas
\(\overline{\langle J_uT(\rho),W_u(n^{-1}x)\rangle}
=\psi_{u,\infty}(n)
\overline{\langle J_uT(\rho),W_u(x)\rangle}\).
These are exactly the same factor.  For \(n\in R_r\), the
first covariance and (PGR.41) show that \(f_r\) is left
\(R_r\)-invariant.  All these identities hold also for absolute
values.

Choose representatives \(n\) of \(N_{u,\rho}/R_r\).
Integration of (PGR.41) over \(\Omega_{u,\rho}\), followed by
\(y=n^{-1}x\), tiles an \(R_r\)-fundamental domain by the
sets \(n^{-1}\Omega_{u,\rho}\).  Apply this first to
\(|f_r|\).  Tonelli, the countability of the rational group, and
the assumed absolute slab bound show
\[
 \sum_{t\ {\rm active}}\int_{F_t}|f_{r_t}|<\infty.
\tag{PGR.42}
\]
Changing any tiled domain to an arbitrary strict domain \(F_t\)
is legitimate by the same measurable intersection partition as in
PGR10.  Only left Haar invariance is used.  Finite (PGR.42) permits
the complex-valued Fubini interchange and proves (PGR.40).

One can choose the representative of each orbit to be the prescribed
representative of the original term: it is active by PGR14.
For clarity, the operator version of its choice independence has the
same cancellation.  Under \(r\mapsto n r n'\), changing
variable \(x=ny\) in the rank-one orbital integral multiplies
that operator by
\(\psi_{u,\infty}(n)\psi_{u',\infty}(n')\).
Its finite factor \(\chi(k_r)\) acquires the inverse multiplier,
by (PGR.19).  Thus the phase-weighted operator is unchanged.
The supplied \texttt{Amp\allowbreak Well\allowbreak Def} and \texttt{OOFund\allowbreak Indep}
assert precisely the permitted independence for every original
carrier and every original strict domain, including the prescribed
\texttt{Fund\allowbreak Of}.  We retain these original premises throughout;
our tiling does not replace either by a special selection of domains
or by a different quotient measure.
\end{proof}
% END PGR16

% BEGIN PGR17
\paragraph{PGR17. Exact term evaluation and the full entry identity.}
\label{claim:PGR17}
For an original term \(t\), keep its prescribed representative
\(r_t\) and its prescribed domain \(F_t=\texttt{Fund\allowbreak Of}(t)\).
Its operator is exactly
\[
 O_t=\int_{F_t}
      \bigl[\xi\longmapsto
       \langle W_{u'}(r_t^{-1}x),\xi\rangle W_u(x)\bigr]
                                         \,d\mu_\infty(x).
\tag{PGR.43}
\]
Let \(B_{t,\rho\rho'}\) be zero if the term is inactive for
these columns, and otherwise set
\[
 B_{t,\rho\rho'}=s_{u,u'}v_K\chi(k_{r_t})O_t.
\tag{PGR.44}
\]
These are the literal original \texttt{blk} operators.  Under
the retained \texttt{Abs\allowbreak Conv\allowbreak T}, every block series is absolutely
convergent in operator norm, and
\[
 s_{u,u'}\langle P_uT,P_{u'}T'\rangle
   =\sum_{\rho,\rho'}
      \left\langle
       \left(\sum_t B_{t,\rho\rho'}\right)J_{u'}T'(\rho'),
                         J_uT(\rho)\right\rangle
                   \quad\hbox{for all }T,T'.
\tag{PGR.45}
\]
This is the full original \texttt{Entry\allowbreak Unfold} on this setting.
\begin{proof}
The supplied \texttt{OOInt\allowbreak All} makes (PGR.43) a Bochner integral
of the original rank-one operator, in a finite-dimensional normed
operator space.  For fixed vectors \(v=J_uT(\rho)\) and
\(v'=J_{u'}T'(\rho')\), its matrix coefficient is
\[
 \left\langle O_t v',v\right\rangle
  =\int_{F_t}
       \langle v',W_{u'}(r_t^{-1}x)\rangle
       \langle W_u(x),v\rangle\,d\mu_\infty(x).
\tag{PGR.46}
\]
Indeed evaluation followed by this inner product is a continuous
real-linear functional of the operator, so commutes with the
Bochner integral.  Conjugate linearity in the first variable gives
\(\langle\langle W',v'\rangle W,v\rangle
=\langle v',W'\rangle\langle W,v\rangle\), proving (PGR.46).
In particular an active term satisfies
\[
 s_{u,u'}v_K\int_{F_t}f_{r_t}(x)\,d\mu_\infty(x)
       =\langle B_{t,\rho\rho'}v',v\rangle.
\tag{PGR.47}
\]
Here the scalar in (PGR.44) is conjugated on the left side of the
inner product: \(\overline{\chi(k_{r_t})}=\chi(k_{r_t})^{-1}\).
This accounts for the entire phase in (PGR.39).  The factor
\(s_{u,u'}v_K\) is real, so is not altered.  Equation (PGR.47)
is the exact original \texttt{Term\allowbreak Val}.

Insert the support identity PGR15 into (PGR.32), then apply the
orbit unfolding PGR16 separately for each pair of columns.  PGR13
supplies its absolute hypothesis for every such pair.  The resulting
scalar series is absolutely convergent by (PGR.42).  Formula
(PGR.47) identifies every term, including zero terms, with its
prescribed block coefficient.  Thus the left side of (PGR.45) equals
\(\sum_{\rho,\rho'}\sum_t
\langle B_{t,\rho\rho'}v',v\rangle\).

We now justify the operator sums in (PGR.45); this is a different
convergence issue from the absolute unfolding just completed.  Write
\(A_t\) for the original \texttt{a\allowbreak Term}, the finite block
operator on \(\bigoplus_{\rho'}V_{u'}\) with blocks
\(B_{t,\rho\rho'}\), and \(\operatorname{op}(t)\) for its
original Hilbert--Schmidt encoding.  The first clause of
\texttt{Abs\allowbreak Conv\allowbreak T} supplies
\(\sum_t\|\operatorname{op}(t)\|<\infty\).
In these fixed finite-dimensional spaces there is a finite constant
\(C\), independent of \(t\), with
\[
 \|B_{t,\rho\rho'}\|\leq\|A_t\|
                        \leq C\|\operatorname{op}(t)\|.
\tag{PGR.48}
\]
The first inequality composes with a norm-one coordinate inclusion
and projection (and is trivial for a zero coordinate space).  For
the second, choose orthonormal bases, recover the finitely many
matrix entries from the Hilbert--Schmidt encoding, and bound the
operator norm by the sum of their absolute values; Cauchy--Schwarz
supplies a fixed dimension-dependent constant.  This also proves
the assertion when an operator space is zero, without requiring a
positive dimension.  Hence both \(\sum_tA_t\) and every block
series converge absolutely.  Continuous evaluation, the inner
product, and the finite sums over columns commute with these
convergent sums.  This proves (PGR.45) with exactly its displayed
operator series.
\end{proof}
% END PGR17

% BEGIN PGR18
\paragraph{PGR18. The full original universal \texttt{KExpansion\allowbreak C}.}
\label{claim:PGR18}
For every original \(D\), every supplied \(X\), every
\(Y:\texttt{Full}(X)\), and every ordered pair of words
\(u,u'\), the original assumptions
\[
 \begin{gathered}
 \texttt{Fold\allowbreak OK}(u),\quad\texttt{Fold\allowbreak OK}(u'),\quad
 \texttt{In\allowbreak HH}(u),\quad\texttt{In\allowbreak HH}(u'),\\
 \texttt{Fund\allowbreak Exists},\quad\texttt{OOInt\allowbreak All},\quad
 \texttt{OOFund\allowbreak Indep},\quad\texttt{Amp\allowbreak Well\allowbreak Def},\quad
 \texttt{Abs\allowbreak Conv\allowbreak T}
 \end{gathered}
\tag{PGR.49}
\]
imply the original \texttt{KExpansion} for
\(X.Fu(u),X.Fu(u'),Y.cs(u),Y.cs(u'),Y.FN(u),Y.FN(u')\),
the given character \(Y.hK\), and the representatives
\(Y.Dt(u),Y.Dt(u')\).  More explicitly, there exists a
continuous complex-linear operator \(B_{u,u'}\) such that
\[
 \operatorname{HasSum}_{t}\,A_t=B_{u,u'},\qquad
 s_{u,u'}\langle P_uT,P_{u'}T'\rangle
       =\langle B_{u,u'}J_{u'}T',J_uT\rangle
                                  \quad\hbox{for every }T,T'.
\tag{PGR.50}
\]
Every object in (PGR.50) is the original supplied object.  The two
words may coincide; no lower bound on their weights or their gap
has been added.
\begin{proof}
Define \(B_{u,u'}=\sum_tA_t\).  Estimate (PGR.48) gives an
absolutely convergent series in the complete finite-dimensional
space of continuous operators.  Equivalently, its net of finite
partial sums converges, which is precisely the \texttt{Has\allowbreak Sum}
assertion.  Each block projection is continuous, so the blocks of
\(B_{u,u'}\) are the block sums in (PGR.45).  The inner product
of finite direct sums now turns (PGR.45) into the second identity
of (PGR.50).

We record explicitly how the scope matches the original interface.
The arbitrary \texttt{Full} supplies both column systems, both
row systems, both rational coset selections, and both unipotent
domains.  PGR1--PGR10 prove the required finite phase, splitting,
trivial column stabilizer, and support-cell statements for this
arbitrary setting.  The two \texttt{Fold\allowbreak OK} inputs give the common
averaging set and the absolute bound in PGR12--PGR13, with their
original quantifiers over all column vectors and all domains.
The two \texttt{In\allowbreak HH} inputs identify the periods in the actual
automorphic Hilbert space.  The four orbital inputs in (PGR.49)
are retained at their full original scopes, including every
carrier pair and every allowed section.  PGR14--PGR17 compare to
their prescribed representatives, domains and integrands.
Finally \texttt{Abs\allowbreak Conv\allowbreak T} supplies the operator-norm convergence
used here.  Its other clauses remain supplied rather than being
silently weakened in the target statement.

The normalization is exactly (PGR.5), using the given
\(\mu_N,\mu_H,F_N\) for each word.  Those measures and
domains were never replaced.  The unnormalized scalar identity
was proved first, so even if a real-valued mass in the literal
interface is zero, the total inverse convention simply makes both
normalized sides zero.  No proof step divides by that mass.
Only the finite level volume was divided out in PGR13, and its
strict positivity and finiteness hold for every supplied setting.
This verifies all the quantifiers in \texttt{KExpansion\allowbreak C}.
\end{proof}
% END PGR18

\subsection{Transport to the actual period operator}

% BEGIN PGR19
\paragraph{PGR19. Linearity, the antilinear Gram transport, and projected periods.}
\label{claim:PGR19}
For every word and every permitted \(D_u\), the actual
\texttt{Pu\allowbreak Linear} assertion holds: under its original
\texttt{In\allowbreak HH} premise, \(P_u\) is complex linear.
Give \(C_u\) its original \texttt{Col\allowbreak T} inner product from
PGR6, and regard \(P_u\) as a map into the actual ambient
quotient \(L^2\) space \(\mathcal H\).  Put
\[
 \mathcal A_u=(\sqrt{c_u})^{-1}P_u.
\tag{PGR.51}
\]
Under the hypotheses of PGR18, the exact Gram transport is
\[
 B_{u,u'}=J_u\,\mathcal A_u^*\mathcal A_{u'}\,J_{u'}^{-1}.
\tag{PGR.52}
\]
The two occurrences of antilinearity in (PGR.52) are essential.
For any subspaces \(E_u\subset C_u\), \(E_{u'}\subset C_{u'}\),
with inclusions \(\iota_u,\iota_{u'}\), and any orthogonal
projection \(Q\) on a closed subspace of \(\mathcal H\),
the Gram operator of the projected periods is exactly
\[
 \begin{split}
 (Q\mathcal A_u\iota_u)^*(Q\mathcal A_{u'}\iota_{u'})
   &=\iota_u^*\mathcal A_u^*\mathcal A_{u'}\iota_{u'}\\
   &\quad-\iota_u^*\mathcal A_u^*(1-Q)
                              \mathcal A_{u'}\iota_{u'}.
 \end{split}
\tag{PGR.53}
\]
Thus a subsequent target projection may be suppressed only after
its displayed correction term has been proved zero.
\begin{proof}
For any fixed finite support decomposition, (PGR.3) is linear in
\(T\), and PGR7 says that its value is independent of that
decomposition.  Off support it is zero.  Hence the actual scalar
seed is linear in \(T\) at every point.  Under \texttt{In\allowbreak HH},
all the scalar Poincar\'e sums converge absolutely.  One can
therefore add them and multiply them by arbitrary complex scalars
term by term.  Their \(L^2\) representatives satisfy the same
identities because \texttt{In\allowbreak HH} supplies their \(L^2\)
membership.  This proves the full implication defining
\texttt{Pu\allowbreak Linear}; no folding or gap hypothesis is needed for
this linearity argument.  Its range in the prescribed automorphic
subspace is the range assertion already supplied by \texttt{In\allowbreak HH}.

PGR6 makes \(C_u\) a finite-dimensional Hilbert space with
\(J_u\) an antilinear isometry onto the complete finite direct
sum of compact-type spaces.  Thus the linear map \(P_u:C_u\to
\mathcal H\) is continuous, and its Hilbert adjoint exists.
We use the ambient \(\mathcal H\) here; no unstated completeness
of a proper automorphic subspace is required.

Each real-valued Haar mass in (PGR.5) is nonnegative, as are
\(v_K\) and the natural-number dimension.  Total division
therefore gives \(c_u\geq0\).  For two such numbers the
square-root product and inverse identities, also when either is
zero, give
\[
 \langle\mathcal A_uT,\mathcal A_{u'}T'\rangle
             =s_{u,u'}\langle P_uT,P_{u'}T'\rangle.
\tag{PGR.54}
\]
Set \(C=\mathcal A_u^*\mathcal A_{u'}\).  The adjoint
identity, the original \texttt{Col\allowbreak T} inner product, and PGR18
give, for all \(T,T'\),
\[
 \langle J_u C T',J_uT\rangle
   =\langle T,CT'\rangle
   =\langle\mathcal A_uT,\mathcal A_{u'}T'\rangle
   =\langle B_{u,u'}J_{u'}T',J_uT\rangle.
\]
Surjectivity of \(J_u\) and nondegeneracy of the inner product
identify the first vectors.  Surjectivity of \(J_{u'}\)
proves (PGR.52).  In particular, this is not the formula obtained
by treating the Riesz map as complex linear.

For (PGR.53), use \(Q^*=Q=Q^2\) and the rule for the adjoint
of a composite.  The left side is
\(\iota_u^*\mathcal A_u^*Q\mathcal A_{u'}\iota_{u'}\).
Substitute \(Q=1-(1-Q)\).  This proves (PGR.53) with no
additional range assertion.  If \(Q\) fixes the relevant period
images, its correction vanishes; otherwise the correction must be
kept in any subsequent comparison.
\end{proof}
% END PGR19

% BEGIN PGR20
\paragraph{PGR20. The full original \texttt{KTie} consumer on the supplied finite datum.}
\label{claim:PGR20}
Retain an arbitrary original \(h_F:\texttt{Fin\allowbreak X}(Y)\), and let
\(s_Y=\texttt{s\allowbreak Of}(h_F)\) be its actual constructed spectral
datum.  For every pair of words, the original two
\texttt{Dim\allowbreak Tie} identities, \texttt{KExpansion}, and
\texttt{Abs\allowbreak Conv\allowbreak T} imply
\[
     \texttt{KExpansion}(s_Y,u,u')
                     \quad\hbox{and}\quad
     \texttt{KAbs\allowbreak Conv}(s_Y,u,u').
\tag{PGR.55}
\]
This is the full original \texttt{KTie}; in particular PGR18
now supplies its analytic expansion input at universal scope.
The fields of \texttt{Fin\allowbreak X} remain the original ones: the two
column bijectivity fields, \texttt{In\allowbreak HH} and period linearity
for all words, finite dimension of the actual automorphic space,
finite fibers of the map to lattice terms, pair independence,
the real-value field, the coefficient-scale supplier, and volume
positivity.  No existence assertion for \texttt{Fin\allowbreak X}, and no
isomorphism from a different spectral datum, is inferred here.
\begin{proof}
We give the transport and regrouping rather than treating
\texttt{KTie} as a consequence of notation.  The column spaces
of \(s_Y\) are the original \texttt{Col\allowbreak T} spaces, with the
same inner products as PGR6.  Its period maps are precisely
\(P_u\), now viewed in the supplied finite-dimensional actual
automorphic space.  Let \(B\) be the operator provided by the
assumed \texttt{KExpansion}, and let
\[
       \mathcal T(A)=J_u^{-1} A J_{u'}
\tag{PGR.56}
\]
for a continuous complex-linear operator between the two column
coordinate sums.  This is again a complex-linear operator on the
original \texttt{Col\allowbreak T} spaces.  As a function of \(A\),
\(\mathcal T\) is additive and real linear; it conjugates
complex scalar multiplication.  Both \(J\)'s are isometries
and surjective, so
\(\|\mathcal T(A)\|=\|A\|\).
In particular it is a continuous additive map and preserves
convergent sums.

The original normalizing constant of \(s_Y\) is
\(Y.\operatorname{volN}(u)Y.\operatorname{volH}(u)/
\operatorname{dimPi}(u)\).  Its numerator is exactly
\(v_K\mu_{N,u}(F_{N,u})\mu_{H,u}(H_{u,\infty})\),
with the same real-valued mass convention as before.  The
\texttt{Dim\allowbreak Tie} assumption identifies its denominator with
\(\dim_{\mathbb C}V_u\).  Thus it equals \(c_u\) exactly,
and likewise for \(u'\).  Nonnegativity identifies the source's
real power \((c_uc_{u'})^{-1/2}\) with
\((\sqrt{c_uc_{u'}})^{-1}\), including the zero convention.
Testing against arbitrary column vectors as in PGR19 now proves
\[
                 K_{s_Y}(u,u')=\mathcal T(B).
\tag{PGR.57}
\]
The finite-dimensionality supplied by \(h_F\) makes all
adjoints in this original definition legitimate.  This calculation
uses no identification of \(s_Y\) with the initial datum \(D.s\).

Let \(\mathfrak T\) be the original two-row orbital term set
and \(\mathcal R\) the original lattice-flag term set.  The
map \(q:\mathfrak T\to\mathcal R\) is the source's
\texttt{to\allowbreak Term}; it transports the primed flag and forgets the
second row presentation.  For each \(t\in\mathcal R\),
the actual \(s_Y\)-term operator is defined to be
\[
             A^R_t=\sum_{\tau:\ q(\tau)=t}\mathcal T(A_\tau).
\tag{PGR.58}
\]
This is the literal finite sum, since \(h_F.\texttt{fib\allowbreak Fin}\)
supplies finiteness for every fiber.  The absolute convergence
clause and (PGR.48) give
\(\sum_{\tau\in\mathfrak T}\|A_\tau\|<\infty\),
and hence the same assertion for \(\mathcal T(A_\tau)\).
Applying the continuous additive map \(\mathcal T\) to the
original \texttt{Has\allowbreak Sum} gives sum \(\mathcal T(B)\).
Partition this absolutely convergent series into its fibers under
\(q\).  The triangle inequality gives
\[
 \sum_{t\in\mathcal R}\|A^R_t\|
       \leq\sum_{\tau\in\mathfrak T}\|\mathcal T(A_\tau)\|<\infty,
 \qquad
       \operatorname{HasSum}_{t\in\mathcal R}A^R_t=\mathcal T(B).
\tag{PGR.59}
\]
Empty fibers contribute zero.  One may prove the regrouping directly
by choosing a finite set of \(\tau\)'s with arbitrarily small
norm tail; every finite regrouping containing their finitely many
images has the same bound on its remaining tail.  Thus the assertion
is independent of any enumeration or choice of fiber representatives.
Equations (PGR.57)--(PGR.59) are exactly the two conclusions in
(PGR.55).  The additional fields of \(\texttt{Fin\allowbreak X}\) remain
supplied for the construction of \(s_Y\), even where this
particular regrouping uses only a subset of them.
\end{proof}
% END PGR20

\subsection{What this supplies, and what remains}

PGR18 supplies the full original universal \texttt{KExpansion\allowbreak C},
and PGR20 supplies its original \texttt{KTie} transport with the
actual common data.  The proof has also reconstructed the finite
phase, integral Levi extraction, trivial column stabilizers, the
original flag return predicate, absolute unfolding, and the exact
matrix-coefficient evaluation.  Each of these intermediate arguments
is retained above so that the identity can be read without following
the source's chain of wrappers.

This is an identity under the original analytic inputs.  It does not
produce \texttt{Full}, all-word folding and automorphic-range
suppliers, \texttt{Abs\allowbreak Conv\allowbreak C}, or the remaining fields needed to
construct \texttt{Fin\allowbreak X}.  It does not construct a common-family
\texttt{SIso}, remove the projected-period correction in (PGR.53),
or establish an upper bound for the resulting Gram operator.
The exact archimedean Jacobian and limiting-integral producers
\texttt{Jac\allowbreak Gauge} and \texttt{Rhs\allowbreak WLim}, the independent product
measure statement \texttt{Prod\allowbreak Exact}, and the remaining mass,
counting and spectral comparisons stay separate.  No whole global
spectral obligation is closed by this chapter alone.


\section{Determinant jets, their short windows, and the actual integral row tests}
\label{sec:determinant-jet-v057}

\status{This chapter proves the full original
\texttt{Formal\allowbreak Series\allowbreak Bounds.\allowbreak Thm\allowbreak P}, the literal row identity
\texttt{Bi\allowbreak Value}, and \texttt{Bi\allowbreak Content} and \texttt{Bi\allowbreak Window}.
It then proves the original conversion from a supplied Cartier normal
form to its actual perturbation, including all four fields of
\texttt{Mto\allowbreak HM}, and constructs the integral germs and the exact
change-of-coordinates identity in \texttt{block\_\allowbreak of\_\allowbreak normal\_\allowbreak form}.
The final window uses the same normal form, coordinates and original
row gauge throughout.  Existence of the original normal form
\texttt{Row\allowbreak Data.\allowbreak Lemma\allowbreak M} is retained at the last application; it is not
replaced by a new assumption about an unrelated power series.
No whole global obligation is asserted here.}

We want an integral row test whose value is not too divisible by a
prime.  There are two separate issues.  The first is an estimate for
an explicitly perturbed exponential in matrix variables.  The second
is to show that this exponential is the restriction of the prescribed
row gauge, and that its derivatives are values on actual integral
germs.  A scalar estimate by itself would not settle the second issue.

The short-window argument in the first part was prompted by the
adjoint-functor reading edition, revision
\texttt{94107d5f1a355a4ae0bd521123dba9ffd62e00f5},
\texttt{tex/\allowbreak arithmetic.\allowbreak tex}, lines 6432--6528.  We give the coefficient,
factorial and interpolation details.  The remaining parts reconstruct
the original row and coordinate arguments.  All formal-source matches
use revision
\texttt{7cbf7a0c32fca1badf05931879a37c3f8a2b03b4}.
The classical ordinary Cayley identity is the one published input in
the determinant calculation; its exact statement is specified below.

\subsection{Conventions and the local window}

Let $K$ be a characteristic-zero field with an additive
nonarchimedean valuation
$v:K\longrightarrow\mathbb R\cup\{+\infty\}$.
Thus $v(0)=+\infty$, $v(xy)=v(x)+v(y)$, and
$v(x+y)\ge\min(v(x),v(y))$.  We impose neither completeness nor
discreteness.  Fix a prime $p\ge5$ with $v(p)=1$, an odd integer
$h\ge1$, and integers $n\ge1$, $0\le r\le n$.  Put
\begin{gather}
 q=p^h,\qquad \rho=\frac1{p^2-1},\qquad
 \gamma=\frac p{p^2-1},\qquad
 e=1+\frac{q-p}{p^2-1},\qquad c=-p^{-e},\label{djw:parameters}\\
 a=\rho+\frac1{2q}-\frac1{8q},\qquad
 b=\rho+\frac12-\frac1{8q}.
 \label{djw:weights}
\end{gather}
Writing $h=2u+1$ gives
$e=1+p\sum_{j<u}p^{2j}$, so $e$ is a positive integer.

For independent $r$ by $r$ matrices $X,Y$, write $|A|,|B|$ for
the total degrees of multiindices on their entries, and $A!,B!$ for
the products of their factorials.  Let
$U=\sum u_{A,B}X^AY^B$ satisfy
\begin{align}
 v(u_{A,B})+a|A|+b|B|&\ge\frac1{p-1},
       &&\text{for every }A,B,\label{djw:hm1}\\
 u_{A,B}\ne0&\Longrightarrow |B|\ge1.
       \label{djw:hm2}
\end{align}
The second condition says $U(X,0)=0$.  In particular its formal
exponential exists coefficient by coefficient.  All derivatives below
are ordinary derivatives, so a monomial of exponent $A$ contributes
the factorial $A!$, rather than a divided-power normalization.
Define
\begin{equation}
 \Lambda_y=
 \left.\det(\partial_X)^{qy}\det(\partial_Y)^y
 \exp\left(c\sum_{i,j}X_{ij}^{q}Y_{ij}+U(X,Y)\right)
 \right|_{X=Y=0}.
 \label{djw:lambda}
\end{equation}
For a polynomial $P$ and a series $F$ in the same variables, write
\[
 \langle P,F\rangle=P(\partial)F(0)
       =\sum_A P_A A!F_A.
\]
The sum is finite.  Empty determinants and empty products equal one.

% BEGIN DJW1
\paragraph{DJW1: integer and factorial valuations.}
\label{claim:DJW1}
The restriction of $v$ to nonzero integers is the ordinary
$p$-adic valuation.  If a multiindex $A$ has at most $r^2$ coordinates,
$|A|\le M$ and $M\ge1$, then
\begin{equation}
 v(A!)\ge\frac{|A|}{p-1}-r^2(\log_p M+1).
 \label{djw:multifactorial}
\end{equation}
For
\[
 C_r(m)=\prod_{t=1}^{m}\prod_{j=0}^{r-1}(t+j)
       =\prod_{j=0}^{r-1}\frac{(m+j)!}{j!},\qquad C_r(0)=1,
\]
and $r\ge1$, one has
\begin{equation}
 v(C_r(m))\le\frac{rm}{p-1}+r(\log_p r+1).
 \label{djw:capacity}
\end{equation}
Also $v(m!)\le m/(p-1)$ for every $m\ge0$.

\begin{proof}
Repeated addition gives $v(z)\ge0$ for every integer $z$.
If $p\nmid z$, Bezout supplies integers $u,w$ with $uz+wp=1$.
If $v(z)>0$, both terms on the left have positive valuation, a
contradiction.  Removing the power of $p$ from any nonzero integer
proves the first assertion, including its sign.

Counting the multiples of $p,p^2,\ldots$ in a factorial gives
\[
 v(m!)=\sum_{j\ge1}\left\lfloor\frac m{p^j}\right\rfloor
       =\frac{m-s_p(m)}{p-1},
\]
where $s_p(m)$ is its base-$p$ digit sum.  For $m\ge1$,
$s_p(m)\le(p-1)(\log_p m+1)$; $s_p(0)=0$.
Sum this estimate over the coordinates of $A$ to obtain
\eqref{djw:multifactorial}, allowing zero coordinates and $A=0$.
The same formula gives
\[
 v(C_r(m))=
 \frac{rm+\sum_{j<r}s_p(j)-\sum_{j<r}s_p(m+j)}{p-1}.
\]
Discard the last, nonpositive, sum and bound each remaining digit
sum by $(p-1)(\log_p r+1)$.  The upper bound on $v(m!)$ follows
by discarding its nonnegative digit sum.
\end{proof}
% END DJW1

% BEGIN DJW2
\paragraph{DJW2: the determinant factorial and its adjoint identity.}
\label{claim:DJW2}
Let $D=\det X$ with $r\ge1$.  The ordinary Cayley identity is
\begin{equation}
 \det(\partial_X)D^m=
       \prod_{j=0}^{r-1}(m+j)D^{m-1}\qquad(m\ge1).
 \label{djw:cayley}
\end{equation}
Its iteration gives
$\langle D^m,D^m\rangle=C_r(m)$.
If $0\le s\le y$ and $Q$ is homogeneous of degree $rqs$, then
\begin{equation}
 \frac{\langle D^{qy},D^{q(y-s)}Q\rangle}{C_r(qy)}
       =\frac{\langle D^{qs},Q\rangle}{C_r(qs)}.
 \label{djw:adjoint}
\end{equation}

\begin{proof}
We use \eqref{djw:cayley} as the classical polynomial identity in
Caracciolo--Sokal--Sportiello, Theorem 2.1.  It is an identity over
$\mathbb Q$ and hence over every characteristic-zero $K$.
No measure, metric or analytic continuation enters it.

For a single variable, direct comparison of factorials gives
$\langle P,X_jQ\rangle=\langle\partial_jP,Q\rangle$.
Linearity and multiplication give the corresponding identity for
$D$.  Apply it $q(y-s)$ times, and use \eqref{djw:cayley} on the
left polynomial.  Its scalar is $C_r(qy)/C_r(qs)$, which is nonzero
in $K$.  This proves \eqref{djw:adjoint} and, on taking $s=0$,
the stated determinant factorial.
\end{proof}
% END DJW2

% BEGIN DJW3
\paragraph{DJW3: the weighted exponential and the Frobenius error.}
\label{claim:DJW3}
Put $E=\exp U$ and introduce an independent matrix $Z$.
Then
\begin{equation}
 v(E_{A,B})+a|A|+b|B|\ge0,\qquad E(X,0)=1.
 \label{djw:weightedexp}
\end{equation}
For
$\Psi(X,Z)=\det(X^{[q]}-p^eZ)-D^q$, where
$X^{[q]}=(X_{ij}^q)$, every coefficient of a monomial of
$Z$-degree $k$ in $\Psi^s$ is divisible by
\begin{equation}
 p^{ek+(s-k)_+},
 \label{djw:psidiv}
\end{equation}
and its $X$-degree is $q(rs-k)$.

\begin{proof}
A contribution from $U^j/j!$ has weighted valuation at least
$j/(p-1)-v(j!)\ge0$.  At a fixed bidegree only finitely many $j$
occur, because every monomial of $U$ has positive right degree.
The ultrametric inequality proves \eqref{djw:weightedexp}.

Every $Z$ in a determinant expansion supplies a factor $p^e$.
The part of $\Psi$ free of $Z$ is divisible by $p$, since in
characteristic $p$ the Frobenius map commutes with determinants:
$\det(X^{[q]})=(\det X)^q$.
Among $s$ factors of $\Psi$, at most $k$ can have positive
$Z$-degree if the total degree is $k$.  At least $(s-k)_+$ of
them therefore supply an additional $p$.  This proves
\eqref{djw:psidiv} before and after equal monomials are combined.
Every factor has weighted degree $rq$ if an $X$ has weight one
and a $Z$ has weight $q$; the asserted $X$-degree follows.
\end{proof}
% END DJW3

% BEGIN DJW4
\paragraph{DJW4: an exact Newton expansion of the normalized jets.}
\label{claim:DJW4}
For $s\ge0$, let $Q_s(X)$ be the homogeneous part of degree $rqs$
of
$[\Psi(X,\partial_Y)^sE(X,Y)]_{Y=0}$.
Then
\begin{equation}
 \frac{\Lambda_y}{c^{ry}C_r(qy)}
   =\sum_{s=0}^{y}\binom ys a_s,\qquad
 a_s=\frac{\langle D^{qs},Q_s\rangle}{C_r(qs)},\qquad a_0=1.
 \label{djw:newton}
\end{equation}

\begin{proof}
Conjugating $\partial_{Y_{ij}}$ by
$\exp(c\sum X_{ij}^qY_{ij})$ replaces it by
$\partial_{Y_{ij}}+cX_{ij}^q$.
The entries of $X$ commute with these derivatives, and
$c^{-1}=-p^e$.  Consequently the full right operator becomes
\[
 c^{ry}\bigl(D^q+\Psi(X,\partial_Y)\bigr)^y.
\]
Its two summands commute.  Expand by the binomial theorem and set
$Y=0$.  The left polynomial $D^{qy}$ sees only total $X$-degree
$rqy$, so in the term with index $s$ it sees precisely $Q_s$.
Apply DJW2 to remove $D^{q(y-s)}$ from that term.  This is
\eqref{djw:newton}.  Since $E(X,0)=1$, we have $Q_0=1$ and
$a_0=1$.  Every sum used for a fixed coefficient is finite; no
convergence in the valued field is needed.
\end{proof}
% END DJW4

% BEGIN DJW5
\paragraph{DJW5: an explicit lower bound for every late Newton coefficient.}
\label{claim:DJW5}
For $s\ge1$ the coefficients in DJW4 satisfy
\begin{equation}
 v(a_s)\ge
 \frac s8-r^2\bigl(3\log_p r+2\log_p s+h+3\bigr).
 \label{djw:newtonbound}
\end{equation}

\begin{proof}
Fix a term of $Z$-degree $k$ in $\Psi^s$.
Its $X$-degree is $q(rs-k)$, so its contribution to $Q_s$ takes
a coefficient of $E$ of bidegree $(qk,k)$.
Differentiation in $Y$ contributes a multiindex factorial of
degree $k$.  By DJW1 and DJW3 the resulting coefficient has
valuation at least
\[
 ek+(s-k)_++\frac{k}{p-1}
       -r^2(\log_p(rs)+1)-aqk-bk.
\]
The exact identity
\[
 qa+b=e+\frac1{p-1}-\frac{1+q^{-1}}8
\]
reduces this to
\[
 (s-k)_++\frac{k(1+q^{-1})}{8}
       -r^2(\log_p(rs)+1)
 \ge \frac s8-r^2(\log_p(rs)+1).
\]
For $k\le s$ use $s-k+k/8\ge s/8$; for $k\ge s$ use
$k/8\ge s/8$.  Taking finite sums proves this lower bound for
each coefficient of $Q_s$.

The coefficients of $D^{qs}$ are integers.  Its factorial pairing
with $Q_s$ involves only factorials of total degree $rqs$.
DJW1 gives
\[
 v(\langle D^{qs},Q_s\rangle)
 \ge \frac s8-r^2(\log_p(rs)+1)
       +\frac{rqs}{p-1}-r^2(\log_p(rqs)+1).
\]
Subtract the upper bound \eqref{djw:capacity} with $m=qs$.
Since $\log_pq=h$ and $r\le r^2$, the result is
\eqref{djw:newtonbound}.  A zero coefficient or pairing has
valuation $+\infty$ and causes no exception.
\end{proof}
% END DJW5

% BEGIN DJW6
\paragraph{DJW6: the stated numerical threshold is sufficient.}
\label{claim:DJW6}
Let
\begin{equation}
 S(p,h,r)=\left\lceil64r^2
          \bigl(h+3+3\log_p(r+1)\bigr)\right\rceil
          \qquad(r\ge1).
 \label{djw:windowlength}
\end{equation}
The right side of \eqref{djw:newtonbound} is strictly positive
for every integer $s\ge S(p,h,r)$.

\begin{proof}
Set $A=h+3+3\log_p(r+1)\ge4$ and $s_0=64r^2A$.
The function
\[
 F(t)=\frac{t}{8r^2}-h-3-3\log_p r-2\log_p t
\]
has positive derivative for $t\ge s_0$:
$1/(8r^2)-2/(t\log p)>0$ there.
With $b_0=\log_p(r+1)$, direct substitution gives
\begin{align*}
 F(s_0)&\ge7A-4b_0-2\log_p64-2\log_pA\\
       &>7A-\frac{4A}{3}-6-2A
        =\frac{11A}{3}-6>0.
\end{align*}
We used $b_0\le A/3$, $\log_p64<3$ for $p\ge5$, and
$\log_pA\le A$ for $A\ge1$.  Monotonicity and the ceiling
now prove the assertion.
\end{proof}
% END DJW6

% BEGIN DJW7
\paragraph{DJW7: integral interpolation on every translated interval.}
\label{claim:DJW7}
Let $K$ be any field with an additive nonarchimedean valuation,
$S\ge1$, and $a_s\in K$ with $a_0=1$ and $v(a_s)>0$ for
every $s\ge S$.  For each $Y_0\ge0$, some
$Y_0\le y<Y_0+S$ satisfies
\[
 v\left(\sum_{s=0}^{y}\binom ys a_s\right)\le0.
\]
This assertion does not require characteristic zero or a normalized
prime.

\begin{proof}
The integer matrix
$V_{ij}=\binom{Y_0+i}{j}$, $0\le i,j<S$, factors as $BC$,
where $B_{ik}=\binom ik$ is lower triangular with diagonal one
and $C_{kj}=\binom{Y_0}{j-k}$ is upper triangular with diagonal
one.  The factorization is the elementary binomial convolution.
Thus $V^{-1}$ has integer entries, over $\mathbb Z$ and after
mapping to $K$.  Every integer has nonnegative valuation.

If all $S$ values in question had positive valuation, their finite
tails with $s\ge S$ would also have positive valuation.  Subtract
those tails, treating $\binom ys$ as zero when $s>y$.
The resulting vector of truncated sums is
$V(a_0,\ldots,a_{S-1})^t$ and has every entry of positive valuation.
Multiplication by $V^{-1}$ preserves this property, which would
give $v(a_0)>0$.  This contradicts $v(1)=0$.
\end{proof}
% END DJW7

% BEGIN DJW8
\paragraph{DJW8: the full original short-window theorem.}
\label{claim:DJW8}
Under \eqref{djw:parameters}--\eqref{djw:hm2}, if $r\ge1$,
each $Y_0\ge0$ admits $Y_0\le y<Y_0+S(p,h,r)$ with
\begin{equation}
 v(\Lambda_y)\le
 ry\bigl((q+1)\gamma-1\bigr)+n(\log_p n+1).
 \label{djw:thmp}
\end{equation}
If $r=0$, then $\Lambda_y=1$ for every $y\ge0$.
This is the complete original
\texttt{Formal\allowbreak Series\allowbreak Bounds.\allowbreak Thm\allowbreak P}, including its zero-size conjunct.

\begin{proof}
For $r\ge1$, use DJW4--DJW7 to find a normalized value with
valuation at most zero.  DJW1 then gives
\begin{align*}
 v(\Lambda_y)
 &\le-ery+\frac{rqy}{p-1}+r(\log_pr+1)\\
 &=ry\bigl((q+1)\gamma-1\bigr)+r(\log_pr+1)\\
 &\le ry\bigl((q+1)\gamma-1\bigr)+n(\log_pn+1).
\end{align*}
The middle identity is $q/(p-1)-e=(q+1)p/(p^2-1)-1$.
The last inequality follows from $1\le r\le n$ and monotonicity
of the two nonnegative factors.

For $r=0$ there are no variables.  The only monomial has right
degree zero, so \eqref{djw:hm2} forces $U=0$.  Both determinant
operators are identities and the exponential is one.  In particular
we do not try to use a zero-length interval in this case.
\end{proof}
% END DJW8

\subsection{The exact row identity}

The next algebraic statements have a wider scope than DJW8.
In this subsection $K$ is any characteristic-zero field, and the
integers $l,i,n,r,y$ are arbitrary nonnegative integers with $r\le n$.
There is no prime, valuation, oddness or separated-row hypothesis.
For these general indices use the original integer
\[
 e_{l,i}=1+l\sum_{j<(i-1)/2}(l^2)^j,
 \qquad q=l^i,\qquad c=-\bigl(l^{e_{l,i}}\bigr)^{-1},
\]
where the sum uses integer truncation and the field inverse is zero
at zero.  This agrees with \eqref{djw:parameters} in the prime,
odd-distance case.  Let $Z,W$ be independent $n$ by $n$ matrices.
Restriction $\operatorname{res}_r$ sets every variable outside the
leading $r$ by $r$ blocks to zero.

For completeness the source defines a truncated coefficient operation
on every series $H$, even one with nonzero constant term, by
\[
 \operatorname{E}(H)_A
       =\sum_{j=0}^{|A|}[X^A]H^j/j!.
\]
When $H(0)=0$ this is the ordinary formal exponential.  We use
the symbol $\operatorname{E}$ only to state its unrestricted
restriction law precisely.

% BEGIN DJW9
\paragraph{DJW9: restriction, support and genuine formal exponentials.}
\label{claim:DJW9}
Restriction is a ring homomorphism and
\[
 \operatorname{res}_r\operatorname{E}(H)
       =\operatorname{E}(\operatorname{res}_rH)
\]
for every $H$.  If $H(0)=H'(0)=0$, then
$\exp(H+H')=\exp(H)\exp(H')$.
Series supported on a specified subset of variables remain so under
sum, product, powers and $\operatorname{E}$.  A substitution by
series with zero constant term preserves order of vanishing and
commutes with ordinary formal exponentiation of a zero-constant
series.

\begin{proof}
The coefficient of a product is the finite sum over decompositions
of its exponent.  An exponent supported in the retained variables
can be decomposed only into exponents supported there.  This proves
multiplicativity.  The retained exponent has the same degree in
the small and large variable sets, so the upper limit $|A|$ in
$\operatorname{E}$ is unchanged.  Applying multiplicativity to
each $H^j$ proves the unrestricted restriction law.

For zero-constant series, the coefficient of degree at most $N$
in a power with exponent $>N$ is zero.  The usual binomial identity,
divided by its factorial, therefore proves the exponential product
identity by finite sums in every degree.  The assertions about
support follow from the same coefficient decompositions.

Under a zero-constant substitution a monomial of degree $d$ has
no terms of degree less than $d$.  Hence to compute a coefficient
of degree $N$ one may first truncate the substituted series and
the exponential at degree $N$.  Both are then finite polynomial
identities.  This proves the last two assertions without any
topological hypothesis on $K$.
\end{proof}
% END DJW9

% BEGIN DJW10
\paragraph{DJW10: determinant coefficients with prescribed diagonal rows.}
\label{claim:DJW10}
Every monomial of $(\det Z)^A(\det W)^B$ has degree $A$ in
each row of $Z$ and $B$ in each row of $W$.
If in every row $\rho\ge r$ a monomial uses only the diagonal
entries, with exponents $A$ and $B$, then its coefficient is the
coefficient of its remaining exponent in
$(\det Z_{<r,<r})^A(\det W_{<r,<r})^B$.
In particular this coefficient vanishes if the remaining monomial
uses a column $\ge r$.

\begin{proof}
Each permutation monomial of a determinant uses exactly one entry
in each row and column.  Taking products proves the degree assertion.
Set all off-diagonal entries in rows $\rho\ge r$ equal to zero.
The resulting matrix is block upper triangular, with its lower
diagonal entries on the diagonal of the lower block.  Its determinant
is the leading-block determinant times their product, with sign $+1$.
Raise to the stated powers and compare coefficients.
Specialization kills precisely the monomials using a deleted variable,
so it does not change any coefficient under consideration.  This
argument proves the source's full \texttt{Det\allowbreak Pow\allowbreak Diag} over every
commutative coefficient ring, including $r=0$, $r=n$ and zero powers.
\end{proof}
% END DJW10

% BEGIN DJW11
\paragraph{DJW11: the full diagonal-block jet formula.}
\label{claim:DJW11}
Suppose $G$ uses only rows $\rho<r$ and $H$ only rows $\rho\ge r$.
Assume that, among exponents having degree at most $A$ in each
right-hand $Z$ row and at most $B$ in each right-hand $W$ row,
the only nonzero coefficient of $H$ is
\[
 \prod_{\rho=r}^{n-1}Z_{\rho\rho}^{A}W_{\rho\rho}^{B},
 \qquad\text{with coefficient }1.
\]
Then
\begin{equation}
 \langle(\det Z)^A(\det W)^B,GH\rangle
  =(A!B!)^{n-r}
    \langle(\det Z_{<r,<r})^A(\det W_{<r,<r})^B,
                \operatorname{res}_rG\rangle.
 \label{djw:blockjet}
\end{equation}
In particular this proves the full source \texttt{Block\allowbreak Jet} on taking
$A=Ly$, $B=y$ for arbitrary $L,y\ge0$.

\begin{proof}
For a monomial in the differential polynomial, DJW10 fixes each
row degree.  In the coefficient convolution of $GH$, the exponent
on the last rows comes entirely from $H$.  It is therefore within
the stated row-degree limits.  The hypothesis on $H$ forces exactly
the displayed diagonal exponent, with coefficient one.
DJW10 then identifies the polynomial coefficient on the first rows
with its leading-block coefficient and discards any remaining extra
columns.  The factorial of the complete exponent is the factorial
of its first-block exponent times $(A!B!)^{n-r}$, since the supports
are disjoint.  Summing gives \eqref{djw:blockjet}.
This coefficient proof also covers empty blocks and $A=0$ or $B=0$;
no differentiation by a positive number of variables was presumed.
\end{proof}
% END DJW11

Let $H_\rho,T_\rho,T'_\rho\in K[[Z,W]]$ satisfy the five
original \texttt{Bi\allowbreak Hyp} requirements:
each $H_\rho$ uses only row $\rho$ and has zero constant term;
for $\rho\ge r$, $T_\rho$ uses only that row's $Z$ variables
and equals $Z_{\rho\rho}$ to degree one, while $T'_\rho$
uses only that row's $W$ variables and equals $W_{\rho\rho}$
to degree one; finally
\begin{equation}
 \operatorname{res}_r\sum_{\rho<r}H_\rho
       =c\sum_{a,b<r}Z_{ab}^{q}W_{ab}+U.
 \label{djw:biR3}
\end{equation}
Here ``equals to degree one'' includes equality of the constant
and every linear coefficient, not just the indicated diagonal
coefficient.  Define the literal row product
\[
 F_y=\prod_{\rho=0}^{n-1}
       \left(\begin{cases}1,&\rho<r,\\
           T_\rho^{qy}(T'_\rho)^y,&\rho\ge r
       \end{cases}\right)\exp(H_\rho).
\]

% BEGIN DJW12
\paragraph{DJW12: the literal row-value identity \texttt{Bi\allowbreak Value}.}
\label{claim:DJW12}
For every $y\ge0$ and all the algebraic indices just specified,
\begin{equation}
 \langle(\det Z)^{qy}(\det W)^y,F_y\rangle
       =((qy)!y!)^{n-r}\Lambda_y.
 \label{djw:bivalue}
\end{equation}
The $\Lambda_y$ on the right is the original small-block series
formed from the same $U$ in \eqref{djw:biR3}; neither of the
valuation hypotheses \eqref{djw:hm1}--\eqref{djw:hm2} is needed
for this identity.

\begin{proof}
Factor $F_y=GH$, where
$G=\exp(\sum_{\rho<r}H_\rho)$ and $H$ is the product over
the remaining rows.  DJW9 justifies this factorization and its
row supports.  A nonzero monomial in $T_\rho^{qy}$ has
$Z$-row degree at least $qy$, with equality only for
$Z_{\rho\rho}^{qy}$ and coefficient one.  This follows by choosing
the linear term from every factor; every other term has larger degree.
The same statement holds for $(T'_\rho)^y$.
Every monomial of $\exp(H_\rho)$ has nonnegative degrees on both
sides, and its constant coefficient is one.  Thus, within the degree
limits $qy,y$, the product on the last rows has exactly the diagonal
coefficient required by DJW11.  A nonconstant exponential term cannot
survive these limits.  This remains valid if either power is zero.

Apply DJW11 and then DJW9.  The restricted first-row factor is
\[
 \operatorname{res}_rG
   =\exp\left(\operatorname{res}_r\sum_{\rho<r}H_\rho\right)
   =\exp\left(c\sum_{a,b<r}Z_{ab}^{q}W_{ab}+U\right).
\]
Its exponent has zero constant term because the left side of
\eqref{djw:biR3} does.  The remaining jet is precisely $\Lambda_y$.
This proves every quantifier of \texttt{Bi\allowbreak Value}, including $r=0$,
$n=0$, $l=0$ and $i=0$ when their algebraic data are supplied.
\end{proof}
% END DJW12

Return now to the valued prime and odd-distance hypotheses of the
first subsection.  For brevity write
\[
 B_r(y)=ry((q+1)\gamma-1)+n(\log_p n+1).
\]

% BEGIN DJW13
\paragraph{DJW13: the exact factorial content loss.}
\label{claim:DJW13}
If $v(\Lambda)\le B_r(y)$, then
\begin{equation}
 v\bigl(((qy)!y!)^{n-r}\Lambda\bigr)
       \le B_n(y)+\frac{(n-r)(q+1)y}{p-1}.
 \label{djw:content}
\end{equation}
This is the full original \texttt{Bi\allowbreak Content} for every scalar
$\Lambda$ satisfying its bound.

\begin{proof}
DJW1 bounds the two factorial valuations by $qy/(p-1)$ and
$y/(p-1)$.  Also
\[
 (q+1)\gamma-1\ge0,
\]
because $q\ge p$ and
$(p+1)p/(p^2-1)>1$.  Hence $B_r(y)\le B_n(y)$.
Add these inequalities using multiplicativity of the valuation.
The assumed finite upper bound also excludes $\Lambda=0$;
there is no subtraction of an infinite valuation.
\end{proof}
% END DJW13

% BEGIN DJW14
\paragraph{DJW14: \texttt{Bi\allowbreak Window} and \texttt{content\_\allowbreak and\_\allowbreak window}.}
\label{claim:DJW14}
Suppose $1\le r\le n$, $n-r\le(h+1)m$, and the literal
\texttt{Bi\allowbreak Hyp} data satisfy \eqref{djw:hm1}--\eqref{djw:hm2}
for their same $U$.  Every $Y_0\ge0$ has an index
$Y_0\le y<Y_0+S(p,h,r)$ such that
\begin{equation}
 v\bigl(\langle(\det Z)^{qy}(\det W)^y,F_y\rangle\bigr)
 \le B_n(y)+\frac{(h+1)m(q+1)y}{p-1}.
 \label{djw:biwindow}
\end{equation}
Together with DJW13, this supplies the full original
\texttt{content\_\allowbreak and\_\allowbreak window} theorem.

\begin{proof}
Choose $y$ by DJW8 for the small block.  Substitute the exact
identity DJW12 and apply DJW13.  The final comparison replaces
$n-r$ by $(h+1)m$ in a nonnegative expression.  This is exactly
the source's \texttt{bi\allowbreak Window\_\allowbreak of} implication with both its
inputs now proved; it adds no geometric premise.
\end{proof}
% END DJW14

\subsection{The perturbation attached to the original normal form}

We now fix the original datum
$\mathcal X=\texttt{Row\allowbreak Data}(K,V,d,I,J)$.
Its field is algebraically closed of characteristic zero, $d>0$,
and $V=(v,p)$ is the original valued-prime datum.
There are two finite sets of tangent labels $I_0,I_1$, their
Cartier data $G_0,G_1$, and the original leaf gauges
\[
 H_{g,\alpha,\alpha'}\in
 K[[z_w:w\in I_0\amalg I_1]].
\]
The datum contains the assigned types, the ranks, integral pairings,
Cartier identities, and comparison of integral-germ spaces between
coordinate systems.  None of those fields is silently changed here.
We use in particular its axis-vanishing property: a gauge coefficient
supported on just one of the two sides is zero.  Thus each such
gauge has zero constant term.

The actual integers in this datum are denoted $n,j,j'$ and
$m=j+j'$.  Fix $g\in\{0,1\}$, a type $\alpha\in\mathbb Z/d\mathbb Z$,
and $h\ge0$.  Until DJW21, these algebraic constructions do not
require $h$ odd or $p\ge5$.
Put $q=p^h$ and use the original integer $e_{p,h}$ above.
For row vectors $u_c$ on $I_g$ and $v_b$ on $I_{1-g}$ write
\[
 H(u,v)(x,y)=H_{g,\alpha,\alpha+h}
       \left(\sum_cu_c x_c,\sum_bv_b y_b\right).
\]
This is the literal linear pullback of the given gauge, not a
replacement series chosen for the estimate.

Here is the exact normal-form premise used below.  The given arrays
$U_0,U'_0$ each have $n$ rows.  An original
\texttt{MConcl} consists of $0\le r\le n$, vectors $e_c$ for
$c<r$, vectors $v'_b$ for $b<n$, and a matrix $\Phi$ of size
$n$ by $r$, with the following properties:
\begin{enumerate}
\item $n\le r+(h+1)j'+j$.
\item The $e_c$ are integral and have an integral right inverse:
there is $w_{x,c}$ with
$\sum_x e_{c,x}w_{x,c'}=\delta_{cc'}$.
\item $e=A U_0$ for an integral $r$ by $n$ matrix $A$.
\item $v'=Y U'_0$ for an integral $n$ by $n$ matrix $Y$ with
$v(\det Y)=0$.
\item $\Phi$ is integral and its first $r$ rows are $-I_r$.
\item Every coefficient $R_{A,B}$ of
\begin{equation}
 R(x,y)=H(e,v')(x,y)
        -p^{-e_{p,h}}\sum_{b<n,c<r}\Phi_{bc}x_c^q y_b
 \label{djw:actualremainder}
\end{equation}
satisfies
$v(R_{A,B})+a|A|+b|B|\ge1/(p-1)$, with weights
\eqref{djw:weights} for the present $p,h$.
\item $R_{A,0}=0$ for every $A$.
\end{enumerate}
These are exactly the original fields, with its first pair of
integer inequalities separated for readability.  When $h=1$ the
additional original \texttt{MPlus\allowbreak Concl} means
$v(R_{A,B})\ge0$ whenever $|A|<p^2$ and $|B|<p$.
Saturation here is the displayed integral right inverse; no
Noetherian or discrete-valuation interpretation is substituted.

% BEGIN DJW15
\paragraph{DJW15: all four original perturbation conclusions.}
\label{claim:DJW15}
Restrict \eqref{djw:actualremainder} by setting $y_b=0$ for
$b\ge r$ and denote the result by $R_0(x,y)$.
For independent $r$ by $r$ matrices $Z,W$ put
\begin{equation}
 U(Z,W)=\sum_{\rho<r}R_0(Z_{\rho,*},W_{\rho,*}).
 \label{djw:actualU}
\end{equation}
This is exactly
\texttt{restricted\allowbreak Curve\allowbreak System\allowbreak Series} of the original normal form.
It satisfies all of \texttt{Mto\allowbreak HM}:
\begin{align}
 \sum_{\rho<r}H(e,v'_{<r})(Z_{\rho,*},W_{\rho,*})
   &=-p^{-e_{p,h}}\sum_{\rho,c<r}Z_{\rho c}^qW_{\rho c}+U,
        \label{djw:actualmain}\\
 v(U_{A,B})+a|A|+b|B|&\ge\frac1{p-1},\qquad U_{A,0}=0.
        \label{djw:actualHM}
\end{align}
If $h=1$ and \texttt{MPlus\allowbreak Concl} holds, the decomposition
\eqref{djw:actualU} has each summand supported in one row and
integral in all bidegrees $|A|<p^2$, $|B|<p$.

\begin{proof}
The restriction is the substitution $y_b\mapsto0$ for $b\ge r$.
It commutes with the actual linear pullback of the gauge by DJW9's
coefficient argument.  Since the first $r$ rows of $\Phi$ are
$-I_r$, it takes the explicit main term in
\eqref{djw:actualremainder} to
$-p^{-e_{p,h}}\sum_{c<r}x_c^qy_c$.
Thus the restricted remainder is precisely the difference defining
the original restricted series.  Substitute each row and sum to
obtain \eqref{djw:actualmain}, including its negative sign.

A coefficient of $R_0$ is exactly a coefficient of $R$ with zeros
inserted in the deleted positions.  Its left and right total degrees
are unchanged.  Hence it retains the weighted inequality and
the first-axis vanishing of $R$.  Row substitution is an injective
renaming of variables.  A coefficient in its image retains those
same two degrees; a coefficient outside its image is zero.
For a fixed matrix exponent, every summand in
\eqref{djw:actualU} therefore has valuation at least the same
number $1/(p-1)-a|A|-b|B|$.  The valuation inequality for finite
sums proves \eqref{djw:actualHM}, even where several summands
contribute or cancel.  First-axis coefficients vanish in every
summand, proving the second assertion without relying on such
cancellation.

When $h=1$, use the additional coefficient bound on $R$ before
restriction.  The two degree cutoffs are again unchanged by
restriction and injective row renaming.  This proves the exact
\texttt{HSeparated\allowbreak Row\allowbreak Sum} conjunct, with the displayed summands
as its witnesses.  If $r=0$, the sum is empty and every assertion
has the same literal meaning.  We have imposed neither the
alternating-type condition nor oddness in proving this conversion.
\end{proof}
% END DJW15

\subsection{Constructing the actual integral row tests}

Write $\mathcal O=\{a\in K:v(a)\ge0\}$.
The integral germs of a coordinate system $t=(t^g_x)$ are, by
the original definition, the series
\begin{equation}
 \gamma\bigl((t^g_x)_{g,x}\bigr),\qquad
       \gamma\in\mathcal O[[z_w:w\in I_0\amalg I_1]],
 \label{djw:intgerm}
\end{equation}
where each $t^g_x$ is inserted into its own side of the disjoint
variable set.  This definition does not assert that the coefficients
of $t^g_x$ in logarithmic coordinates are themselves integral.

% BEGIN DJW16
\paragraph{DJW16: completing split integral rows without a discrete valuation.}
\label{claim:DJW16}
Suppose $A\in M_{r,n}(\mathcal O)$,
$B\in M_{n,r}(\mathcal O)$ and $AB=I_r$.
For every $r\le j\le n$ there are integral matrices
$A_j\in M_{j,n}(\mathcal O)$ and $B_j\in M_{n,j}(\mathcal O)$
with $A_jB_j=I_j$ and first $r$ rows of $A_j$ equal to $A$.
For $j=n$ this gives inverse square matrices
$\bar A,\bar B\in M_n(\mathcal O)$ and
$v(\det\bar A)=v(\det\bar B)=0$.

\begin{proof}
It suffices to add one row to a split pair with $j<n$.
Over $K$, the kernel of $A_j:K^n\to K^j$ contains a nonzero
vector $u$.  Choose an index $k$ at which $v(u_k)$ is minimal.
The finite minimum is finite, so $u_k\ne0$.
Set $s_x=u_x/u_k$.  Every $s_x$ is integral, $s_k=1$, and
$A_js=0$.  Define the new row
\[
 r_x=\delta_{kx}-\sum_{c<j}(B_j)_{kc}(A_j)_{cx}.
\]
It is integral.  Direct multiplication gives
\[
 rB_j=0,\qquad rs=1,\qquad A_js=0.
\]
Append $r$ to $A_j$ and append $s$ as a column to $B_j$.
These three identities together with $A_jB_j=I_j$ prove that
the enlarged matrices still multiply to the identity.
Induction supplies every intermediate size, including the empty
initial set of rows.

For $j=n$, taking determinants gives
$\det\bar A\det\bar B=1$.
Both determinants are integral, so both have valuation zero;
the matrices are inverses over $K$ and already have integral
entries.  The proof used only a minimum in a finite list of
valuations.  No uniformizer, completeness or residue-field
finiteness is needed.  This proves the original
\texttt{complete\_\allowbreak rows}, with all its intermediate sizes.
\end{proof}
% END DJW16

% BEGIN DJW17
\paragraph{DJW17: actual standard Cartier coordinates.}
\label{claim:DJW17}
For either original Cartier datum $G_g$ one can choose its system
$S_g$ and its associated formal coordinates $t^g_x$ so that
\[
 (S_g).\mathrm{IsT}(t^g),\qquad
 t^g_x=z_x+\text{terms of total degree at least two}.
\]
These choices depend on the datum, not on the later integer $y$.

\begin{proof}
The original Cartier axiom supplies a typed curve with any integral
typed initial tangent vector.  Apply it to each standard basis vector.
The resulting system has identity tangent matrix, hence unit
determinant.  If $\alpha_{c,s,x}$ is its tangent-side coefficient
at Cartier index $s$, then $\alpha_{c,0,x}=\delta_{cx}$.
The coordinate equations are the formal equalities
\begin{equation}
 \sum_{s\ge0,c}\alpha_{c,s,x}(t_c)^{p^s}=z_x.
 \label{djw:cartiercoord}
\end{equation}
For any fixed degree the sum is finite, since $p\ge2$ and the
coordinates have zero constant term.  Separate the $s=0$ terms
and write this as $t_x=z_x-Q_x(t)$, where every monomial of
$Q_x$ has degree at least two.

Starting from zero, iterate $t\mapsto z-Q(t)$.
If two zero-constant tuples agree below degree $N$, their images
agree below degree $N+1$.  Indeed each difference
$u^a-v^a$ factors by $u-v$, and its remaining $a-1\ge1$
factors have positive degree.  Thus every coefficient of the
iterates stabilizes after finitely many steps.  The stabilized
tuple solves \eqref{djw:cartiercoord}, and the same argument
proves uniqueness among zero-constant tuples.  Its degree-one
part is $z$.

Taking coefficients in \eqref{djw:cartiercoord} gives precisely
the source's finite equation \texttt{Sys.\allowbreak Is\allowbreak T}; terms with
$p^s$ above the coefficient degree vanish.  This constructs the
actual systems and coordinates rather than merely positing a
formal change of variables.
\end{proof}
% END DJW17

% BEGIN DJW18
\paragraph{DJW18: linear substitution and the ordinary differential pairing.}
\label{claim:DJW18}
Let $L:z\mapsto w$ be a linear substitution between finite sets
of formal variables over a characteristic-zero field.  For every
polynomial $P$ in the $z$ variables and every series $F$ in the
$w$ variables,
\begin{equation}
 \langle P,F\circ L\rangle
       =\langle P(L^t\partial_w),F\rangle.
 \label{djw:linearjet}
\end{equation}
Here the right side means the polynomial obtained by replacing
each differential letter $\partial_{z_j}$ by
$\sum_iL_{ij}\partial_{w_i}$.  If, row by row, a square matrix
of these letters becomes $P_0A^t$, its determinant becomes
$\det A\det P_0$.

\begin{proof}
For one derivative the coefficient formula for substitution gives
the ordinary chain rule
$\partial_{z_j}(F\circ L)
=\sum_iL_{ij}(\partial_{w_i}F)\circ L$.
It is enough to check this on monomials; the factorial coefficients
agree directly.  Iterate the identity and set $z=0$, then extend
linearly over $P$.  In every prescribed degree all calculations
involve finite polynomials, so the argument applies to formal
series.  Determinant multiplicativity gives the last assertion.
In particular the scalar is $\det A$, not its inverse.  This
orientation will fix the exact scalar in the actual row value.
\end{proof}
% END DJW18

For the rest of the construction take injective maps
$\iota:[n]\hookrightarrow I_g$ and
$\iota':[n]\hookrightarrow I_{1-g}$, and specialize the
given normal form to
\[
 (U_0)_{b,x}=\delta_{x,\iota(b)},\qquad
 (U'_0)_{b,x}=\delta_{x,\iota'(b)}.
\]
These are the same selected letters as in the original
\texttt{block\_\allowbreak of\_\allowbreak normal\_\allowbreak form} theorem.  No further
type or alternating hypothesis is needed once this original
normal form is supplied.

% BEGIN DJW19
\paragraph{DJW19: the actual row substitution and all five \texttt{Bi\allowbreak Hyp} fields.}
\label{claim:DJW19}
For the given normal form, there are systems $S_g$, coordinates
$t^g$, row exponents $H_\rho$ and row coordinates
$T_\rho,T'_\rho$ satisfying the original \texttt{Bi\allowbreak Hyp} with
the exact perturbation \eqref{djw:actualU}.
All these objects are chosen before $y$.

\begin{proof}
Let $w$ be the integral right inverse of $e$ in the normal form.
The matrix $B_{b,c}=w_{\iota(b),c}$ satisfies $AB=I_r$:
substitute $e_{c,x}=\sum_bA_{cb}\delta_{x,\iota(b)}$ in
$\sum_xe_{c,x}w_{x,c'}=\delta_{cc'}$.
Apply DJW16 to obtain square inverse matrices $\bar A,\bar B$
extending the first $r$ rows of $A$.
Also put $Z_0=Y^{-1}$.  It is integral because the adjugate of
the integral matrix $Y$ is integral and $v(\det Y)=0$.
Define
\[
 E_{c,x}=\sum_b\bar A_{cb}\delta_{x,\iota(b)}.
\]
Its first $r$ rows are $e$.  Moreover
$E_{c,\iota(b)}=\bar A_{cb}$ and
$v'_{c,\iota'(b)}=Y_{cb}$, by injectivity of the selected letters.
Choose the standard systems and coordinates from DJW17.

For each $\rho<n$, the actual row substitution is
\begin{align*}
 z_{(g,x)}&\longmapsto\sum_cE_{c,x}Z_{\rho c},\\
 z_{(1-g,x)}&\longmapsto\sum_cv'_{c,x}W_{\rho c}.
\end{align*}
Call it $L_\rho$, and define
\begin{align*}
 H_\rho&=H_{g,\alpha,\alpha+h}\circ L_\rho,\\
 T_\rho&=
 \left(\sum_b\bar B_{b\rho}t^g_{\iota(b)}\right)
       \left(z_x=\sum_cE_{c,x}Z_{\rho c}\right),\\
 T'_\rho&=
 \left(\sum_b(Z_0)_{b\rho}t^{1-g}_{\iota'(b)}\right)
       \left(z_x=\sum_cv'_{c,x}W_{\rho c}\right).
\end{align*}
Every substituted variable has zero constant term, so these formal
substitutions are defined.

The first series uses only row $\rho$ and has constant term zero,
by the axis-vanishing of the original gauge.  The second uses only
the left side of that row.  Its degree-one part is
\[
 \sum_{b,c}\bar B_{b\rho}E_{c,\iota(b)}Z_{\rho c}
       =\sum_c(\bar A\bar B)_{c\rho}Z_{\rho c}
       =Z_{\rho\rho}.
\]
Every higher term of $t$ remains of degree at least two under
linear substitution.  The analogous computation with $YZ_0=I$
gives the precise right diagonal coordinate.

Finally restrict $\sum_{\rho<r}H_\rho$ to the leading blocks.
The first $r$ rows of $E$ are $e$, and the surviving right
columns are exactly $v'_{<r}$.  Thus the restricted sum is the
left side of \eqref{djw:actualmain}.  DJW15 identifies it with
the required main term plus the very same $U$.
This proves all five \texttt{Bi\allowbreak Hyp} fields, including the exact
restriction equality.  No choice in this construction depends
on a derivative order or a window origin.
\end{proof}
% END DJW19

To state the final value identity, use independent ambient variables
$z_{\rho,w}$ for $\rho<n$ and $w\in I_0\amalg I_1$.
Let
\[
 D_g=\det(z_{\rho,(g,\iota(c))})_{\rho,c<n},\qquad
 D_{1-g}=\det(z_{\rho,(1-g,\iota'(c))})_{\rho,c<n},
 \qquad P_y=D_g^{qy}D_{1-g}^{y}.
\]
These are the original selected-letter determinants.  The actual
row value for germs $f_\rho$ is
\begin{equation}
 \mathcal V_y(f)=
 \left\langle P_y,
  \prod_{\rho<n}f_\rho(z_{\rho,*})
              \exp H_{g,\alpha,\alpha+h}(z_{\rho,*})\right\rangle.
 \label{djw:actualvalue}
\end{equation}
Thus its row-kind is the same
$\operatorname{some}(g,\alpha,\alpha+h)$ in every row.

% BEGIN DJW20
\paragraph{DJW20: integral germs and the exact original value comparison.}
\label{claim:DJW20}
For every $y\ge0$, there are actual germs
$f_{\rho,y}\in\operatorname{IntGerm}(V,t)$ such that
\begin{equation}
 \langle(\det Z)^{qy}(\det W)^y,F_y\rangle
   = (\det\bar A)^{qy}(\det Y)^y\,
                     \mathcal V_y(f_{*,y}).
 \label{djw:exactscalar}
\end{equation}
The scalar has valuation zero.  In particular
\begin{equation}
 v(\mathcal V_y(f_{*,y}))
 \le v\bigl(\langle(\det Z)^{qy}(\det W)^y,F_y\rangle\bigr),
 \label{djw:actualcomparison}
\end{equation}
with equality for the displayed witnesses.  DJW19 and this statement
are the full original \texttt{block\_\allowbreak of\_\allowbreak normal\_\allowbreak form} conclusion
for its supplied normal form and restriction identity.
More generally, for every supplied $h_1$ whose main-term expression
equals the same restricted gauge, the construction satisfies
\texttt{Bi\allowbreak Hyp} and \eqref{djw:exactscalar} with that $h_1$;
no weighted estimate on $h_1$ is needed for this statement.

\begin{proof}
Define integral polynomials on the original tangent-label variables by
\[
 \gamma_{\rho,y}(z)=
 \begin{cases}
 1,&\rho<r,\\[2pt]
 \displaystyle
 \left(\sum_b\bar B_{b\rho}z_{(g,\iota(b))}\right)^{qy}
 \left(\sum_b(Z_0)_{b\rho}z_{(1-g,\iota'(b))}\right)^y,
      &\rho\ge r.
 \end{cases}
\]
All their coefficients lie in $\mathcal O$: their linear
coefficients do, and this ring is closed under finite sums and
products.  Put $f_{\rho,y}=\gamma_{\rho,y}(t)$ with the same
coordinate system chosen in DJW19.  Formula \eqref{djw:intgerm}
proves their required integrality in the original sense.

Apply the row substitutions $L_\rho$ simultaneously to the series
inside \eqref{djw:actualvalue}.  By construction each
$f_{\rho,y}$ becomes the corresponding coordinate powers in $F_y$,
and its gauge becomes $H_\rho$.  DJW9 permits composition of the
zero-constant substitutions and commutation with the exponential.
Thus the substituted product is exactly $F_y$.

Apply DJW18 to its derivative polynomial.
The matrix of left derivative letters is the selected-letter
matrix multiplied on the right by $\bar A^t$, and the right
one is multiplied on the right by $Y^t$.  Their determinants
therefore contribute $\det\bar A$ and $\det Y$, respectively.
Taking the powers $qy$ and $y$ proves \eqref{djw:exactscalar}.
Both determinants have valuation zero, by DJW16 and the normal-form
premise.  Taking valuations gives the equality and hence
\eqref{djw:actualcomparison}, even when the value is zero.

For the literal source theorem, the supplied series $h_1$ need
only satisfy its displayed restriction identity, rather than be
named by \eqref{djw:actualU}.  The construction just given uses
that identity in DJW19 and no other property of $h_1$.
Consequently it proves the source's full statement for every such
$h_1$, not only for our canonical choice.  The canonical choice
is used in DJW21 to obtain the weighted estimates from DJW15.
\end{proof}
% END DJW20

% BEGIN DJW21
\paragraph{DJW21: a short window on the actual theta-row data.}
\label{claim:DJW21}
Assume now $p\ge5$, $h$ odd and $n\ge1$, and keep the original
normal form and selected letters of DJW19.
There is one choice of the original systems $S$ and coordinates
$t$, before any $y$ or $Y_0$, with the following property.
Zero is an admissible order, and every $Y_0\ge0$ has
$Y_0\le y<Y_0+S(p,h,n)$ and integral germs $f_{\rho,y}$ with
\begin{equation}
 v(\mathcal V_y(f_{*,y}))\le
 \begin{cases}
 0,&y=0,\\[2pt]
 \displaystyle
 ny((q+1)\gamma-1)+n(\log_p n+1)
       +\frac{(h+1)m(q+1)y}{p-1},&y>0.
 \end{cases}
 \label{djw:actualwindow}
\end{equation}
This is the original \texttt{Windows} conclusion on the actual
selected-letter polynomial and its exact \texttt{block\allowbreak Bd},
conditional only on the stated original normal-form premise.

\begin{proof}
Use DJW15 to supply both weighted hypotheses for the exact $U$,
and DJW19--DJW20 for the actual coordinate and germ comparison.
The normal-form rank inequality gives
\[
 n-r\le(h+1)j'+j\le(h+1)(j+j')=(h+1)m.
\]
If $r\ge1$, DJW14 supplies a suitable $y$ in a window of length
$S(p,h,r)$.  This is at most $S(p,h,n)$, since $r^2$ and
$h+3+3\log_p(r+1)$ are nonnegative increasing functions of $r$.
Combine its estimate with \eqref{djw:actualcomparison}.

If $r=0$, take $y=Y_0$ itself.  DJW8 gives $\Lambda_y=1$,
and DJW12--DJW13 give the same bound after the factorial loss.
The required interval has positive length because $n\ge1$.
The hypothesis on the scalar in DJW13 holds since
$B_0(y)=n(\log_p n+1)\ge0$.

In both cases the case $y=0$ deserves its exact value rather than
the rough logarithmic bound.  The polynomial $P_0$ is one, each
$\gamma_{\rho,0}$ and $f_{\rho,0}$ is one, and every gauge
has zero constant term.  Formula \eqref{djw:actualvalue} is
therefore $\mathcal V_0=1$, of valuation zero.  The same argument
shows explicitly that order zero belongs to the admissible set,
as required by \texttt{Windows}.  All systems and coordinates
were fixed before the order was selected.  This proves the exact
window and its original quantifier order.
\end{proof}
% END DJW21

\subsection{The next join and the retained normal-form construction}

The original \texttt{window\_\allowbreak letter\allowbreak Det} theorem first obtains the
normal form from \texttt{Row\allowbreak Data.\allowbreak Lemma\allowbreak M}.  The theorem proved here
begins with that exact normal form and constructs every remaining
object explicitly.  Thus, once that original existence theorem is
supplied, DJW21 applies to the same selected letters and proves its
window.  In the alternating setting, an odd shift changes which
side is big.  Hence both selected type spaces have rank at least
$n$, and the first $n$ letters of either original flag give the
injective selections used above.  Their indicator rows are saturated:
the transpose indicators are an integral right inverse.  Renaming the
leading minors as these selected-letter determinants gives the two
original flag versions of \texttt{Bi\allowbreak Supply}.

This last description is a conditional assembly, not a claim that
the normal form has now been constructed.  Its original producer
\texttt{Row\allowbreak Data\allowbreak Main\allowbreak Term\allowbreak Normal\allowbreak Form\allowbreak And\allowbreak Level\allowbreak One} uses the adapted
Cartier-chain construction, residue rank estimates, a pairing-kernel
estimate, the weighted gauge remainder, and first-axis vanishing.
Those are the exact upstream tasks retained here.  In particular the
arbitrary series used in the first part has been replaced in the
application by the prescribed restricted gauge, but the common global
theta family and its geometric realization have not been chosen by
this chapter.  The previously closed
\texttt{GLOBAL\_\allowbreak LOCAL\_\allowbreak GAIN} is not reopened or counted as a new
whole closure.

\paragraph{Exact source map.}
All paths below are relative to \texttt{src/\allowbreak } at the stated source
pin.  The metadata records the exact line ranges and byte hashes.
\begin{itemize}
\item The formal series, integer exponent and derivative conventions
are in
\texttt{Formalism/\allowbreak Incoming\allowbreak Bounds/\allowbreak Statement\allowbreak Thm\allowbreak P.\allowbreak lean}, lines 8--42.
The targeted theorem is the distinct, precise
\nolinkurl{Formalism.L2.FormalSeriesBounds.ThmP} in
\texttt{Formalism/\allowbreak Local\allowbreak To\allowbreak Global/\allowbreak Character\allowbreak Germs/\allowbreak }
\texttt{Biadditive\allowbreak In\allowbreak Single\allowbreak Variables.\allowbreak lean}, lines 108--114.
Its row interfaces are at lines 179--224.
\item The window calculation and literal content/window assembly are
in \texttt{Development/\allowbreak Abelian\allowbreak Variety\allowbreak Theory/\allowbreak Proofs/\allowbreak Unramified\allowbreak Depth/\allowbreak }
\texttt{Sum\allowbreak Eq\allowbreak Of\allowbreak Is\allowbreak Homogeneous.\allowbreak lean}, especially lines 771--954 and
961--1085.  The row identity is proved in
\texttt{Block\allowbreak Value\allowbreak On\allowbreak Rows\allowbreak Eq.\allowbreak lean}, lines 477--596, using
\texttt{Det\allowbreak Pow\allowbreak Diagonal\allowbreak Rows\allowbreak Block\allowbreak Jet\allowbreak Restrict.\allowbreak lean}.
\item The exact normal-form and restricted-series definitions are in
\texttt{Formalism/\allowbreak Local\allowbreak To\allowbreak Global/\allowbreak Character\allowbreak Germs/\allowbreak }
\texttt{Lattice\allowbreak Inclusion\allowbreak Forms\allowbreak Injective.\allowbreak lean}, lines 83--166.
\texttt{Main\allowbreak Term\allowbreak Implies\allowbreak Row\allowbreak Hypothesis.\allowbreak lean}, lines 145--194, proves
the full conversion reconstructed in DJW15.
\item The integral completion and differential substitution are in
\texttt{Constant\allowbreak Coeff\allowbreak Subst\allowbreak Same.\allowbreak lean}, especially lines 275--350.
The complete actual row construction is
\texttt{Block\allowbreak Of\allowbreak Normal\allowbreak Form.\allowbreak lean}, lines 317--434.
The consumer is \texttt{Odd\allowbreak Level\allowbreak Block\allowbreak Bound\allowbreak Attained.\allowbreak lean},
lines 89--147, followed by its two flag installations.
\end{itemize}

\paragraph{Published identity and attribution.}
The accepted classical input is S. Caracciolo, A. D. Sokal and
A. Sportiello, \emph{Algebraic/combinatorial proofs of Cayley-type
identities for derivatives of determinants and pfaffians},
Advances in Applied Mathematics \textbf{50} (2013), 474--594,
Theorem 2.1; arXiv:1105.6270, version 2, page 7.
The adjoint-functor reader supplied the useful short organization of
the local window calculation.  The original mathematical argument,
definitions and target results are attributed to Alp\"oge and Claude.


\section{Exact section growth from the classical abelian-variety index theorem}
\label{sec:SGV}

\status{Seven detailed claims supply the complete original
\texttt{sec\allowbreak M\_\allowbreak finite\_\allowbreak and\_\allowbreak growth} output.  This is a classical
replacement proof using PBC26--PBC27 and the explicitly stated Milne
inputs below.  It installs the independently reviewed comparison
argument without changing any earlier proof.}

The point of this chapter is to prove the precise numerical input
needed by GSL19 on the original section modules.  We choose one very
ample power of the given line bundle and compute all its positive
section dimensions exactly.  In particular, there is no need for this
consumer to reconstruct the more general theorem about growth of
filtered polynomial images.  That general theorem remains useful in
other settings; it is not claimed here.

\subsection*{The original objects and the explicit classical input}

Fix an algebraically closed field $K$, with no restriction on its
characteristic, and put $S=\operatorname{Spec}K$.  Let $P$ be an
original $\operatorname{PolAV}(S)$, and write
$A=P.A$, $L=P.L$, $g=P.g$ and $a:A\to S$ for the structural map.
The record gives a group scheme whose structural map is smooth, proper
and geometrically connected; it specifies smooth relative dimension
$g$.  The datum $L$ is an actual line bundle with the original
rigidification, symmetry, relative ampleness and $\operatorname{KTriv}$
fields.  We keep this entire record.  The proof below needs its
ampleness, not an additional identification of its translation kernel
with a Picard-scheme kernel.

Let $T=\operatorname{baseOver}(S)$ be the identity object over $S$.
The actual product in the definition of the sections is
\[
 X=(T\otimes A).\mathrm{left},\qquad
 \pi=(\operatorname{snd}(T,A)).\mathrm{left}:X\longrightarrow A,
 \qquad
 f=(\operatorname{fst}(T,A)).\mathrm{left}:X\longrightarrow S.
\]
Put $R_0=\Gamma(T,\mathcal O_T)$ and, for every integer $j$, put
\[
 V_j=\operatorname{SecM}(A,L,j,T).
                                                        \tag{SGV.1}
\]
These are exactly the source's modules, including its scalar action
through the first projection.

We use three stated parts of classical abelian-variety theory from
J.~S.~Milne, \emph{Abelian Varieties}, corrected 2022 version of the
1986 Storrs article,
\url{https://www.jmilne.org/math/xnotes/AVs.pdf}.
For an abelian variety $B$ over a field and an invertible sheaf $M$,
write $K_M$ for its translation stabilizer, with Milne's reduced
closed-subscheme convention.  Proposition~9.1 on printed page~15
implies that an ample $M$ with a nonzero section has
$\dim K_M=0$.  Theorem~13.3(c) on printed page~25 says that this last
condition concentrates the nonzero cohomology in one degree.
Theorem~13.3(b) on the same page gives
\[
 \chi(B,\mathcal O_B(D))=\frac{(D^{\dim B})}{(\dim B)!}
                                                        \tag{SGV.2}
\]
for a Cartier divisor $D$.  These results hold in arbitrary
characteristic.  They are retained as published classical inputs;
their proofs are not being attributed to this chapter.  We do not
use Theorem~13.3(a), the degree formula for a polarization.

The elementary scheme dictionary uses smoothness over a field implying
regularity, regular local rings being domains, and the fiber-dimension
meaning of smooth relative dimension.  References are the Stacks
Project tags \nolinkurl{04QM}, \nolinkurl{00NN} and
\nolinkurl{01V4}; the component and connectedness argument is given
explicitly below.  The two inherited human proofs used are PBC26
(very ample powers on every actual affine test) and PBC27 (finiteness
of the actual sections, including every integer power).

% BEGIN SGV1
\paragraph{SGV1: the actual identity product is an abelian variety of dimension \(g\).}
\label{claim:SGV1}
The specified projection $\pi:X\to A$ is an isomorphism over $S$.
With the transported group law, $X$ is a geometrically integral
abelian variety, and its dimension is the given integer $g=P.g$.

\begin{proof}
The defining square is Cartesian with lower map
$\operatorname{id}_S$.  Its universal property supplies
$\beta:A\to X$ from the pair $(a,\operatorname{id}_A)$.
The second projection gives $\pi\beta=\operatorname{id}_A$.
Both projections of $\beta\pi$ agree with the corresponding
projections of $\operatorname{id}_X$, so the uniqueness part of the
same universal property gives $\beta\pi=\operatorname{id}_X$.
Thus these are mutually inverse maps over the specified base.  In
particular, transporting multiplication, identity and inverse through
$\pi$ gives the existing group scheme in this identity-product model.

The original $\operatorname{isAb}$ fields imply that $X\to S$ is
smooth, proper and geometrically connected.  Properness gives finite
type and separatedness over the field, so $X$ is Noetherian.  Its local
rings are regular, hence domains.  Two distinct irreducible components
cannot meet: at an intersection point they would give distinct
minimal primes of the same local ring.  More explicitly, take an
affine neighborhood of that point.  The distinct component primes
contained in its point prime remain distinct on localization; if an
element of one entered the other after localization, multiplying it
by an element outside the point prime would contradict primality.
A domain has only one minimal prime, giving the contradiction.

There are finitely many irreducible components.  They are disjoint
closed sets, hence each is also open.  Connectedness and the identity
point therefore imply that there is exactly one component.  Every
local ring is reduced, so $X$ is reduced as well and is integral.

The argument applies after every extension of $K$: smoothness,
properness, geometric connectedness and the identity section all
persist.  Thus $X$ is geometrically integral.  It is a complete
connected group variety, which is the abelian-variety convention in
Milne's article.  Finally the original smooth relative dimension $g$
transports to $f$.  Over a field every component of its unique fiber
has dimension $g$; this fiber is $X$ itself.  Hence the dimension in
(SGV.2) is exactly $P.g$.
\end{proof}
% END SGV1

% BEGIN SGV2
\paragraph{SGV2: the literal section module and its scalar field.}
\label{claim:SGV2}
Let $\mathcal L$ be the invertible sheaf defined by the original
chart datum $L$.  For every $j\in\mathbb Z$ there is a natural
coordinate-preserving identification
\[
 V_j\simeq H^0\bigl(X,(\pi^*\mathcal L)^{\otimes j}\bigr).
                                                        \tag{SGV.3}
\]
It is linear after transporting scalars through the canonical field
isomorphism $e:K\simeq R_0$.  Consequently the corresponding finite
dimensions over $K$ and $R_0$ agree.

\begin{proof}
The original datum consists of opens $U_i$ and transition units
$g_{ij}$ satisfying the cover and cocycle identities.  Glue
$\mathcal O_{U_i}$ by these units.  The resulting sheaf is invertible,
and a global section is exactly a tuple $a_i$ with
$a_i=g_{ij}a_j$ on each overlap, with the restrictions understood.
The gluing construction preserves sums and multiplication by
functions on the base.

By definition $\operatorname{LPow}(A,L,j,T)$ is the
singleton-factor $\operatorname{LinComb}$ of $L$ along $\pi$,
with exponent $j$.  Its charts are indexed by functions
$\operatorname{Fin}(1)\to L.\iota$.  Evaluation at $0$ identifies
these with the original chart indices.  The chart opens are
$\pi^{-1}U_i$, and the transition units are $(\pi^*g_{ij})^j$.
They are exactly those of $(\pi^*\mathcal L)^{\otimes j}$,
using the dual bundle when $j<0$.

The submodule $\operatorname{SecM}$ imposes precisely these overlap
equations.  The source's maps $\operatorname{SecM.toSec}$ and
$\operatorname{SecM.ofSec}$ leave every chart coordinate unchanged,
so the preceding description proves (SGV.3) as a bijection of the
actual sections, not merely as an equality of dimensions.

For $b\in R_0$ the source action on a coordinate over an open
$W\subseteq X$ is multiplication by $f^*b|_W$.
The canonical isomorphism
$e:K\simeq\Gamma(\operatorname{Spec}K,\mathcal O)$ identifies
this action with the ordinary $K$-action on the right of (SGV.3).
Thus the identification preserves addition and the specified scalar
action.  A basis over one field transports to a basis over the other,
which proves the dimension assertion.  No sign restriction on $j$
was used.
\end{proof}
% END SGV2

% BEGIN SGV3
\paragraph{SGV3: one fixed very ample power, with nonzero sections in every positive power.}
\label{claim:SGV3}
There is an integer $N\ge1$ such that
$H=(\pi^*\mathcal L)^{\otimes N}$ is very ample.
The integer $N$ is chosen before any positive $t$.  For every
$t\ge1$, the sheaf $H^{\otimes t}$ is ample and has a nonzero
global section.  Its actual coordinate sections identify with $V_{Nt}$.

\begin{proof}
Apply PBC26 to the displayed identity Cartesian square, with $R=K$.
It supplies $N\ge1$ and actual sections $s_0,\ldots,s_m$ of
$\operatorname{pw}(L.\operatorname{pull}(\pi),N)$ satisfying
$\operatorname{VeryAmpleBy}(f,\operatorname{pw}(\pi^*L,N),s)$.
The associated invertible sheaf is $H$.

Here the original coordinate form of very ampleness gives classical
very ampleness directly.  The opens $D_i$ on which $s_i$ is nonzero
cover $X$, and $s_i$ trivializes $H$ on $D_i$.  The supplied ratio
functions satisfy $r_j=s_j/s_i$ there, so $r_i=1$.
These ratios define the morphism $X\to\mathbb P^m_K$ given by the
sections.  Its inverse image of the $i$th standard affine chart is
exactly $D_i$.  The specified polynomial evaluation map to
$\Gamma(D_i,\mathcal O)$ is surjective.  Since it sends its $i$th
variable to $1$, it factors through that relation and gives a
surjective map from the polynomial ring in the remaining $m$ chart
coordinates.  Each map to the standard affine chart is therefore a
closed immersion.  Closed immersion is local on the target, so they
give a closed immersion into $\mathbb P^m_K$ whose pullback of
$\mathcal O(1)$ is $H$.

Choose an $i$ for which $s_i$ is nonzero at the identity point of $X$.
Such a choice is possible because the $D_i$ cover.  For every $t\ge1$,
$s_i^{\otimes t}$ is nonzero at that same point: it is the tensor
power of a nonzero vector in a one-dimensional $K$-space.
This proves $H^0(X,H^{\otimes t})\ne0$.  Composition of the displayed
closed immersion with the degree-$t$ Veronese embedding proves that
$H^{\otimes t}$ is very ample and hence ample, in any characteristic.

Lastly, iterating the power-chart operation gives transition units
$((\pi^*g_{ij})^N)^t=(\pi^*g_{ij})^{Nt}$.  Identifying the singleton
chart indices as in SGV2 gives the actual coordinate correspondence
with $V_{Nt}$.  The restriction maps and scalar action are unchanged.
This is the coordinate content of the source's mutually inverse
$\operatorname{toPw}$ and $\operatorname{ofPw}$ maps.
\end{proof}
% END SGV3

% BEGIN SGV4
\paragraph{SGV4: concentration in degree zero for every positive power of \(H\).}
\label{claim:SGV4}
For every integer $t\ge1$,
\[
 H^q(X,H^{\otimes t})=0\quad(q>0),\qquad
 h^0(X,H^{\otimes t})=\chi(X,H^{\otimes t}).
                                                        \tag{SGV.4}
\]

\begin{proof}
By SGV1 the scheme is an abelian variety over the original field.
By SGV3 the sheaf $H^{\otimes t}$ is ample and has a nonzero
global section.  Milne's Proposition~9.1 therefore gives dimension
zero for its translation stabilizer.  Theorem~13.3(c) says that its
nonzero cohomology occurs in exactly one degree.  Since degree zero
is already nonzero, this degree must be zero.  Thus every higher
cohomology group vanishes, and the alternating sum defining the Euler
characteristic reduces to $h^0$.

This uses only the dimension-zero condition on the reduced translation
stabilizer in Milne's statement.  It does not require identifying it
with the original $\operatorname{KTriv}$ functor, and it does not use a
separability or prime-to-characteristic assumption on $t$.
\end{proof}
% END SGV4

% BEGIN SGV5
\paragraph{SGV5: the exact positive-dimensional growth formula.}
\label{claim:SGV5}
If $g>0$, then for every integer $t\ge1$,
\[
 h^0(X,H^{\otimes t})=t^g h^0(X,H).
                                                        \tag{SGV.5}
\]

\begin{proof}
Use the nonzero section $s_i$ chosen in SGV3.  On the integral scheme
$X$ it defines an effective Cartier divisor $D$ with
$H\simeq\mathcal O_X(D)$.  To see the Cartier condition explicitly,
its coordinate on any nonempty trivializing open is not zero: a zero
coordinate would force the section to vanish on that open, then on
every other chart through nonempty overlaps and the transition units.
The local coordinate rings are domains, so these nonzero coordinates
are nonzerodivisors.  Their principal ideals fit together because the
coordinates differ by units.  This is exactly the asserted effective
Cartier divisor, with the original bundle as its associated sheaf.

Tensor powers now give $H^{\otimes t}\simeq\mathcal O_X(tD)$.
By SGV4, Milne's Riemann--Roch formula and multilinearity of the
intersection product,
\[
 \begin{aligned}
 h^0(X,H^{\otimes t})
   &=\chi(X,\mathcal O_X(tD))
     =\frac{((tD)^g)}{g!}\\
   &=t^g\frac{(D^g)}{g!}
     =t^g\chi(X,H)
     =t^g h^0(X,H).
 \end{aligned}
\]
All intersections and Euler characteristics here are integer-valued;
the displayed fractions are interpreted in $\mathbb Q$.  We never
divide by $g!$ inside $K$.  Therefore this calculation remains valid
when the characteristic divides $g!$ or $t$.
\end{proof}
% END SGV5

% BEGIN SGV6
\paragraph{SGV6: the zero-dimensional case.}
\label{claim:SGV6}
If $g=0$, then $X\simeq\operatorname{Spec}K$, and (SGV.5) holds
for every $t\ge1$, with $h^0(X,H)=1$.

\begin{proof}
An integral zero-dimensional scheme has just its generic point:
any further point would give a strict chain of irreducible closed
subsets of length one.  Thus an affine neighborhood of the unique
point of $X$ is all of $X$.  Its coordinate ring is a domain of
dimension zero, hence a field.  It is a finite-type $K$-algebra by
SGV1, so Zariski's lemma makes it a finite algebraic extension of $K$.
Algebraic closedness gives the equality with $K$.

An invertible sheaf on $\operatorname{Spec}K$ is a one-dimensional
$K$-vector space.  Every tensor power again has one-dimensional
sections.  Since $t^0=1$, this proves the formula, including $t=1$,
without using intersection theory in dimension zero.
\end{proof}
% END SGV6

% BEGIN SGV7
\paragraph{SGV7: the complete original section-finiteness and growth output.}
\label{claim:SGV7}
For every original $P$ over every algebraically closed field $K$,
$V_j$ is finite over $R_0$ for every $j\in\mathbb Z$.
There exist integers $N,d\ge1$, chosen before $t$, such that
\[
 \dim_{R_0}V_{Nt}=d\,t^g\qquad(t\ge1).
                                                        \tag{SGV.6}
\]
In particular, the full original conclusion is supplied:
\[
 \begin{gathered}
  (\forall j\in\mathbb N)\quad V_j\text{ is finite over }R_0,\\
  \exists N,c,C\in\mathbb N,\quad N,c\ge1,\quad
  \forall t\in\mathbb N,\ t\ge1:\quad
  t^g\le c\dim_{R_0}V_{Nt},\qquad
  \dim_{R_0}V_{Nt}\le C t^g.
 \end{gathered}
                                                        \tag{SGV.7}
\]

\begin{proof}
PBC27 gives finiteness of the actual $V_j$ for every integer $j$.
Choose the single $N$ from SGV3 and put $d=h^0(X,H)$.
This is a finite integer by SGV2 and the $j=N$ instance of PBC27.
It is positive because SGV3 gives a nonzero section.

SGV5 and SGV6 give (SGV.5) in every dimension.  The actual
power-coordinate comparison in SGV3 identifies the left side with
$V_{Nt}$, and SGV2 transports its dimension to the precise coefficient
field $R_0$.  This proves (SGV.6).  Since $d\ge1$, the choices
$c=1$ and $C=d$ prove both inequalities in (SGV.7).
Every choice precedes the universal quantifier over $t$.
The integers may depend on the given $K$ and $P$, as in the original
source; no uniformity across different polarized varieties is claimed.
There is no characteristic restriction on $K$, $N$ or $t$.
\end{proof}
% END SGV7

\subsection*{What this supplies next}

Equation (SGV.7) is exactly the finite-sections and growth output of
\nolinkurl{CubeAssembly.sections_growth}, and hence the growth input
called \texttt{hgrow} by
\nolinkurl{finrank_sections_lineBundlePow_eq_pow}.  The source obtains
it from several general geometric and polynomial-growth suppliers.
The present argument reaches the same original output directly using
the explicit classical input, while preserving the actual identity
product, section coordinates and scalar field.

For GSL19, SGV1 supplies the needed integrality of the original
geometric fiber and SGV7 supplies its fourth input.  The rank-one
statement $\dim V_1=1$ and the actual torsion-indexed section
equivalences remain to be proved.  In particular, the constant $d$ in
(SGV.6) is not declared to be $N^g$; establishing all ranks $j^g$
requires those further inputs.  The full all-cover eventual Serre
vanishing theorem remains supplied by SHV17, and its homogeneous-polynomial
companion remains supplied by SHV21.  The later descent theorem giving
all-cover vanishing for every positive power $L^n$ is a separate input,
which SGV does not supply.  Cube and square generators, translated-power
generators, the global $\operatorname{SecBaseChangeSt}$,
$\operatorname{SecBaseChangeAnySt}$, $\operatorname{SecTwistSt}$ and
moduli objects also remain outside this chapter.  No selected whole
obligation is closed here.

\subsection*{Exact original source correspondence}

All source paths below use commit
\nolinkurl{7cbf7a0c32fca1badf05931879a37c3f8a2b03b4} of
\url{https://github.com/CocoSgt/riemann-hypothesis-lean}.
Let $\mathcal A$ denote
\texttt{src/\allowbreak Development/\allowbreak Abelian\allowbreak Variety\allowbreak Theory/\allowbreak Proofs/\allowbreak Abelian\allowbreak Scheme\allowbreak Law\allowbreak Proofs/\allowbreak },
and $\mathcal R$ denote
\texttt{src/\allowbreak Reciprocity\allowbreak Theory/\allowbreak Reductive\allowbreak Groups/\allowbreak Arithmetic\allowbreak Geometry/\allowbreak }.
\begin{enumerate}
 \item The actual record is
 $\mathcal R$\texttt{Extension\allowbreak Descent/\allowbreak Nonzero\allowbreak Locus.\allowbreak lean}, lines59--73;
 its abelian-scheme fields are in
 $\mathcal R$\texttt{Semiabelian\allowbreak Schemes/\allowbreak Abelian\allowbreak Scheme\allowbreak Quotients/\allowbreak Is\allowbreak Abelian\allowbreak Scheme.\allowbreak lean},
 lines13--17.
 \item The line-bundle charts are
 $\mathcal R$\texttt{Semiabelian\allowbreak Schemes/\allowbreak Cubical\allowbreak Structures/\allowbreak Cube\allowbreak Symmetric\allowbreak Swap\allowbreak First.\allowbreak lean},
 lines17--78.  The literal power, action and section module are
 $\mathcal R$\texttt{Extension\allowbreak Descent/\allowbreak Theta\allowbreak Pt.\allowbreak lean}, lines53--57, and
 $\mathcal R$\texttt{Extension\allowbreak Descent/\allowbreak Theta\allowbreak VV.\allowbreak lean}, lines18--60.
 \item The target and actual power-coordinate maps are in
 $\mathcal A$\texttt{Finite\allowbreak Of\allowbreak Not\allowbreak Linear\allowbreak Independent.\allowbreak lean}, lines280--350;
 its exported output is
 $\mathcal A$\texttt{Cube\allowbreak Theorem\allowbreak And\allowbreak Sections\allowbreak Finite\allowbreak Growth\allowbreak Abelian\allowbreak Variety.\allowbreak lean},
 lines38--43.
 \item The exact next rank consumer is
 $\mathcal A$\texttt{Finrank\allowbreak Sections\allowbreak Line\allowbreak Bundle\allowbreak Pow\allowbreak Eq\allowbreak Pow.\allowbreak lean}, lines228--260.
 The separate inputs at the later global join appear in
 $\mathcal A$\texttt{Laws\allowbreak Of\allowbreak H0Eq\allowbreak One\allowbreak Closed\allowbreak Fields.\allowbreak lean}, lines64--93.
\end{enumerate}


\section{Faithful cyclotomic invariants and the actual etale functor}
\label{sec:faithful-etale-v057}

\status{This chapter proves the original faithful-invariants statement,
constructs a fixed basis on the prescribed integral Witt tensor lattice,
constructs the Frobenius change of basis into the original coefficient
ring, and returns the full original etale functor, including its maps on
all lattice morphisms. The qualitative perfectoid finite-extension
theorem is used as a precisely stated classical foundation for the first
step. All later averaging and conjugacy arguments are written out.
The original good layer and period records remain supplied data.
Herr cohomology, the quantitative comparison and the common global
denominator family remain separate.}

The problem is to make the action of the cyclotomic kernel disappear
from a basis without changing either the lattice or its coefficient
embedding. There are two steps. A finite averaging argument followed by
an integral limit produces a fixed basis. Its Frobenius matrix initially
has entries in a larger fixed ring. A second, nonlinear change of basis
puts those entries in the specified smaller ring. The previously proved
IFD25 concerns descent of a matrix already satisfying an equation; it
does not by itself construct this second change of basis. Below we give
that construction before using IFD25 to descend morphisms.

Fix a prime $p$; the argument includes $p=2$. Let
$\overline{\mathbb Q}_p\subset C$ be the original algebraic closure
and its completion. Put
\[
 S=\mathcal O_C^\flat=\varprojlim_{x\mapsto x^p}\mathcal O_C/p,
 \qquad F=\operatorname{Frac}(S),\qquad W_0=W(F).
                                                        \tag{FET.1}
\]
For a finite coefficient extension $LL/\mathbb Q_p$ put
$O=\mathcal O_{LL}$ and $\mathcal W=O\otimes_{\mathbb Z_p}W_0$.
For a continuous $O$-lattice $T$ with $G_K$-action put
$M_T=\mathcal W\otimes_O T$ and $H=G_K\cap\ker\chi$.
The notation $T$ refers to the integral lattice itself. Its scalar
extension to $LL$ is not substituted for it.

The original records $P,AI$ prescribe the actions on $C,S,W_0$ and
$\mathcal W$. Their identity with the actions just described is part of
the proof, rather than a change in the input data. Later, on a good
layer, $E=R.\mathrm{EE}$ and $j=I.\mathrm{ioE}:E\to\mathcal W$
will denote the original coefficient ring and map.

\subsection{Why the tilt separates the original finite extensions}

% BEGIN FET1
\paragraph{FET1: the prescribed actions are the canonical actions.}
\label{claim:FET1}
For every original $P,AI$, the action $AI.\mathrm{actC}$ is the
unique isometric extension of the given action on
$\overline{\mathbb Q}_p$. The action $AI.\mathrm{actT}$ is its
coordinatewise action on the inverse limit in (FET.1). Moreover,
$\ker(P.\chi_Q)$ fixes every $p$-power root of unity, and is exactly
the usual cyclotomic kernel.

\begin{proof}
Two isometries of $C$ which agree on the dense subfield
$\overline{\mathbb Q}_p$ agree on a convergent sequence approaching
each point of $C$. The fields \texttt{act\allowbreak C\_\allowbreak ext} and
\texttt{act\allowbreak C\_\allowbreak norm} therefore determine the first action uniquely.
For $x\in S$, every coordinate of $AI.\mathrm{actT}(g)x$ is
specified by \texttt{act\allowbreak T\_\allowbreak coeff}: lift the corresponding residue
coordinate of $x$ to $\mathcal O_C$, apply the already determined
isometry, and reduce modulo $p$. Equality of all coordinates proves
the second assertion.

For a primitive $p^m$-th root $\zeta$, the field
\texttt{chi\allowbreak Q\_\allowbreak cyc} gives an integer $a$ whose distance from
$P.\chi_Q(g)$ is at most $p^{-m}$ and for which $g\zeta=\zeta^a$.
The ordinary cyclotomic exponent has the same residue modulo $p^m$.
For all $m$ these residues determine one element of $\mathbb Z_p$,
so it equals $P.\chi_Q(g)$ in $\mathbb Q_p$.
This also follows from the original integral-character result WPR2.
In particular a value $1$ fixes every root, and the converse holds by
uniqueness of the residues. No alternate action or character has been
introduced.
\end{proof}
% END FET1

% BEGIN FET2
\paragraph{FET2: the completed cyclotomic extension of a local field.}
\label{claim:FET2}
For a finite extension $L/\mathbb Q_p$ in the original algebraic
closure, let
\[
 L_\infty=L(\mu_{p^\infty}),\qquad
 E_L=\overline{L_\infty}^{\,C}.
\]
Then $E_L$ is a perfectoid field. In particular its ring of integers
is complete for the $p$-adic topology and Frobenius on
$\mathcal O_{E_L}/p$ is onto.

\begin{proof}
Write $E_0=\widehat{\mathbb Q_p(\mu_{p^\infty})}$.
First one can see its perfectoid property directly. The integral
ring at level $m$ is $\mathbb Z_p[\zeta_{p^m}]$; the usual Eisenstein
polynomial for $\zeta_{p^m}-1$ proves this assertion, including the
initial trivial levels when $p=2$. Modulo $p$, a polynomial in
$\zeta_{p^m}$ has a $p$-th root at level $m+1$: reduce its
coefficients to $\mathbb F_p$ and replace $\zeta_{p^m}$ by
$\zeta_{p^{m+1}}$. Every integral element of the completion is
congruent modulo $p$ to an integral element at a finite level, by
density. Thus Frobenius on $\mathcal O_{E_0}/p$ is onto.
The valuations of $\zeta_{p^m}-1$, namely
$1/(p^{m-1}(p-1))$, show that the value group is nondiscrete.
Together with completeness these are the perfectoid-field conditions.

We now use the classical finite-extension theorem for perfectoid
fields: every finite extension of a perfectoid field, with its natural
complete valuation topology, is perfectoid. This is Scholze,
\emph{Perfectoid spaces}, Theorem~3.7(i); this named foundation is
the only deep input in the present separation argument.
The compositum $LE_0$ inside $C$ has finite degree over $E_0$ and
is complete, hence closed in $C$. It contains $L_\infty$.
Conversely $E_L$ contains both $L$ and $E_0$, hence $LE_0$.
Thus $E_L=LE_0$, so the theorem applies to the actual embedded field.
Equivalently, Bhatt--Morrow--Scholze, \emph{Integral $p$-adic Hodge
theory}, Example~3.6, states this integral perfectoid example directly.
There is no Galois or unramified assumption on $L$.
\end{proof}
% END FET2

% BEGIN FET3
\paragraph{FET3: recovering integral elements from their tilt coordinates.}
\label{claim:FET3}
Let $E\subset C$ be a closed perfectoid subfield containing
$\mathbb Q_p$. Its integral tilt
$S_E=\varprojlim_{x\mapsto x^p}\mathcal O_E/p$ embeds in $S$.
For $x=(\overline x_m)\in S_E$ and arbitrary lifts
$x_m\in\mathcal O_E$, the limit
\[
                     x^\sharp=\lim_m x_m^{p^m}             \tag{FET.2}
\]
exists, is independent of the lifts, and reduces to
$\overline x_0$ modulo $p$. Every $a\in\mathcal O_E$ admits an
expansion
\[
                      a=\sum_{n\geq0}p^n x_n^\sharp,
                      \qquad x_n\in S_E.                 \tag{FET.3}
\]
If an isometric automorphism of $C$ fixes every image of $S_E$ in
$S$, it fixes $E$ pointwise.

\begin{proof}
The injection $\mathcal O_E/p\hookrightarrow\mathcal O_C/p$
follows from $p\mathcal O_C\cap\mathcal O_E=p\mathcal O_E$:
divide an element of the intersection by $p$ in $E$ and use its
valuation. Passing to inverse limits is therefore injective.
If $u-v\in p^a\mathcal O_E$, with $a\geq1$, then the binomial
formula gives $u^p-v^p\in p^{a+1}\mathcal O_E$. The terms with
$1\leq i<p$ have a factor $p$, while $pa\geq a+1$ handles the
last term, also for $p=2$. Consequently compatible residue lifts
satisfy
\[
 x_{m+1}^{p^{m+1}}-x_m^{p^m}\in p^{m+1}\mathcal O_E.
\]
Changing a lift has the same estimate. Completeness proves (FET.2)
and its stated residue property.

Surjectivity of Frobenius permits successive choices of roots of
any residue class, so the zeroth projection
$S_E\to\mathcal O_E/p$ is onto. Given $a_0=a$, choose $x_0$
whose sharp has the same residue, and set
$a_1=(a_0-x_0^\sharp)/p\in\mathcal O_E$. Repeat.
The error after $n$ terms is $p^n a_n$, proving (FET.3).
Formula (FET.2) is equivariant for isometric automorphisms, even
when the automorphism is not assumed to preserve $E$ in advance.
Fixing all the $x_n$ therefore fixes every sum in (FET.3).
Finally every element of $E$ becomes integral after multiplication
by a sufficiently large power of $p$. This proves the last assertion.
The sharp map has not been treated as additive.
\end{proof}
% END FET3

% BEGIN FET4
\paragraph{FET4: full original faithful cyclotomic invariants.}
\label{claim:FET4}
For every finite $L/\mathbb Q_p$ and every
$g\in G_{\mathbb Q_p}$, if $AI.\mathrm{actT}(g)$ fixes every
element of $S$ fixed by $G_L\cap\ker\chi$, then $g\in G_L$.
In particular the full original predicate \texttt{Faith\allowbreak H AI} holds.

\begin{proof}
Every $h\in G_L\cap\ker\chi$ fixes $L_\infty$ pointwise.
Its isometric extension fixes the closure $E_L$ by FET1, hence
fixes every coordinate of $S_{E_L}\subset S$.
The assumption on $g$ therefore fixes all of $S_{E_L}$.
FET2 and FET3 imply that $g$ fixes $E_L$, and hence $L$.
This is precisely membership in $G_L$.
The original statement additionally assumes $\chi(g)=1$; the
proof shows that this antecedent is unnecessary for this conclusion.
It remains retained when the statement is returned to its original
consumer. The construction has used the original $AI$ action and
the original embedded $L$, for every supplied $P,AI$.
\end{proof}
% END FET4

\subsection{An integral fixed basis, obtained by finite averages}

For this subsection fix $K,LL,T$. Let $\mathcal N$ be the family
of subgroups $N=H\cap G_L$, where $L/\mathbb Q_p$ is finite
Galois and contains $K$. These are open normal subgroups of $H$,
closed under passing to a common smaller member. The coefficient
action on $W_0$ will be denoted $\sigma_h$, and the semilinear
action on $M_T$ by $\tau_h$.

% BEGIN FET5
\paragraph{FET5: finite quotients act by distinct integral-tilt embeddings.}
\label{claim:FET5}
For $N=H\cap G_L\in\mathcal N$, the maps
\[
 S^N\longrightarrow F,\qquad x\longmapsto h(x),
                   \quad hN\in H/N,                      \tag{FET.4}
\]
are distinct. They have an infinite common set of values obtained
from elements of $S^N$: the powers of $\epsilon-1$, where
$\epsilon$ is the prescribed cyclotomic tilt element.

\begin{proof}
The maps are independent of a chosen representative because $N$
fixes $S^N$. If two maps are equal, $h'^{-1}h$ fixes $S^N$.
Here $N=G_L\cap\ker\chi$, since $L$ contains $K$.
FET4 gives $h'^{-1}h\in G_L$, so it belongs to $N$.
Thus the two cosets are equal.

The prescribed roots defining $\epsilon$ are all fixed by $H$,
so $\epsilon$ and $\epsilon-1$ are fixed. They are the same roots
as in FET1, by the original \texttt{zeta\_\allowbreak eq} field.
The element $\epsilon-1$ is nonzero: a sufficiently high primitive
$p$-power root has $0<v_p(\zeta-1)<1$, and its residue coordinate
is not $1$. It is a nonunit of $S$, because its zeroth residue
coordinate is $0$. In the domain $S$, a nonzero nonunit cannot
have equal distinct nonnegative powers; cancellation would make
one positive power equal to $1$. Their images in $F$ are therefore
an infinite set fixed by every map in (FET.4).
\end{proof}
% END FET5

% BEGIN FET6
\paragraph{FET6: a determinant can be made nonzero by one common average.}
\label{claim:FET6}
Let $A$ be a commutative ring, $C_0$ a field, and
$s_q:A\to C_0$ distinct unital ring maps for a finite nonempty
set $Q$. Suppose an infinite set $B\subset C_0$ consists of
common values: for each $b\in B$ there is $a\in A$ with
$s_q(a)=b$ for every $q$. For square matrices $B_q$ of a fixed
finite size over $C_0$, if $\det B_{q_0}\ne0$ for one $q_0$,
then
\[
              \det\left(\sum_q s_q(y)B_q\right)\ne0
                   \quad\hbox{for some }y\in A.           \tag{FET.5}
\]
The size-zero case is included, with determinant $1$.

\begin{proof}
Distinct ring maps into a field are linearly independent as
functions. For completeness, in a shortest nonzero linear relation
choose two distinct maps, evaluate at $ab$, and subtract one of
their values at $b$ times the relation at $a$. Choosing $b$ where
the two maps differ gives a shorter nonzero relation, a contradiction.
It follows that the vectors $(s_q(a))_q$, as $a$ ranges over $A$,
span $C_0^Q$: a proper span would have a nonzero annihilating
linear functional and hence such a relation.

Choose finitely many $a_1,\ldots,a_m$ whose vectors span, and
consider
\[
 P(X_1,\ldots,X_m)=
 \det\left(\sum_q\sum_{j=1}^m X_j s_q(a_j)B_q\right).
\]
There are coefficients whose vector of inner sums is the unit
vector at $q_0$. At them $P=\det B_{q_0}\ne0$, so $P$ is not
the zero polynomial. A nonzero polynomial over a field cannot
vanish on $B^m$ when $B$ is infinite; prove this by induction
on $m$, using the finite root bound in one variable.
Choose $b_j\in B$ where it does not vanish, and choose common
preimages $c_j\in A$. The element $y=\sum_j c_j a_j$ gives
$s_q(y)=\sum_j b_j s_q(a_j)$ and proves (FET.5).
This proof does not assume that the infinite set is a subfield.
\end{proof}
% END FET6

% BEGIN FET7
\paragraph{FET7: a unit determinant over the original Witt ring.}
\label{claim:FET7}
Let $N\in\mathcal N$, let $Q=H/N$, and choose representatives
$h_q$. For matrices $B_q\in M_d(\mathbb Z_p)$ with one unit
determinant, there exists $y\in W_0^N$ such that
\[
                  \det\left(\sum_q\sigma_{h_q}(y)B_q\right)
                        \in W_0^\times.                  \tag{FET.6}
\]

\begin{proof}
Reduce the matrices modulo $p$ into $F$ and use FET5--FET6
with $A=S^N$. This gives $a\in S^N$ with nonzero determinant
in $F$. Take the Teichmuller lift $y=[a]\in W(F)$.
Functoriality of Witt vectors gives
$\sigma_h([a])=[h(a)]$, so $y$ is $N$-fixed and the determinant
in (FET.6) has the prescribed nonzero zeroth Witt coordinate.
For a perfect field $F$, $W(F)$ is $p$-adically complete and an
element is a unit exactly when its reduction in $F$ is nonzero.
These are the actual Witt-ring properties already established in
WIS2. Thus the determinant is a unit. Neither $|Q|$ nor $p$
has been inverted.
\end{proof}
% END FET7

% BEGIN FET8
\paragraph{FET8: exact trace-one averaging coefficients.}
\label{claim:FET8}
For every $N\in\mathcal N$ there is $b_N\in W_0^N$ with
\[
                      \sum_{q\in H/N}\sigma_{h_q}(b_N)=1.
                                                               \tag{FET.7}
\]

\begin{proof}
Apply FET7 to the one-by-one matrices all equal to $1$.
For the resulting $y$, the sum $u=\sum_q\sigma_{h_q}(y)$ is
a unit. Left multiplication by any $h\in H$ permutes the cosets,
and the right discrepancies of representatives belong to $N$.
Since $y$ is $N$-fixed, this proves $\sigma_h(u)=u$.
The inverse $u^{-1}$ is then fixed as well.
The element $b_N=yu^{-1}$ is $N$-fixed and has trace $1$.
The same coefficient, embedded as $1\otimes b_N$ in
$\mathcal W$, works for $M_T$.
\end{proof}
% END FET8

% BEGIN FET9
\paragraph{FET9: finite averaging improves invariance at a chosen precision.}
\label{claim:FET9}
Let $J_m=p^{m+1}M_T$. Call $x\in M_T$ smooth modulo these
filtrations if for each $m$ some $N\in\mathcal N$ satisfies
$\tau_\nu x-x\in J_m$ for all $\nu\in N$.
If this holds and $N$ has that property at precision $m$, then
for any $c\in\mathcal W^N$ the finite sum
\[
                    A_{N,c}(x)=\sum_{q\in H/N}\tau_{h_q}(cx)
                                                               \tag{FET.8}
\]
is smooth modulo every $J_k$ and is invariant modulo $J_m$
under all of $H$.

\begin{proof}
For $h\in H$, write $hh_q=h_{q'}\nu$ with $\nu\in N$.
Subtract the term indexed by $q'$ after applying $\tau_h$.
Its error is $\tau_{h_{q'}}(c(\tau_\nu x-x))\in J_m$.
Here semilinearity, $\sigma_\nu(c)=c$ and stability of $J_m$
are precisely what is needed. Summing proves invariance modulo
$J_m$.

For smoothness at $J_k$, choose a normal member $N'\subset N$
which fixes $x$ modulo $J_k$. Such a member exists by the
directed property of $\mathcal N$. For $\nu\in N'$, each
conjugate $h_q^{-1}\nu h_q$ lies in $N'$ and in $N$.
It therefore fixes $c$ exactly and $x$ modulo $J_k$.
Subtracting each of the finitely many terms of (FET.8) shows
that $N'$ fixes the average modulo $J_k$. This argument does
not require continuity of arbitrary Witt coefficients.
\end{proof}
% END FET9

% BEGIN FET10
\paragraph{FET10: the averaging limit is exactly invariant.}
\label{claim:FET10}
If $x_0\in M_T$ is smooth modulo the $J_m$, if $N_0$ fixes
it modulo $pM_T$, and if $c_0\in\mathcal W^{N_0}$, then there
exists $x\in M_T^H$ such that
\[
                   x\equiv A_{N_0,c_0}(x_0)\pmod {pM_T}.
                                                               \tag{FET.9}
\]

\begin{proof}
Set $x^{(0)}=A_{N_0,c_0}(x_0)$. FET9 makes it invariant
modulo $J_0$ and smooth. Suppose $x^{(m)}$ is invariant modulo
$J_m$. Choose $N$ fixing it modulo $J_{m+1}$ and average
with the trace-one coefficient in FET8 to define $x^{(m+1)}$.
It is invariant modulo $J_{m+1}$, and
\[
 x^{(m+1)}-x^{(m)}
 =\sum_q\sigma_{h_q}(1\otimes b_N)
                       (\tau_{h_q}x^{(m)}-x^{(m)})\in J_m.
\]
Every iterate is smooth by FET9. The original ring $\mathcal W$
and its finite free module $M_T$ are $p$-adically complete and
separated (WIS2--WIS4). Hence $x^{(m)}$ tends to an $x$ with
$x-x^{(m)}\in J_m$. For each fixed $h$, its semilinear action
preserves every $J_m$. Thus $\tau_hx-x\in J_m$ for every $m$,
and separation makes this difference zero. The congruence at
$m=0$ gives (FET.9).
\end{proof}
% END FET10

% BEGIN FET11
\paragraph{FET11: one initial average is an automorphism.}
\label{claim:FET11}
Choose an $O$-basis $b_1,\ldots,b_r$ of $T$ and a
$\mathbb Z_p$-basis $c_1,\ldots,c_d$ of the same $T$.
There is $N_0\in\mathcal N$ fixing every $b_i$ modulo $pT$.
There is also $y\in W_0^{N_0}$ for which the $\mathcal W$-linear
operator
\[
 U=\sum_{q\in H/N_0}(1\otimes\sigma_{h_q}(y))\rho(h_q)
                         \quad\hbox{on }M_T               \tag{FET.10}
\]
is an automorphism.

\begin{proof}
The lattice action is continuous in its finite
$\mathbb Z_p$-coordinates. Its reduction on $T/pT$ therefore
has an open kernel. An open normal subgroup $G_L$ of
$G_{\mathbb Q_p}$ can be chosen inside the corresponding
neighborhood and inside $G_K$; intersecting with $H$ gives
$N_0$. The same argument works at every $p$-power precision.

In the basis $c$, every $\rho(h_q)$ has an invertible integral
matrix $B_q\in\mathrm{GL}_d(\mathbb Z_p)$. FET7 makes
$\sum_q\sigma_{h_q}(y)B_q$ invertible over $W_0$.
Under the natural identification
$M_T\simeq W_0\otimes_{\mathbb Z_p}T$, this is the matrix of
the underlying $W_0$-linear map $U$. Thus $U$ is bijective.
It commutes with $O$, since all $\rho(h_q)$ do, and with
$W_0$ by its definition. Its inverse therefore commutes with
$\mathcal W$, so it is a $\mathcal W$-automorphism.
This use of a larger $\mathbb Z_p$-basis is only a way to test
invertibility; the eventual basis has the original $O$-rank $r$.
\end{proof}
% END FET11

% BEGIN FET12
\paragraph{FET12: the fixed vectors form a basis on the unchanged lattice.}
\label{claim:FET12}
There is a basis $e_1,\ldots,e_r$ of $M_T$ over $\mathcal W$
fixed by every $h\in H$, where $r=\operatorname{rank}_O T$.

\begin{proof}
For each $1\otimes b_i$, apply FET10 with the common $N_0,y$
of FET11. Smoothness at every precision follows from the
continuity argument there. This gives fixed $e_i$ with
$e_i-U(1\otimes b_i)\in pM_T$.
The columns $U(1\otimes b_i)$ form a $\mathcal W$-basis.
In that basis the matrix of the $e_i$ is $1+pZ$ for an
integral matrix $Z$. The convergent geometric matrix series
$\sum_{n\geq0}(-pZ)^n$ is its inverse.
Thus the $e_i$ are a basis on $M_T$ itself. This argument
uses no division by $p$, no change of $O$ and no enlargement
of $T$. If $r=0$, the empty basis proves the assertion directly.
\end{proof}
% END FET12

% BEGIN FET13
\paragraph{FET13: the full original cyclotomic-kernel fixed-basis theorem.}
\label{claim:FET13}
For every prime $p$, supplied $P,AI$, finite extensions $K,LL$ of $\mathbb Q_p$ and
original lattice $T$, there is a basis indexed by
$\operatorname{Fin}(\operatorname{rank}_{\mathcal O_{LL}}T)$
of the original $\mathrm{TW}(\mathrm{RW}(LL),T)$ satisfying
the original \texttt{Fix\allowbreak Basis AI T} predicate.

\begin{proof}
The identifications in (FET.1) are the literal definitions of
$\mathrm{OCb}$, $\mathrm{Cb}$ and $\mathrm{WCb}$. The original
$\mathrm{RW}(LL)$ is $\mathcal W$, and its lattice tensor and
action are exactly $M_T,\tau$. FET1 identifies the given actions;
FET4 supplies faithful invariants for every finite extension.
The family $\mathcal N$ used above is cofinal in the required
identity neighborhoods by the Krull topology and finite Galois
closures. All matrices use the original coefficient embedding.
FET12 therefore gives every quantified field of the exported
\texttt{exists\_\allowbreak cyclotomic\allowbreak Kernel\allowbreak Fixed\allowbreak Basis\_\allowbreak witt\allowbreak Tensor\allowbreak Lattice}.
The original implication with antecedent \texttt{Faith\allowbreak H} is
also proved, since FET5 onward only needs that antecedent.
No good-layer, crystalline or Wach-lattice assumption was used
in this fixed-basis theorem.
\end{proof}
% END FET13

\subsection{Constructing the Frobenius conjugacy into the coefficient ring}

Now fix one original good layer $K,s,LL,\Gamma,R,I,\mathrm{IoW}$.
Write $E=R.\mathrm{EE}$, $j=I.\mathrm{ioE}$,
$\mathcal S=\mathcal W^H$, and let $\phi$ be the prescribed
Frobenius. Keep the integral subring
$\mathcal B=O\otimes_{\mathbb Z_p}W(S)$ with its original
inclusion $c:\mathcal B\to\mathcal W$. Put
\[
 t=c(1\otimes([\epsilon_s]-1)),\qquad
 \mathcal N_m=t^m c(\mathcal B)+p\mathcal W\quad(m\geq0).
                                                               \tag{FET.11}
\]
As in IFD1, $t$ is a unit of $\mathcal W$, but (FET.11)
uses $c(\mathcal B)$ and is not the ideal
$t^m\mathcal W+p\mathcal W$. For matrices write
$X\in\mathcal N_m$ when all its entries lie there.

% BEGIN FET14
\paragraph{FET14: the exact two filtrations used in conjugacy.}
\label{claim:FET14}
The actual rings just fixed have the following properties.
The subgroups $\mathcal N_m$ decrease,
$\mathcal N_a\mathcal N_b\subset\mathcal N_{a+b}$,
$\phi(\mathcal N_m)\subset\mathcal N_{pm}$, and
$\bigcap_m\mathcal N_m=p\mathcal W$.
Every finite matrix $X$ over $\mathcal W$ has a common pole bound
$k$ with $t^kX\in\mathcal N_0$.
If $X\in M_r(\mathcal S)$, then for each $m$ there are $N,u,A$
such that
\[
 X-u\in\mathcal N_m,\qquad
                \phi^N(u)-j(A)\in pM_r(\mathcal W),
                         \quad A\in M_r(E).               \tag{FET.12}
\]
The ring $\mathcal S$ is $p$-saturated in $\mathcal W$, closed
under inverses of invertible matrices, and admits sums of sequences
with $i$-th term in $\mathcal S\cap\mathcal N_{i+m}$.
Both $\mathcal S$ and $E$ admit $p$-weighted sums with their
usual congruences, and $\mathcal W$ is $p$-adically separated.

\begin{proof}
IFD1--IFD2 give the multiplication, cancellation and Frobenius
properties of the precise subgroups (FET.11).
The intersection statement follows from WIS23: each zeroth Witt
coordinate belongs to every $\tau_s^m S$, hence is zero by
valuation separation. Thus the element lies in $p\mathcal W$.
Every element of $F$ acquires integral zeroth Witt coordinate
after multiplication by a power of the nonzero $\tau_s$.
Using the fixed finite coefficient basis gives the pole bound for
each entry, and the maximum gives one bound for a finite matrix.

Apply IFD16 to the finitely many entries of $X$, at the same
precision $m$. Choose a common maximum of their Frobenius
iteration exponents. Persistence under Frobenius of $j(E)+p\mathcal W$
(IFD14) gives (FET.12). It is not necessary that $u$ be fixed.
WIS20 gives saturation and $p$-weighted fixed-ring completeness,
and WIS25 gives the stated near-parameter sums at every shift.
If a fixed matrix has an inverse over $\mathcal W$, applying any
$h\in H$ to the inverse equation and using uniqueness shows that
the inverse is fixed. The original $p$-adic completeness of $E$
gives its $p$-weighted sums. WIS2--WIS4 give separation of
$\mathcal W$. These are previously proved properties of the
same $j,t,\phi,H$, not additional analytic assumptions.
\end{proof}
% END FET14

% BEGIN FET15
\paragraph{FET15: the near-parameter contraction solves a matrix equation.}
\label{claim:FET15}
Let $A,B\in M_r(\mathcal S)$ have common pole bound $k$,
and put $m=2k+1$. The additive operator
$\mathcal T(Z)=A\phi(Z)B$ sends
$\mathcal N_{i+m}$ into $\mathcal N_{i+m+1}$.
For $y\in M_r(\mathcal S)\cap\mathcal N_m$ there is
$z\in M_r(\mathcal S)\cap\mathcal N_m$ with
\[
                         z=y+\mathcal T(z)+pR
                      \quad\hbox{for some }R\in M_r(\mathcal W).
                                                               \tag{FET.13}
\]

\begin{proof}
Frobenius raises the filtration index from $i+m$ to $p(i+m)$.
Multiplication by the two matrices loses at most $2k$ by the
cancellation rule in IFD1. Since
$p(i+2k+1)-2k\geq i+2k+2$, the asserted contraction follows.
The sums of $\mathcal T^i(y)$ exist in the fixed ring by FET14.
Let $z$ be such a sum. Its error after $a$ terms lies in
$\mathcal N_{a+m}$, and $z\in\mathcal N_m$ by taking the
empty partial sum. Subtract the finite telescoping identity
from $z-y-\mathcal T(z)$. Contraction shows that this difference
belongs to every $\mathcal N_a$. FET14 puts it in
$pM_r(\mathcal W)$, proving (FET.13).
This is convergence modulo the zeroth Witt coordinate; it is
deliberately separate from the later $p$-adic iteration.
\end{proof}
% END FET15

% BEGIN FET16
\paragraph{FET16: Frobenius iteration makes the approximation sufficiently accurate.}
\label{claim:FET16}
For $F_0\in\mathrm{GL}_r(\mathcal W)$ with entries fixed by
$H$, there are $N\geq0$, $A\in M_r(E)$ and $k\geq0$ such that,
putting $F_1=\phi^N(F_0)$ and $G_1=F_1^{-1}$,
\[
 t^k F_1, t^kG_1, t^k j(A)\in\mathcal N_0,\qquad
                       F_1-j(A)\in\mathcal N_{3k+1}.
                                                               \tag{FET.14}
\]
Also $F_0$ and $F_1$ are Frobenius-conjugate by an invertible
matrix with entries in $\mathcal S$.

\begin{proof}
Choose a common pole bound $k_0$ for $F_0,F_0^{-1}$.
Apply (FET.12) at precision $3k_0+1$, obtaining $N,u,A$.
Frobenius multiplies the error precision by $p^N$.
Since $\phi(t)\equiv t^p\pmod {p\mathcal W}$, applying it
$N$ times to a pole bound multiplies that bound by $p^N$.
Set $k=p^N k_0$. Then
\[
 F_1-j(A)\in\mathcal N_{p^N(3k_0+1)}
                        \subseteq\mathcal N_{3k+1}.
\]
Subtracting this error from $F_1$ gives the same pole bound
for $j(A)$. Invariance and invertibility of $F_1,G_1$ follow
from FET14 and functoriality of Frobenius.

The product
$V_0=F_0\phi(F_0)\cdots\phi^{N-1}(F_0)$, with the empty
product $1$, satisfies
$F_0\phi(V_0)=V_0\phi^N(F_0)$.
It is invertible and fixed, proving the last assertion.
The choice of approximation precision preceded the unknown
iteration exponent $N$; the displayed inequality explains why
this order of choices is valid.
\end{proof}
% END FET16

% BEGIN FET17
\paragraph{FET17: the first conjugacy has coefficient values modulo $p$.}
\label{claim:FET17}
Every invertible fixed matrix $F_0$ is Frobenius-conjugate,
over $\mathcal S$, to an invertible fixed matrix $F_2$ with
\[
                              F_2=j(A)+pR_0
                             \quad(A\in M_r(E)).           \tag{FET.15}
\]
In fact $A$ is invertible over $E$.

\begin{proof}
Use FET16 and put $y=-(F_1-j(A))G_1$.
This is fixed and lies in $\mathcal N_{2k+1}$.
FET15 for the operator $Z\mapsto j(A)\phi(Z)G_1$ gives a
fixed $z\in\mathcal N_{2k+1}$ with
\[
 z=-(F_1-j(A))G_1+j(A)\phi(z)G_1+pR_1.
\]
Multiplying on the right by $F_1$ yields the exact identity
\[
 (1+z)F_1=j(A)\phi(1+z)+pR_1F_1.                           \tag{FET.16}
\]
The matrix $1+z$ is invertible. Indeed the geometric series
in $-z$, interpreted using FET14, gives $W$ with
$(1+z)W=1+pR$. The right side is invertible by the $p$-adic
geometric series in $-pR$. Hence $1+z$ has a right inverse,
and a square matrix over a commutative ring with a right inverse
is invertible (take determinants and the adjugate). Its inverse
$V_1$ is fixed.

Set $F_2=(1+z)F_1\phi(V_1)$. Equation (FET.16) proves
(FET.15), while
$F_1\phi(V_1)=V_1F_2$ gives the required conjugacy.
The matrix $F_2$ is invertible, as a product of invertible
matrices. Then
$j(A)F_2^{-1}=1-pR_0F_2^{-1}$ is invertible, so $j(A)$ is
invertible. WIS18 reflects units through $j$; applied to
$j(\det A)=\det j(A)$, it shows that $A$ is invertible over
the original $E$.
\end{proof}
% END FET17

% BEGIN FET18
\paragraph{FET18: the previously proved quotient theorem supplies every correction.}
\label{claim:FET18}
For $A\in\mathrm{GL}_r(E)$ and every $y\in M_r(\mathcal S)$
there are $z\in M_r(\mathcal S)$ and $d\in M_r(E)$ with
\[
                        y+j(A)\phi(z)-zj(A)=j(d).          \tag{FET.17}
\]

\begin{proof}
Apply the surjectivity in IFD24 to the additive operator
$\Phi(z)=j(A)\phi(z)j(A^{-1})$ on matrices modulo $j(M_r(E))$.
In entry coordinates this is exactly the coefficient matrix
$A\otimes(A^{-1})^{\mathsf t}$ used in IFD25, so IFD24
applies without a new hypothesis. Choose $z$ such that
$\Phi(z)-z$ represents the negative of the class of
$yj(A^{-1})$. Then
$yj(A^{-1})+\Phi(z)-z=j(d_0)$ for some $d_0$ over $E$.
Multiply on the right by $j(A)$ and put $d=d_0A$.
This proves (FET.17) exactly; an error modulo $p$ would already
suffice for the next step.
\end{proof}
% END FET18

% BEGIN FET19
\paragraph{FET19: the integral nonlinear change of basis converges.}
\label{claim:FET19}
For the fixed matrix $F_2=j(A)+pR_0$ in FET17 there are
$V\in\mathrm{GL}_r(\mathcal W)$ with entries in $\mathcal S$
and $A_\infty\in M_r(E)$ such that
\[
 V\equiv1\pmod {p\mathcal W},\qquad
                   F_2\phi(V)=Vj(A_\infty).               \tag{FET.18}
\]

\begin{proof}
Start with $V_0=1$, $D_0=A$. Inductively maintain
$V_n\equiv1\pmod p$, $D_n\equiv A\pmod p$, fixed entries
of $V_n$, and
\[
 F_2\phi(V_n)-V_nj(D_n)=p^{n+1}Y_n.
\]
The left side is fixed. Saturation in FET14, applied repeatedly,
puts every entry of $Y_n$ in $\mathcal S$. Use FET18 to choose
$Z_n,d_n$ with
$Y_n+j(A)\phi(Z_n)-Z_nj(A)=j(d_n)$.
Set
\[
 V_{n+1}=V_n+p^{n+1}Z_n,\qquad
 D_{n+1}=D_n+p^{n+1}d_n.
\]
Expanding the new defect and dividing its expression by
$p^{n+1}$ gives
\[
 Y_n+F_2\phi(Z_n)-V_nj(d_n)-Z_nj(D_n)
                                  -p^{n+1}Z_nj(d_n).
\]
Modulo $p$ this is zero by the choice of $Z_n,d_n$ and the
three maintained congruences. Thus the defect is divisible by
$p^{n+2}$, and all the induction assertions continue.

The $p$-weighted sums in FET14 give limits $V,A_\infty$ in
$M_r(\mathcal S),M_r(E)$ with the stated congruences at every
precision. All ring maps preserve powers of $p$; the finite
matrix operations do as well. Comparing the limiting defect
with the $n$-th defect therefore puts it in every
$p^nM_r(\mathcal W)$. Separation gives zero.
Finally $V=1+pZ$ is invertible by its integral geometric
series. This proves (FET.18) with the fixed coefficients and
without a continuity assertion beyond these explicit congruences.
\end{proof}
% END FET19

% BEGIN FET20
\paragraph{FET20: full original Frobenius-matrix decompletion.}
\label{claim:FET20}
On every original good layer, every
$F_0\in\mathrm{GL}_r(\mathcal W)$ whose entries satisfy
\texttt{Fix\allowbreak H} admits
\[
 V\in\mathrm{GL}_r(\mathcal W),\quad A'\in\mathrm{GL}_r(E),
 \qquad V_{ij}\in\mathcal S,\qquad
                          F_0\phi(V)=Vj(A').              \tag{FET.19}
\]
This proves the full original
\texttt{etale\allowbreak Phi\allowbreak Matrix\_\allowbreak decompletion\_\allowbreak to\_\allowbreak wach} endpoint.

\begin{proof}
FET16 conjugates $F_0$ to $F_1$, FET17 conjugates $F_1$ to
$F_2$, and FET19 conjugates $F_2$ to $j(A')$.
If $F\phi(V)=VF'$ and $F'\phi(W)=WF''$, then
$F\phi(VW)=VWF''$. Multiply the three changes of basis in
that order. Each is invertible and fixed, so the resulting $V$
has those properties. Also
$j(A')=V^{-1}F_0\phi(V)$ is invertible. Unit reflection WIS18
again gives $A'\in\mathrm{GL}_r(E)$.
All approximation hypotheses were discharged in FET14 from
WIS/IFD on the actual good layer. No new hypothesis on the
invariant matrix remains. The empty matrix case has $V=A'=1$
in the size-zero matrix ring and satisfies the same conclusion.
\end{proof}
% END FET20

% BEGIN FET21
\paragraph{FET21: full original etale basis on the original tensor lattice.}
\label{claim:FET21}
For every original lattice $T$ on a good layer there is a basis
$e$ of $M_T$ fixed by $H$ whose Frobenius matrix has every
entry in $j(E)$. Thus the original predicate \texttt{Et\allowbreak Basis I T e}
holds at its full scope.

\begin{proof}
Take the fixed basis of FET13. Write
$\Phi(e_j)=\sum_i F_{ij}e_i$, where
$\Phi=\phi\otimes\mathrm{id}_T$.
The Frobenius images span $M_T$: they include
$\Phi(1\otimes t)=1\otimes t$ for every $t\in T$.
Thus the square matrix $F$ has a right inverse and is invertible.
Frobenius commutes with $H$, as one checks on a pure tensor
and on Witt coordinates. The vectors $\Phi(e_j)$ are fixed,
so uniqueness of coordinates in the fixed basis gives
$F_{ij}\in\mathcal S$.

Apply FET20 and define $e'_j=\sum_i V_{ij}e_i$.
The $e'_j$ are a basis and remain fixed. Direct expansion gives
\[
 \Phi(e')=eF\phi(V)=eVj(A')=e'j(A').
\]
Thus the coordinates belong to the image of the original
$I.\mathrm{ioE}$. The matrix orientation is specified in this
calculation; no inverse transpose or changed scalar embedding
has been substituted. This proves every field of \texttt{Et\allowbreak Basis}.
\end{proof}
% END FET21

\subsection{One etale construction, including every lattice morphism}

Choose the basis in FET21 for every lattice $T$, once and for all.
For $r_T=\operatorname{rank}_O T$ define
\[
 D_T=E^{r_T},\qquad
 j_T(x)=\sum_i j(x_i)e_{T,i}\in M_T,
 \qquad \Phi(e_T)=e_Tj(A_T),\quad A_T\in\mathrm{GL}_{r_T}(E).
                                                               \tag{FET.20}
\]
Choice of these bases is permitted by the original existence
statement. The maps below are forced by the embeddings $j_T$, so
they remain coherent despite this choice.

% BEGIN FET22
\paragraph{FET22: the coordinate image and its descent test.}
\label{claim:FET22}
The map $j_T$ is injective, is $E$-semilinear through $j$, and
has $\mathcal W$-spanning image. If a finite list
$x_1,\ldots,x_q\in M_T$ is fixed by $H$ and satisfies
\[
       \Phi(x_j)=\sum_k j(B_{kj})x_k,\qquad B\in\mathrm{GL}_q(E),
                                                               \tag{FET.21}
\]
then every $x_j$ lies in the image of $j_T$.

\begin{proof}
The coordinates of $j_T(x)$ are $j(x_i)$, so injectivity follows
from WIS18. Scalar compatibility is the same coordinate identity.
The image contains all basis vectors $e_{T,i}=j_T(\mathbf e_i)$,
so it spans after extending scalars.
For (FET.21), let $U$ be the matrix of the $x_j$ in the fixed
basis. Its entries are fixed. Taking coordinates of (FET.21)
gives
\[
                       j(A_T)\phi(U)=Uj(B).
\]
IFD25, with zero inhomogeneous term, puts every entry of $U$
in $j(E)$. The corresponding coefficient columns give preimages
under $j_T$. This is descent of an already specified matrix,
which is the distinct role of IFD25 promised at the start.
\end{proof}
% END FET22

% BEGIN FET23
\paragraph{FET23: Frobenius and the quotient Galois action descend.}
\label{claim:FET23}
The image $j_T(D_T)$ is preserved by $\Phi$ and by every
$g\in G_K$. These actions descend to $D_T$ and the latter
action factors through the prescribed group $\Gamma$.
They are semilinear for the prescribed $\phi_E$ and
$\gamma_E$, satisfy the group laws, and commute.

\begin{proof}
The formula for Frobenius in (FET.20), together with
$\phi j=j\phi_E$, gives its descended map directly.
For $g\in G_K$, the list $x_j=g(e_{T,j})$ is fixed by $H$:
the group $H$ is normal and
$h g(e_{T,j})=g(g^{-1}hg)(e_{T,j})=g(e_{T,j})$.
Commutation with Frobenius and the original identity
$g(j(a))=j(\gamma_E(\operatorname{res}_\Gamma g)a)$ give
(FET.21) with
$B=\gamma_E(\operatorname{res}_\Gamma g)(A_T)$, which is
invertible. FET22 descends each basis image, and semilinearity
then descends every vector. Applying the same argument to
$g^{-1}$ gives an inverse.

The original injective cyclotomic character on $\Gamma$ and
its compatibility with $P.\chi_Q$ give
\[
             \operatorname{res}_\Gamma(g)=1
                        \quad\Longleftrightarrow\quad\chi(g)=1.
\]
For the reverse implication use injectivity of that character;
for the forward one use the injective coefficient-field map.
Thus its kernel is precisely $H$, which fixes $j_T(D_T)$
pointwise. Surjectivity of the prescribed
$\operatorname{res}_\Gamma$ gives an action of $\Gamma$:
choose any lift, with independence just proved. Applying $j_T$
reduces the identity, product and commutation laws to those on
$M_T$. Its injectivity returns those laws on $D_T$.
\end{proof}
% END FET23

% BEGIN FET24
\paragraph{FET24: every original lattice morphism descends coherently.}
\label{claim:FET24}
For every original $O$-linear equivariant lattice morphism
$f:T\to T'$, the tensor map $1\otimes f:M_T\to M_{T'}$
induces a unique $E$-linear map $D(f):D_T\to D_{T'}$ with
\[
                       j_{T'}D(f)=(1\otimes f)j_T.        \tag{FET.22}
\]
These maps commute with Frobenius and $\Gamma$, preserve identity,
composition and addition, and send scalar lattice maps to their
prescribed scalar maps.

\begin{proof}
Apply FET22 in $M_{T'}$ to the list
$x_i=(1\otimes f)e_{T,i}$. It is fixed by $H$ because $f$
is equivariant, and satisfies (FET.21) with $B=A_T$ because
the tensor map commutes with Frobenius. Hence it descends.
The prescribed coefficient identity makes the resulting map
$E$-linear, and injectivity of $j_{T'}$ gives uniqueness.

Apply $j_{T'}$ or $j_{T''}$ to each claimed compatibility.
For Frobenius and $\Gamma$ the two compositions become the
commuting actions on the tensor modules. For composition they
both become $(1\otimes(g\circ f))j_T$; for identity, addition
and scalar maps they become the respective elementary tensor
identities. Injectivity proves all equalities. Thus arbitrary
choices of bases have not caused incompatible choices of maps.
\end{proof}
% END FET24

% BEGIN FET25
\paragraph{FET25: every field of the original etale functor.}
\label{claim:FET25}
For every prime $p$, original $P,AI$, and every original good
layer $K,s,LL,\Gamma,R,I,\mathrm{IoW}$, the structure
$\mathrm{WEt}(I)$ is nonempty.

\begin{proof}
Set its module at $T$ to $D_T=E^{r_T}$ from (FET.20).
It is free of the original rank. Its additive group and its
$E$-module structure are the coordinatewise ones; restriction
along the original maps $O\to\mathrm{LetA}\to E$ gives the
two other module structures and both required scalar towers.
FET23 supplies its Frobenius,
quotient Galois action, their semilinearity and all group and
commutation laws. FET24 supplies its map on every original
lattice morphism, its $E$-linearity, the Frobenius and Galois
compatibilities, identity, composition, addition and scalar laws.
The maps $j_T$ in FET22 supply the original injections, scalar
identities and spanning property. Their Frobenius and Galois
intertwining follow from FET23, and their naturality from
(FET.22). These are exactly the fields of the original
\texttt{WEt} record and the full conclusion of
\texttt{Et\allowbreak Mod.\allowbreak nonempty\_\allowbreak w\allowbreak Et}.

The only inputs at this endpoint are the original good-layer
records. Faithful invariants, the fixed basis, nonlinear
decompletion and morphism descent have all been proved above or
by the named earlier human proofs. No crystalline representation
is needed to define this etale functor on general lattices.
\end{proof}
% END FET25

% BEGIN FET26
\paragraph{FET26: the actual etale input is now available for the Wach return.}
\label{claim:FET26}
Retain the original crystalline-realization and local-tied-data
premises of WPR30. On every supplied original good layer there
now exist the full original etale functor and its associated
original Wach structure. This does not yet assert the later
Herr or quantitative comparison conclusions.

\begin{proof}
FET25 constructs a $\mathrm{WEt}(I)$ on that layer. WPR30,
which proves the original implication from
\texttt{Bhatt\allowbreak Scholze\allowbreak Crystalline\allowbreak Lattices} and \texttt{Loc\allowbreak Tied}
to \texttt{Wach\allowbreak Ex}, applies to this very functor. It gives
the original $\mathrm{NS}$, including its lattice, maps and
comparison fields at the scope already proved in WLD26/WPR30.
The coefficient inclusion, Frobenius, Galois action and lattice
morphisms used by both constructions are the ones just fixed.
The original layerwise quantifier requires this construction for
each allowed layer; it adds no transition-system hypothesis
between distinct outer indices.

The next consumer still invokes the original Herr class maps
and the quantitative comparison. Their existence and bounds do
not follow just from an etale functor or a Wach lattice. Likewise
the common global \texttt{Concl\allowbreak F} family and all of its local
row estimates remain to be supplied before
\texttt{GLOBAL\_\allowbreak DENOMINATOR} can be closed.
\end{proof}
% END FET26

\paragraph{Source dictionary and the next proof.}
The original targets are \texttt{Conv\allowbreak Zero.\allowbreak Faith\allowbreak H},
\texttt{Fixed\allowbreak Basis.\allowbreak main}, \texttt{Mat\allowbreak Decompl.\allowbreak decompletion},
\texttt{Etale\allowbreak Basis.\allowbreak exists\_\allowbreak et\allowbreak Basis}, and
\texttt{Et\allowbreak Mod.\allowbreak nonempty\_\allowbreak w\allowbreak Et}, all on the accepted source pin.
FET14--FET20 prove construction of the nonlinear conjugacy,
whereas IFD25 remains the earlier linear descent theorem used
in FET22. This distinction must remain visible in the roadmap.

For the next step read the actual inputs of
\texttt{Exists\allowbreak Local\allowbreak Wach\allowbreak Functor.\allowbreak lean}, then the Herr class maps
and the quantitative local comparison bodies that it calls.
The signed crystalline-lifting estimates must preserve the
full inverse-image finite local condition, including torsion,
the coefficient-normalized length, both cyclic orders of the
signed clusters, and the same original maps. The present result
does not provide those estimates by changing their hypotheses.

\paragraph{Classical foundation references.}
Peter Scholze, \emph{Perfectoid spaces}, Theorem~3.7(i),
\url{https://www.math.uni-bonn.de/people/scholze/PerfectoidSpaces.pdf}.
Bhargav Bhatt, Matthew Morrow and Peter Scholze,
\emph{Integral $p$-adic Hodge theory}, Definition~3.5 and
Example~3.6,
\url{https://www.math.uni-bonn.de/people/scholze/integralpadicHodge.pdf}.
FET3 supplies the elementary sharp-coordinate argument needed
here, so neither a full tilting equivalence nor an unprinted
fixed-field identification is being used as an additional lemma.

\endgroup
\end{document}
